\documentclass[11pt,reqno]{amsart}
\usepackage[T1]{fontenc}
\usepackage{lmodern,amssymb,mathtools,microtype,booktabs}
\usepackage[a4paper,textwidth=158mm,textheight=246mm,centering,headsep=13pt]{geometry}
\usepackage{xcolor}
\definecolor{linkcolor}{RGB}{0,104,107}
\usepackage[colorlinks=true,allcolors=linkcolor,bookmarksnumbered=true]{hyperref}
\hypersetup{pdftitle={Invariant tensor dimensions as Hamburger moments},pdfauthor={T. Banica and S. Palcoux},pdfsubject={Version 29}}
\newtheorem{theorem}{Theorem}[section]
\newtheorem{proposition}[theorem]{Proposition}
\newtheorem{corollary}[theorem]{Corollary}
\newtheorem{lemma}[theorem]{Lemma}
\newtheorem*{theoremA}{Theorem A}
\newtheorem*{theoremB}{Theorem B}
\newtheorem*{theoremC}{Theorem C}
\newtheorem*{theoremD}{Theorem D}
\newtheorem*{theoremE}{Theorem E}
\theoremstyle{definition}
\newtheorem{definition}[theorem]{Definition}
\newtheorem{example}[theorem]{Example}
\newtheorem{remark}[theorem]{Remark}
\newtheorem{question}[theorem]{Question}
\newcommand{\Hom}{\operatorname{Hom}}
\newcommand{\End}{\operatorname{End}}
\newcommand{\Ext}{\operatorname{Ext}}
\newcommand{\FPdim}{\operatorname{FPdim}}
\newcommand{\one}{\mathbf 1}
\newcommand{\C}{\mathcal C}
\newcommand{\PX}{\mathcal P_X}
\newcommand{\RR}{\mathbb R}
\newcommand{\CC}{\mathbb C}
\newcommand{\ZZ}{\mathbb Z}
\newcommand{\NN}{\mathbb N}
\newcommand{\Ga}{\mathbb G_a}
\newcommand{\Gm}{\mathbb G_m}
\newcommand{\Rep}{\operatorname{Rep}}
\newcommand{\Vect}{\operatorname{Vec}}
\DeclareMathOperator{\supp}{supp}
\DeclareMathOperator{\Gr}{Gr}
\DeclareMathOperator{\rad}{rad}
\DeclareMathOperator{\Lie}{Lie}
\setcounter{tocdepth}{1}
\setlength{\emergencystretch}{1.3em}
\title[Invariant tensor moments]{Invariant tensor dimensions as Hamburger moments}
\author{T. Banica and S. Palcoux}
\keywords{Invariant tensors, Hamburger moment problem, tensor categories, orthogonal polynomials, Jordan blocks, projective covers, fusion rings}
\date{}
\begin{document}


\begin{abstract}
For a self-dual object in a complex tensor category, we characterize the Hamburger moment property of its invariant tensor dimensions by a Cauchy--Schwarz inequality on sums of its tensor powers. Imposing the inequality on all pairs of objects determines the indecomposables in each block and the tensor rules of the unit block. For the additive group, the even invariant dimensions recover the Jordan type. In a finite tensor category of fusion type, Hamburger positivity for the projective cover of the unit forces unimodularity. If the identity and double dual are naturally isomorphic as additive functors, this positivity is equivalent to semisimplicity. A zero Cartan eigenvalue yields an explicit negative Hankel square in the unimodular case. We also prove semisimplicity of every finite categorification of the even $\mathfrak{sl}_2$ Verlinde ring at every odd level. Finally, we give a Stieltjes criterion for $X\otimes X^*$ and a six-dimensional representation whose doubled sequence has a 42-atom representing measure through degree 83 and fails at degree 84.

\end{abstract}
\maketitle
\tableofcontents
\section*{Introduction}

For a finite-dimensional representation $V$ of a compact group $G$, averaging the action on $V^{\otimes n}$ gives
\[
       \dim(V^{\otimes n})^G=\int_G\chi_V(g)^n\,dg .
\]
If $V$ is self-dual, its character is real and the distribution of $\chi_V$ represents every invariant dimension. This connects the dimension data studied by Larsen and Pink \cite{LP} to the classical Hamburger moment problem. For compact quantum groups the fundamental characters of $O_N^+$ for $N\ge2$ and $S_N^+$ for $N\ge4$ have the semicircle and rate-one free Poisson distributions, respectively \cite{Banica,BCS}.

In a complex tensor category $\C$ put
\[
 c_n(X)=\dim\Hom(\one,X^{\otimes n}),\qquad X^{\otimes0}=\one,\qquad
 h(A,B)=\dim\Hom(A,B).
\]
When are all the $c_n(X)$ moments of a positive measure on the real line? What does its existence imply about $\C$? Forey, Fres\'an and Kowalski \cite[Propositions~2.1, 2.2 and~2.5]{FFK} obtain spectral measures from global inequalities for additive invariants of symmetric tensor categories. Rigidity gives a local criterion for a single object in an arbitrary tensor category.

\begin{theoremA}
Let $X\simeq X^*$ be an object of a complex tensor category. Its numbers $c_n(X)$ are the moments of a positive Borel measure on $\RR$ if and only if
\[
             h(A,B)^2\le h(A,A)h(B,B)
\]
for all finite direct sums $A,B$ of tensor powers of $X$, including the zeroth power.
\end{theoremA}

Indeed $c_{i+j}(X)=h(X^{\otimes i},X^{\otimes j})$. The positive and negative coefficient parts of a polynomial with integral coefficients give two actual objects, as in \cite[proof of Proposition~2.1]{FFK}. A negative square for the moment functional thereby gives a pair violating the inequality. This equivalence only needs a complex-linear additive monoidal category with bilinear tensor product, $\End(\one)=\CC$, an object $X$ isomorphic to its left dual, and finite $c_n(X)$ (Theorem~\ref{thm:main}). For the Taft algebra of dimension nine, a self-dual indecomposable and the direct sum of its composition factors have the same ordinary Grothendieck class. The indecomposable has a Hankel determinant equal to $-1$. The direct sum has the character law of a cyclic group (Theorem~\ref{thm:taft-failure}).

The inequality in Theorem A concerns one self-dual $X$. Requiring it for every self-dual $X$ gives \emph{tensor-power Cauchy--Schwarz}. Requiring it for every ordered pair of objects $A,B$ gives \emph{global Cauchy--Schwarz}. The Sweedler category satisfies the former and fails the latter (Proposition~\ref{prop:sweedler}). Theorem~\ref{thm:ext} characterizes the global inequality by vanishing of $\Ext^1(S,T)$ between distinct simples and the bound $\dim\Ext^1(S,S)\le1$. In a nonsemisimple tensor category satisfying it, each simple $S$ has a unique indecomposable $J_{S,r}$ of every length $r$ in its block. Writing $J_r=J_{\one,r}$, we obtain
\[
 J_r\otimes S\simeq J_{S,r}\simeq S\otimes J_r,\qquad
 J_r\otimes J_s\simeq\bigoplus_{j=0}^{\min(r,s)-1}J_{r+s-1-2j}
 \qquad\text{(Theorem~\ref{thm:universal-jordan}).}
\]
Thus the unit block has the split Grothendieck ring and Hom dimensions of $\Rep(\Ga)$. For $\Ga$ the measure is the law of the lifted $SL_2$-character on the torus under $(1+\cos\theta)\,d\theta/(2\pi)$.

\begin{theoremB}
Let $G$ be a connected complex affine algebraic group, and take finite-dimensional rational representations.
\begin{enumerate}
\item For $G=\Ga$, the even invariant dimensions $c_{2n}(X)$ determine the Jordan blocks of $X$.
\item For solvable $G$, every self-dual module has a Hamburger law precisely when $G$ is a torus or the product of a torus with $\Ga$.
\item Each group $\CC^N\rtimes SL_N$ or $\CC^N\rtimes GL_N$ for $N\ge2$, and each complex Euclidean group $E_m=\CC^m\rtimes SO_m$ for $m\ge2$, has a self-dual module without a Hamburger law.
\item For the double cover $\widetilde E_3=\CC^3\rtimes SL_2\subset Sp_4$, every module $\bigoplus_i Y_i|_{\widetilde E_3}\otimes W_i$, with $Y_i\in\Rep(Sp_4)$ and $W_i\in\Rep(SL_2)$ inflated to $\widetilde E_3$, has a law given by a character distribution on $SU(2)\times SU(2)$.
\end{enumerate}
\end{theoremB}

Corollary~\ref{cor:finite-reconstruction} gives a finite reconstruction of the Jordan type in (1). The global inequality is classified for all complex affine algebraic groups of finite type, including disconnected groups. It holds exactly for $R$ and $R\times\Ga$ with $R$ linearly reductive (Theorem~\ref{thm:global-groups}). For tensor-power Cauchy--Schwarz, (2) classifies the connected solvable case. The connected nonsolvable case remains open in general.

Both $E_3$ and its double cover have modules with and without laws. The $14$-dimensional sum of the affine module, its dual and the adjoint module of $E_3$ has a negative Hankel square, which persists under inflation to the cover. The module $E\oplus W$ of Theorem~\ref{thm:euclidean} has a law already on $E_3$. The spinor is genuine on the cover and has the law of the sum of two independent semicircular variables (Theorem~\ref{thm:contraction}).

For a finite tensor category let $P=P(\one)$ be the projective cover of the tensor unit. We call $\C$ \emph{of fusion type} when its ordinary Grothendieck ring is a fusion ring in the simple basis. The moment problem for this single canonical object gives a semisimplicity theorem.

\begin{theoremC}
Let $\C$ be a finite tensor category over $\CC$ of fusion type. If the invariant tensor dimensions of $P(\one)$ are Hamburger moments, then $\C$ is unimodular. If the identity and double-dual functors are naturally isomorphic as additive functors, then these dimensions are Hamburger moments if and only if $\C$ is semisimple.
\end{theoremC}

Write $P=P(\one)$. The nonzero poles and residues of its rational moment transform first force $c_1(P)>0$, so the simple socle of $P$ is the unit. In this unimodular case $P\simeq P^*$ and $c_{n+1}(P)$ is the $n$th moment of the Cartan spectral measure $\nu$ at the unit. A measure $\mu$ representing $c_n(P)$ would satisfy $t\,d\mu=d\nu$, by uniqueness on compact support. Cartan singularity in the nonsemisimple case \cite[Theorem~6.6.1]{EGNO} gives a positive atom of $\nu$ at zero and hence an explicit negative square for $P$ (Corollary~\ref{cor:cartan-failure}). For the quaternion example the first five counts $1,1,2,5,14$ agree with a positive law, but the size-four Hankel determinant is $-3$ (Example~\ref{ex:quaternion-witness}).

There is a related structural question. Which fusion rings admit a nonsemisimple finite categorification? Etingof's answer to the question of the second author \cite{MOEtingof} gives early obstructions. Call a fusion ring \emph{fusion-simple} if it has no proper nontrivial based fusion subring.

\begin{theoremD}
Let $\C$ be a finite tensor category over $\CC$ of fusion type, with Grothendieck ring $R$.
\begin{enumerate}
\item If the identity and double-dual functors are isomorphic as additive functors, and $[P(\one)]$ is not a zero divisor in $R$, for instance if $R\otimes\mathbb Q$ is a field, then $\C$ is semisimple.
\item If $\C$ is braided, then $\C$ is semisimple whenever its M\"uger center is Tannakian, in particular when $\C$ is nondegenerate. For braided $\C$ the same conclusion holds if $R\neq\ZZ[C_2]$ and either $R$ has rank at most three or $R$ is fusion-simple.
\item If $R$ is the even part $\widetilde K_l$ of the Verlinde ring of $\mathfrak{sl}_2$ at level $2l+1$, where $l\ge1$, then $\C$ is semisimple.
\end{enumerate}
\end{theoremD}

Part (2) gives a necessary condition on the ring of a nonsemisimple braided category. Some nonunit simple has integral Frobenius--Perron dimension and a bipartite fusion graph on the full simple basis (Theorem~\ref{thm:bipartite}). This forces semisimplicity for any braided finite realization of 211 of the 303 commutative multiplicity-free tables through rank nine. For (3), the double dual acts on first extensions through a twist of order dividing both $2$ and $2l+3$. The resulting separation of the two signs in the sphericalization proves semisimplicity at every odd level. It also implies that every nonzero object of the center of the corresponding even Temperley--Lieb fusion category has infinite-dimensional Nichols algebra (Corollary~\ref{cor:verlinde-nichols}).

There is a way to recognize the hypothesis of (3) from invariant counts. If a self-dual $X$ has semisimple powers through degree $l+1$ and its counts through degree $2l+2$ agree with those of the marked generator of $\widetilde K_l$, the summands of its powers generate a fusion subcategory with that marked ring. If the ambient category is finite and every simple occurs as a composition factor of a power of $X$, it follows that the ambient category is semisimple (Theorem~\ref{thm:small-recognition}).

Finally, passing from $X$ to $X\otimes X^*$ gives a Stieltjes moment problem.
\begin{theoremE}
Let $X^{**}\simeq X$ in a complex tensor category, and put $Z=X\otimes X^*$. The numbers $c_k(Z)$ are moments of a positive measure on $[0,\infty)$ precisely when the inequality of Theorem A holds both on finite sums of the $Z^{\otimes i}$ and on finite sums of the $Z^{\otimes i}\otimes X$, in each case with $i\ge0$. If $X\simeq X^*$ has a law $\mu$, its pushforward by $t\mapsto t^2$ represents $c_k(Z)$. A doubled law can also exist for an $X$ without a law.
\end{theoremE}

For $X=\Lambda^2(\CC^4)$ restricted to $\CC^3\rtimes SL_3$, write $d_k=c_{2k}(X)$. A probability measure on exactly 42 points in $(0,36)$ represents $d_0,\ldots,d_{83}$, whereas appending $d_{84}$ makes a Hamburger law impossible (Theorem~\ref{thm:exterior-doubled}). Proposition~\ref{prop:long-truncations} gives arbitrarily long positive Stieltjes truncations without a full law, with the object and category varying. For normal objects, Section~\ref{sec:complex} gives a local criterion for complex mixed moments under a radial Carleman hypothesis. It also compares two invariants on the multiplicity-free fusion-ring census. The multiset of spectral laws of the normal basic elements, with one marker for each nonnormal element, has collisions at rank nine. The multiset of tuples of unit coefficients of $b^2,b^3,b^4,b^5$, over all basic elements $b$, distinguishes all 353 tables through rank nine.

Throughout, scalar products are linear in the first variable.

\emph{Acknowledgements.} We thank Claude and GPT for their substantial involvement in this project. An AI-assisted review of an earlier draft suggested the unimodularity argument in Proposition~\ref{prop:fusion-type}(5), the use of the two signs in Theorem~\ref{thm:verlinde}, and the bipartite graph obstruction in Theorem~\ref{thm:bipartite}.

\section{Moments and Hankel matrices}\label{sec:moments}
A positive Borel measure $\mu$ on $\RR$ represents $(c_n)_{n\ge0}$ if every absolute moment is finite and $c_n=\int t^n\,d\mu(t)$. Here $c_0=1$, so $\mu$ is a probability measure. Write $L(t^n)=c_n$ and extend $L$ linearly to $\RR[t]$. For $p=\sum_{i=0}^d a_it^i$,
\[
 L(p^2)=\sum_{i,j=0}^d a_i a_j c_{i+j}=a^{\mathsf T}H_da,
 \qquad H_d=(c_{i+j})_{0\le i,j\le d}.
\]
Hamburger's theorem says that a representing measure exists exactly when every $H_d$ is positive semidefinite. Stieltjes' theorem gives a measure on $[0,\infty)$ exactly when the shifted matrices $(c_{i+j+1})_{0\le i,j\le d}$ are also positive semidefinite \cite[Theorems~3.8 and~3.12]{Sch}.

\begin{example}\label{ex:classical}
A point mass $\delta_a$ gives $c_n=a^n$ and $H_d$ of rank one. The semicircle law $\sqrt{4-t^2}\,dt/(2\pi)$ and the arcsine law $dt/(\pi\sqrt{4-t^2})$ on $[-2,2]$ have vanishing odd moments. Their even moments are, respectively, $\binom{2k}{k}/(k+1)$ and $\binom{2k}{k}$. The arcsine law is the distribution of $2\cos\theta$ for uniform $\theta$, a useful comparison for Section~\ref{sec:groups}. In contrast, the first 4 numbers $(1,0,0,1)$ admit no continuation to a Hamburger sequence. Zero second moment forces a point mass at zero. The same failure is visible in the principal minor $c_2c_4-c_3^2=-1$, whatever the proposed value of $c_4$.
\end{example}

\paragraph{Support and uniqueness.}\label{rem:support}
If a representing measure exists, a nonzero polynomial in the kernel of $H_d$ vanishes on its support. Conversely, for a measure supported on exactly $k$ points the Vandermonde matrix gives $\operatorname{rank}H_d=\min(k,d+1)$ and the first $k$ moments determine the masses on these points \cite[Proposition~3.9]{Sch}. Thus some $H_d$ is singular exactly when the support is finite. More generally, $R=\lim_{n\to\infty}c_{2n}^{1/(2n)}$ is the essential supremum of $|t|$. If $R<\infty$, polynomial approximation on $[-R,R]$ gives uniqueness. Carleman's condition $\sum_{n\ge1}c_{2n}^{-1/(2n)}=\infty$ also suffices for uniqueness \cite[Theorem~4.3(i)]{Sch}. A zero even moment of positive degree forces $\mu=\delta_0$.

\section{Tensor categories, from scratch}\label{sec:tensor}
The model for a complex-linear abelian category is the category of finite-dimensional modules over a complex algebra, with module maps as morphisms. Tensor products and duals add precisely the structure needed to turn invariant dimensions into entries of a Hankel matrix.

\subsection{Linear and abelian categories}
We work over $\CC$ with essentially small additive categories. Hom spaces are complex vector spaces and composition is bilinear. In an abelian category one has kernels, cokernels and short exact sequences, as for modules. A nonzero object $S$ is \emph{simple} if its only subobjects are $0$ and $S$. A complex-linear abelian category is \emph{locally finite} if its Hom spaces are finite-dimensional and every object has a finite filtration with simple quotients. Jordan--H\"older makes their multiplicities $[A:S]$ well defined. Krull--Schmidt also gives a unique decomposition of $A$ into indecomposable direct summands, up to permutation and isomorphism \cite[\S1.5]{EGNO}.

An object is \emph{semisimple} if it is a direct sum of simples, and the category is semisimple if every object is. Schur's lemma gives $\dim\Hom(S,T)=\delta_{S,T}$ for simple objects, where the delta refers to isomorphism classes. In particular, for semisimple $A,B$,
\begin{equation}\label{eq:mult-form}
 h(A,B)=\sum_S[A:S][B:S].
\end{equation}
For general $A$, $h(S,A)$ counts copies of $S$ in the \emph{socle} (the sum of the simple subobjects). A nonsplit length-two extension already separates this count from the composition multiplicity. An object $P$ is \emph{projective} if every epimorphism onto $P$ splits, equivalently if $\Hom(P,-)$ preserves epimorphisms.

\subsection{Tensor products, unit and duals}
A \emph{tensor category} is a locally finite complex-linear abelian category equipped with a bilinear monoidal product $\otimes$, a unit $\one$ with $\End(\one)=\CC$, and left and right duals. We suppress the associativity and unit constraints. Tensoring is exact in either variable and $\one$ is simple \cite[Definition~4.1.1, Proposition~4.2.1 and Theorem~4.3.8(i)]{EGNO}. A tensor category is \emph{finite} if its underlying abelian category is equivalent to the category of finite-dimensional modules over a finite-dimensional algebra. It then has finitely many simple isomorphism classes and enough projectives. A finite semisimple tensor category is a \emph{fusion category}.

For a Hopf algebra $H$ with bijective antipode $S$, its coproduct gives the action on $V\otimes W$, its counit $\varepsilon$ gives the action on $\one=\CC$, and $(af)(v)=f(S(a)v)$ defines the action on the left dual $V^*$. A map $\one\to V$ is determined by an invariant vector, that is, by $v\in V$ such that $av=\varepsilon(a)v$ for all $a\in H$. For a group algebra the tensor action is diagonal and these are the fixed vectors. For an enveloping algebra the Lie algebra acts on tensor products by the Leibniz rule and an invariant vector is killed by every Lie algebra element.

We use $V^*$ for a left dual and ${}^*V$ for a right dual. Evaluation $\mathrm{ev}_V:V^*\otimes V\to\one$ and coevaluation $\mathrm{coev}_V:\one\to V\otimes V^*$ satisfy the triangle identities. In vector spaces these are $f\otimes v\mapsto f(v)$ and $1\mapsto\sum_i e_i\otimes e^i$ for dual bases. Since $(V\otimes W)^*\simeq W^*\otimes V^*$, every power of a self-dual object is self-dual.

\begin{example}\label{ex:tensorexamples}
(1) For a finite group $G$, averaging gives Maschke's theorem and $c_n(V)=|G|^{-1}\sum_{g\in G}\chi_V(g)^n$.
(2) In $\Vect_\Gamma$, the finite-dimensional vector spaces graded by a group $\Gamma$, each simple $\delta_g$ is a line in degree $g$. Hence $\delta_g\otimes\delta_h=\delta_{gh}$ and $\delta_g^*=\delta_{g^{-1}}$.
(3) The compact quantum groups $O_N^+$ ($N\ge2$) and $S_N^+$ ($N\ge4$) have semisimple representation categories with infinitely many simples. The Haar state gives the averaging operation.
(4) Both $\Rep(\Ga)$ and the representation category of a Taft algebra $T_m$, $m\ge2$, are nonsemisimple. The Taft category is finite. In $\Rep(\Ga)$ there is only one simple, the unit, which has no projective cover (Example~\ref{thm:min} and Theorem~\ref{thm:unbounded}).
\end{example}

\begin{theorem}[Frobenius reciprocity]\label{thm:frobenius}
For objects $U,V,W$ of a tensor category there are natural isomorphisms
\begin{equation}\label{eq:rigidity}
 \Hom(U\otimes V,W)\simeq\Hom(U,W\otimes V^*),\qquad
 \Hom(V^*\otimes U,W)\simeq\Hom(U,V\otimes W).
\end{equation}
In particular, for $i,j\ge0$, $\Hom(\one,X^{\otimes(i+j)})\simeq\Hom((X^{\otimes i})^*,X^{\otimes j})$. If $X\simeq X^*$, then $h(X^{\otimes i},X^{\otimes j})=c_{i+j}(X)$.
\end{theorem}
\begin{proof}
The adjunctions insert a coevaluation in one direction and contract an evaluation in the other. The triangle identities make the maps inverse \cite[Proposition~2.10.8(a)]{EGNO}. For example, the second map sends $f:V^*\otimes U\to W$ to $(\mathrm{id}_V\otimes f)(\mathrm{coev}_V\otimes\mathrm{id}_U)$. For $H$-modules and $U=\one$, this restricts the tensor--Hom identification $\Hom_{\CC}(V^*,W)\simeq V\otimes W$ to intertwiners on the left and invariant tensors on the right. Take $V=X^{\otimes i}$, $W=X^{\otimes j}$ to get the moment identity.
\end{proof}

\section{The Cauchy--Schwarz moment criterion}\label{sec:criterion}
\subsection{The objectwise characterization}
By additivity, $h(A,B)=\dim\Hom(A,B)$ extends to a bilinear form on the split Grothendieck ring $K^{\mathrm{sp}}_0(\C)$, whose additive basis consists of the indecomposable classes. In finite tensor categories Wang--Li--Zhang study the related Green-ring pairing $([A],[B])=h(A,B^*)$ \cite[\S3.1]{WLZ}. For self-dual $X$ the two forms agree on its powers. There $h(X^{\otimes i},X^{\otimes j})=c_{i+j}$ by rigidity. On the ordinary Grothendieck ring $\Gr(\C)$ one also imposes exact-sequence relations, and even $h(\one,-)$ may depend on the extension class.

Forey--Fres\'an--Kowalski obtain real spectral measures from an additive invariant satisfying an arithmetic-mean inequality on all pairs in a symmetric category of self-dual objects \cite[Proposition~2.1]{FFK}. For invariant dimensions, rigidity lets us test necessity and sufficiency on sums of powers of a single self-dual object in any tensor category.

\begin{theorem}[Objectwise Cauchy--Schwarz characterization]\label{thm:main}
Suppose $X\simeq X^*$, put $x=[X]$, and let $\PX$ consist of the finite direct sums
of $X^{\otimes i}$, $i\ge0$. The following are equivalent.
\begin{enumerate}
\item The sequence $(c_n(X))_{n\ge0}$ has a representing measure on $\RR$.
\item For all $A,B\in\PX$,
\begin{equation}\label{eq:cs}
                    h(A,B)^2\le h(A,A)h(B,B).
\end{equation}
\item The restriction of $h$ to the polynomial subalgebra
$\RR[x]=\{p(x):p\in\RR[t]\}$
of $K^{\mathrm{sp}}_0(\C)\otimes_{\ZZ}\RR$ is positive semidefinite.
\end{enumerate}
\end{theorem}

\begin{proof}
Put $L_X(t^n)=c_n(X)$, and let
$a(X)=\bigoplus_i(X^{\otimes i})^{\oplus a_i}$ for
$a(t)=\sum_i a_it^i\in\ZZ_{\ge0}[t]$. Additivity gives
\begin{equation}\label{eq:bridge}
 h(a(X),b(X))=L_X(ab).
\end{equation}
By bilinear extension $h(p(x),q(x))=L_X(pq)$ for arbitrary real
polynomials. In particular this restriction is symmetric.
Hamburger's theorem gives (1)$\Leftrightarrow$(3), and integration
against a representing measure gives (2) by Cauchy--Schwarz in
$L^2(\mu)$.

For (2)$\Rightarrow$(1) we use the sign decomposition of
\cite[proof of Proposition~2.1]{FFK}. Write an integral polynomial
$p=a-b$ with $a,b\in\ZZ_{\ge0}[t]$. Applying \eqref{eq:cs}
to $a(X),b(X)$ in \eqref{eq:bridge} yields
\[
 L_X(p^2)\ge
 \bigl(\sqrt{h(a(X),a(X))}-\sqrt{h(b(X),b(X))}\bigr)^2\ge0.
\]
Scaling gives the same inequality for rational coefficients, and
continuity in each fixed degree gives it for real coefficients.
\end{proof}

In the present setting the arithmetic-mean inequality of \cite{FFK} is $2h(A,B)\le h(A,A)+h(B,B)$. Applying it to $A^{\oplus m},B^{\oplus n}$ and varying $m/n$ gives \eqref{eq:cs}. Conversely, a negative square for $L_X$ gives a concrete violating pair. Approximate its polynomial in the same degree by one with rational coefficients, clear denominators, and take the objects corresponding to its positive and negative coefficient parts in the proof above.

The equivalence (1)$\Leftrightarrow$(2) needs only a complex-linear additive monoidal category with bilinear tensor product, $\End(\one)=\CC$, a left dual $X^*\simeq X$ and finite $h(\one,X^{\otimes n})$ for every $n$. We call (2) for every self-dual $X$ \emph{tensor-power Cauchy--Schwarz}, and \eqref{eq:cs} for every ordered pair in $\C$ \emph{global Cauchy--Schwarz}.

If the invariant dimensions have a representing measure and $\operatorname{span}_{\RR}\{1,x,x^2,\ldots\}$ has finite dimension $d$, its Hankel matrices have rank at most $d$ by \eqref{eq:bridge}, so the measure has at most $d$ support points. A fiber functor $F$ gives $c_{2n}\le(\dim F(X))^{2n}$ by faithfulness. In a finite category, $\FPdim(X)$ is the spectral radius of multiplication by $[X]$ in the simple basis. Since $h(\one,Y)\le[Y:\one]\le\FPdim(Y)$ by additivity and $\FPdim$ is multiplicative \cite[\S4.5]{EGNO}, $c_{2n}\le\FPdim(X)^{2n}$. If a representing measure exists, either growth bound makes it compactly supported and unique by Section~\ref{sec:moments}. In Deligne's semisimple $\underline{\Rep}(GL_t)$ the rational diagram presentation gives the same matching bases at each complex $t$ transcendental over $\mathbb Q$. The tensor cube of the standard object plus its dual, $(V\oplus V^*)^{\otimes3}$, has the indeterminate moments of $G^3$ for a centered Gaussian $G$ of variance 2 \cite[Remark~3.8]{FFK}.

The sign character of $\ZZ/2$ has law $(\delta_{-1}+\delta_1)/2$. Section~\ref{sec:doubling} gives the additional Hom inequality that yields a measure on $[0,\infty)$ for the doubled object $Y\otimes Y^*$ when $Y^{**}\simeq Y$.

\subsection{The semisimple case and its spectral measure}\label{sec:fusion}
Suppose all powers of the self-dual $X$ are semisimple. Applying \eqref{eq:mult-form} and Theorem~\ref{thm:frobenius} to their multiplicities $m_i(S)=[X^{\otimes i}:S]$ gives, as in \cite[Corollary~2.6]{FFK},
\begin{equation}\label{eq:gram}
 c_{i+j}(X)=\sum_{S}m_i(S)m_j(S).
\end{equation}
Each multiplicity vector has finite support, so every Hankel matrix is a Gram matrix even when infinitely many simples occur among all powers. The same formula for arbitrary objects in a semisimple category gives global Cauchy--Schwarz.

The support can be unbounded even in a semisimple category. For the standard object of Deligne's orthogonal category at a transcendental complex parameter, a basis of invariants is indexed by pairings. Thus $c_{2j}=(2j-1)!!$ and $c_{2j+1}=0$, the moments of the standard real Gaussian \cite[Remark~3.8 and the following paragraph]{FFK}.

A finite-rank fusion ring $R$ has a distinguished basis $B\ni1$ and a
basis-preserving anti-involution $*$ such that
$\tau(bc)=\delta_{c,b^*}$, where $\tau$ takes the unit coefficient
\cite[\S3.1]{EGNO}. Its structure constants are integers $\geq0$.
Extend $*$ conjugate-linearly to $R_{\CC}$ and put
$\langle u,v\rangle=\tau(uv^*)$. The basis $B$ is orthonormal,
$\tau$ is a trace, and left multiplication $N_x$ has adjoint $N_{x^*}$.
For $x=x^*$ the scalar spectral measure $\nu_x$ of $N_x$ at $1$ satisfies
\[
        \int t^n\,d\nu_x(t)=\langle N_x^n1,1\rangle=\tau(x^n).
\]
This construction applies to noncommutative rings and needs no
categorification. If $R=\Gr(\C)$ for a fusion category and $x=[X]$,
semisimplicity gives $\tau(x^n)=c_n(X)$. The difference between these
2 quantities in a nonsemisimple category will be essential below.

\begin{proposition}\label{thm:fp}
Let $0\ne x=x^*\in\ZZ_{\ge0}B$. If $S$ is the fusion subring generated
by the basis elements occurring in $x$, then
\[
 \max\supp\nu_x=\FPdim(x),\qquad
 \nu_x(\{\FPdim(x)\})=\frac{1}{\FPdim(S)}.
\]
\end{proposition}
\begin{proof}
The union of the supports of $x^n$ ($n\ge0$) is closed under duality and under constituents of products, by positivity. It is therefore the basis of $S$. Every vertex is reachable from $1$, so $N_x$ restricted to $S$ is symmetric and irreducible. Its positive eigenvector $v=(\FPdim b)_{b\in B\cap S}$ has eigenvalue $\FPdim(x)$ \cite[Propositions~3.3.6 and~3.3.11]{EGNO}. Perron--Frobenius gives the top eigenline $\CC v$. All $N_x^n1$ lie in $S$, so the top atom has mass $v_1^2/\|v\|^2=1/\FPdim(S)$, where $\FPdim(S)=\sum_{b\in B\cap S}\FPdim(b)^2$. For a self-dual representation $X$ of a finite group $G$, the top atom has mass $1/|G/\ker X|$, since $X$ tensor-generates $\Rep(G/\ker X)$.
\end{proof}

\begin{example}\label{ex:fusion}
In the Fibonacci ring, $x^2=1+x$ and $N_x=\left(\begin{smallmatrix}0&1\\1&1\end{smallmatrix}\right)$. With $\varphi=(1+\sqrt5)/2$, the spectral law is
\[
     \nu_x=\frac{1}{\varphi+2}\delta_\varphi+
                \frac{\varphi+1}{\varphi+2}\delta_{-1/\varphi}.
\]
The 2 eigenvalues and the zero first moment determine these weights. For the complex 2-dimensional irreducible $\rho$ of $S_3$, the character takes the values $2,0,-1$ on conjugacy classes of sizes $1,3,2$. Hence $\mu_\rho=\delta_2/6+\delta_0/2+\delta_{-1}/3$. The top atoms in these examples have masses $1/(\varphi+2)$ and $1/6$ as predicted by Proposition~\ref{thm:fp}.
\end{example}

Even for an invertible simple the real moment property can fail without
self-duality. A generator of $\Vect_{\ZZ/3}$ has $c_2=0$ and $c_3=1$,
whereas a measure on $\RR$ of second moment zero is concentrated at zero.
On $\CC$ this object has the uniform law on the cube roots of unity when
one prescribes the mixed moments $\tau(x^a(x^*)^b)$.

\begin{remark}[Infinite divisibility]\label{thm:fid}
For $x=x^*\in\ZZ_{\ge0}B$ in a finite-rank fusion ring, $\nu_x$ is freely infinitely divisible precisely when $x=n\cdot1$, $n\ge0$. Indeed,
$\operatorname{Var}(\nu_x)=\tau(x^2)-\tau(x)^2=\sum_{b\ne1}x_b^2$ for $x=\sum_bx_b b$.
Positive variance gives at least two atoms in this finitely supported law. Consecutive atoms bound a gap, so \cite[Theorem~2.5]{BW} gives free indecomposability (see also the atom criterion in \cite{BV}). Classical infinite divisibility gives the same criterion since a bounded classically infinitely divisible law is a point mass. In infinite rank the semicircle and free Poisson laws supply freely infinitely divisible examples \cite{NS}.
\end{remark}

\section{Extensions and the Hom form}\label{sec:global}

The Hom inequality on all objects depends only on the abelian category. We first classify the extensions compatible with this inequality. Tensoring then forces a common multiplication rule for the resulting Jordan blocks.

\subsection{Extensions and Jordan blocks}\label{sec:ext}
An extension of $S$ by $T$ is a short exact sequence
$0\to T\to E\to S\to0$. Its class in $\Ext^1(S,T)$ vanishes exactly
when it splits. For example, an extension of the trivial $\Ga^d$-module by itself has matrices $\left(\begin{smallmatrix}1&\gamma(t)\\0&1\end{smallmatrix}\right)$, with $\gamma$ a regular additive map $\CC^d\to\CC$, hence a linear form. It splits exactly when $\gamma=0$. Choosing $\gamma(t)=t$ for $d=1$ gives the first nonsimple Jordan block.

\begin{lemma}\label{lem:les}
For an extension $\xi\colon0\to T\to E\to S'\to0$, applying Hom in
either variable gives exact sequences
\begin{align*}
 0\to\Hom(S,T)\to\Hom(S,E)\to\Hom(S,S')\xrightarrow{\partial}\Ext^1(S,T),\\
 0\to\Hom(S',S)\to\Hom(E,S)\to\Hom(T,S)\xrightarrow{\partial}\Ext^1(S',S).
\end{align*}
The first boundary takes pullbacks of $\xi$ and the second takes pushouts.
\end{lemma}
\begin{proof}
A map to $S'$ lifts to $E$ precisely when its pullback extension splits.
The pushout characterization is dual.
\end{proof}

\begin{example}[The Jordan block model]\label{thm:min}
A finite-dimensional rational $\Ga$-module has the form
$t\mapsto\exp(tN)$ for a nilpotent endomorphism $N$. In these coordinates
the tensor product is given by $N\otimes1+1\otimes M$ and the dual by
$-N^{\mathsf T}$, so every such module is self-dual. The indecomposables are
$J_r=\CC[u]/(u^r)$ with $N$ acting by $u$. An intertwiner $J_r\to J_s$
is determined by the image of $1$, which can be any element of $J_s$
annihilated by $u^r$. Hence
\[
 h(J_r,J_s)=\min(r,s)
 =\langle\mathbf1_{\{1,\ldots,r\}},\mathbf1_{\{1,\ldots,s\}}\rangle.
\]
On direct sums the Hom form is thus a scalar product of integer vectors.

\end{example}

\subsection{The global Cauchy--Schwarz property}

\begin{theorem}\label{thm:ext}
Let $\mathcal A$ be an essentially small Hom-finite $\CC$-linear abelian
length category. Put $h(A,B)=\dim_{\CC}\Hom(A,B)$.
The following conditions are equivalent.
\begin{enumerate}
\item $h(A,B)^2\le h(A,A)h(B,B)$ for all $A,B$.
\item $h(A,B)=h(B,A)$ for all $A,B$.
\item For simple objects $S,T$ one has $\Ext^1(S,T)=0$ if
$S\not\simeq T$, and $\dim\Ext^1(S,S)\le1$.
\end{enumerate}
Under these conditions every indecomposable is uniserial (its subobjects form a chain) and is determined
by its simple top $S$ and its length $r$. Write $J_{S,r}$ for that object when it
exists. All its composition factors are isomorphic to $S$, and
\begin{equation}\label{eq:min}
   h(J_{S,r},J_{T,s})=\delta_{S,T}\min(r,s).
\end{equation}
\end{theorem}

\begin{proof}
To see that (1) implies (2), let $e=h(A,A)$, $f=h(B,B)$,
$p=h(A,B)$ and $q=h(B,A)$. Take $A\ne0$ and apply (1) in both orders
to $A$ and $A^{\oplus n}\oplus B$. This bounds both $(ne+p)^2$ and
$(ne+q)^2$ by $e(n^2e+n(p+q)+f)$. Comparing the coefficients of $n$
as $n\to\infty$ gives $2p\le p+q$ and $2q\le p+q$.

Suppose (2) holds. A nonsplit extension
$0\to T\to E\to S\to0$ of distinct simples gives
$h(E,S)=1$ and $h(S,E)=0$. To exclude 2 linearly independent
self-extensions $\xi,\eta$ of $S$, use the construction in Gabriel's
argument \cite[proof of Theorem~1.7.1]{CK}. Take the extension
$0\to S\to E\to S^{\oplus2}\to0$ of class $(\xi,\eta)$.
The boundary $\Hom(S,S^{\oplus2})\to\Ext^1(S,S)$ sends
$(a,b)$ to $a\xi+b\eta$, so $h(S,E)=1$. In the second sequence of Lemma~\ref{lem:les}, the boundary
$\End(S)\to\Ext^1(S^{\oplus2},S)$ sends the identity to
$(\xi,\eta)$, so $h(E,S)=2$. Both dimension computations use
$\End(S)=\CC$. This proves (3).

Gabriel's criterion \cite[Theorem~1.7.1]{CK}, \cite[Theorems~7--8]{Eriksen} characterizes uniserial indecomposables by at most one incoming and one outgoing extension-quiver arrow at each simple $S$, counted by $\sum_T\dim\Ext^1(T,S)$ and $\sum_T\dim\Ext^1(S,T)$. Under (3) the subobjects of every indecomposable form a chain, and its top and length determine it \cite[Lemma~1.7.3]{CK}. Consecutive factors form a nonsplit extension, hence are isomorphic. Distinct simple blocks are Hom-orthogonal. Filter $\Hom(J_{S,r},J_{S,s})$ by the length $j$ of the target subobject. Each increment has dimension at most one, since maps to its simple quotient factor through the source's top. It is one for $1\le j\le\min(r,s)$, when that subobject is the length-$j$ quotient of the source, and zero thereafter by image length. This proves \eqref{eq:min}. By Krull--Schmidt, $h(A,B)=\langle u_A,u_B\rangle$ where $u_A(S,j)$ counts summands $J_{S,r}$ with $r\ge j$. Thus a block gives the indicator of its initial interval, and ordinary Cauchy--Schwarz gives (1).
\end{proof}

\begin{corollary}\label{cor:finite-global}
Every semisimple tensor category over $\CC$ satisfies global
Cauchy--Schwarz. A finite tensor category over $\CC$ satisfies it
exactly when it is a fusion category.
\end{corollary}
\begin{proof}
The first assertion follows from the scalar product of simple
multiplicity vectors. In a finite tensor category in characteristic zero,
$\Ext^1(\one,\one)=0$ by \cite[Theorem~4.4.1]{EGNO}.
Condition (3) then gives $\Ext^1(\one,S)=0$ for every simple $S$.
Induction on length gives $\Ext^1(\one,A)=0$ for every $A$.
Thus the unit is projective, which implies semisimplicity by
\cite[Corollary~4.2.13]{EGNO}.
\end{proof}

\begin{example}[Nilpotent endomorphisms in a semisimple category]\label{ex:nil}
Let $\mathcal S$ be a semisimple complex tensor category with any number
of simple isomorphism classes. Objects $(V,N)$, with $V\in\mathcal S$
and $N\in\End(V)$ nilpotent, form a tensor category. Morphisms intertwine
the endomorphisms, kernels and cokernels are taken in $\mathcal S$, and
the tensor and dual endomorphisms are $N\otimes1+1\otimes M$ and $-N^*$.
The evaluation and coevaluation of $\mathcal S$ intertwine these
endomorphisms. Decomposing $V$ into simple isotypic summands identifies
$N$ with one nilpotent matrix on each multiplicity space. Thus an
indecomposable is a Jordan block supported at one simple, and its Hom
dimensions obey \eqref{eq:min}. This gives global Cauchy--Schwarz, while
the length-2 block at the unit is a nonsplit self-extension.
For $\mathcal S=\Rep(R)$ with $R$ linearly reductive, the construction is
$\Rep(R\times\Ga)$. Taking $R$ cyclic of order $m$ gives pairwise
inequivalent nonsemisimple examples with $m$ simples, for every $m\ge1$.
\end{example}

\subsection{Jordan blocks of every length}
The following result complements Corollary~\ref{cor:finite-global}. Nonsemisimple categories satisfying global Cauchy--Schwarz are far from finite: they have no projective objects at all.
\begin{theorem}\label{thm:unbounded}
Suppose that the tensor category $\C$ satisfies global Cauchy--Schwarz and is not semisimple. Then for every simple $S$ and every $r\geq1$ the indecomposable $J_{S,r}$ exists, and $\C$ has no nonzero projective objects.
\end{theorem}

\begin{proof}
Let $\C_S$ be the full subcategory of objects whose composition factors are all isomorphic to $S$. By Theorem~\ref{thm:ext}, every object is a finite direct sum of objects of these subcategories, which are mutually Hom-orthogonal. If $\Ext^1(\one,\one)$ vanished, induction on length would give $\Ext^1(\one,-)=0$, making $\one$ projective and hence making $\C$ semisimple. Thus $\Ext^1(\one,\one)\neq0$.

Suppose there is a projective $P\neq0$. Then $P\otimes P^*$ is projective \cite[Proposition~4.2.12]{EGNO}, and the coevaluation $\one\to P\otimes P^*$ is nonzero. Hence some indecomposable summand of $P\otimes P^*$ receives a nonzero map from $\one$; it is a projective object $J_{\one,r}$. A longer $J_{\one,s}$, $s>r$, would have $J_{\one,r}$ as a quotient, and the surjection would split by projectivity, contradicting indecomposability. Thus every object of $\C_\one$ is a sum of $J_{\one,s}$ with $s\leq r$, each a quotient of $J_{\one,r}$, and $J_{\one,r}$ is a projective generator of $\C_\one$. Since $\otimes$ is exact and $\one\otimes\one=\one$, the subcategory $\C_\one$ is closed under tensor products, duals and extensions, so it is a finite tensor category, with a single simple. By \cite[Theorem~4.4.1]{EGNO} in $\C_\one$, we get $\Ext^1_{\C_\one}(\one,\one)=0$, while all extensions of $\one$ by $\one$ in $\C$ lie in $\C_\one$. This contradiction shows that there are no nonzero projectives.

Finally, if the lengths of the indecomposables of some $\C_S$ were bounded, a longest one $L$ would be projective. Indeed, given an epimorphism $M\to L$, decompose $M$ into indecomposables; the images of the summands lying in $\C_S$ are subobjects of the uniserial $L$, hence form a chain, and they generate $L$, so one summand maps onto $L$; by maximality it has the same length, so it maps isomorphically, and the epimorphism splits. This contradicts the previous paragraph. So lengths are unbounded in each $\C_S$, and since subobjects of uniserial objects are uniserial of every smaller length, every $J_{S,r}$ exists.
\end{proof}

\subsection{Tensor products of the Jordan blocks}
Assume global Cauchy--Schwarz. Put $J_{S,0}=0$ and set
\[
 e_{S,r}=[J_{S,r}]-[J_{S,r-1}]
\]
whenever $J_{S,r}$ exists. By \eqref{eq:min} these elements form an orthonormal
$\ZZ$-basis of the split Grothendieck group. The coordinate
$u_A(S,r)=h(e_{S,r},[A])$ counts the summands of $A$ supported at $S$ of length
at least $r$. In particular the coordinates of an object are $\geq0$ and,
for fixed $S$, nonincreasing in $r$.

\begin{lemma}\label{lem:layer-positive}
For every object $Y$, left and right multiplication by $[Y]$ have integer matrices
with entries $\geq0$ in the basis $e_{S,r}$. If $Y$ is self-dual, the matrices
are symmetric.
\end{lemma}
\begin{proof}
Choose a socle inclusion $T\hookrightarrow J_{T,s}$. The image of the
restriction $\Hom(J_{T,s},A)\longrightarrow\Hom(T,A)$ has dimension
$u_A(T,s)$: its kernel consists of the maps through the quotient
$J_{T,s-1}$, so the dimension follows from \eqref{eq:min}. A monomorphism
$A\hookrightarrow B$ induces an injection between these images, whence
$u_A(T,s)\leq u_B(T,s)$. Apply this to the inclusion
$Y\otimes J_{S,r-1}\hookrightarrow Y\otimes J_{S,r}$ given by exactness,
and likewise for right multiplication. Symmetry for $Y\simeq Y^*$ is
Frobenius reciprocity in an orthonormal basis.
\end{proof}

\begin{theorem}[Jordan tensor products]\label{thm:universal-jordan}
Suppose that $\C$ satisfies global Cauchy--Schwarz and is nonsemisimple.
Write $J_r=J_{\one,r}$, $J_0=0$. For all simples $S$ and $r,s\geq1$,
\[
 J_r\otimes S\simeq J_{S,r}\simeq S\otimes J_r,\qquad
 J_r\otimes J_s\simeq\bigoplus_{j=0}^{\min(r,s)-1}J_{r+s-1-2j}.
\]
In particular the split ring and the Hom form on the unit block are those of
$\Rep(\Ga)$.
\end{theorem}
\begin{proof}
All the indicated blocks exist by Theorem~\ref{thm:unbounded}. The object
$E=J_2$ is self-dual by uniqueness of the indecomposable of its length in
the unit block. Tensoring on the left by $E$ preserves each block and doubles
lengths, since it extends the identity functor by itself. In the basis
$e_{S,r}$ its matrix $M$ on each block is integral, with entries $\geq0$, and symmetric
by Lemma~\ref{lem:layer-positive}. Every column sums to $2$ because the
linear length functional takes value $1$ on every $e_{S,r}$.

On the unit block abbreviate $e_{\one,r}$ to $e_r$. We have $Me_1=e_1+e_2$.
Symmetry and the column sums leave only one unused term in $Me_2$, and
successively in each new column. More precisely, after a path through
$e_1,\ldots,e_r$, its next column is $Me_r=e_{r-1}+e_t$, where $t\geq r$:
all earlier columns are already filled. It is impossible that $t>r+1$,
because $M(e_1+\cdots+e_r)=[E\otimes J_r]$ would then have a positive
coordinate at $t$ and a zero coordinate at $r+1$. Thus the path continues
to $r+1$ or ends with a loop at $r$.

A terminal loop bounds the block lengths in powers of $E$ by $r$. The full subcategory of subquotients of finite direct sums of these powers is a tensor category by exactness and self-duality of $E$. Filter each finite sum of powers by its radical layers. There are at most $r$ layers, each semisimple. Intersecting the filtration with a subobject and taking its images in a quotient preserves this bound, since each new layer is a subquotient of a semisimple object. A longest block there is a projective generator: it surjects onto every smaller block and every epimorphism onto it splits by the uniserial argument of Theorem~\ref{thm:unbounded}. This finite tensor category would contain the nonsplit self-extension $E$ of its unit, contrary to \cite[Theorem~4.4.1]{EGNO}.

Consequently $Me_r=e_{r-1}+e_{r+1}$ for $r\geq2$. Summing the first $r$
columns gives $[E][J_r]=[J_{r+1}]+[J_{r-1}]$ for $r\geq1$. Each $[J_r]$
is therefore a monic polynomial of degree $r-1$ in $[E]$. In $\Rep(\Ga)$ the
block of length $r$ is the restriction of $\mathrm{Sym}^{r-1}(\CC^2)$
to the upper unitriangular subgroup of $SL_2$. Clebsch--Gordan gives the
same polynomials and the displayed products there. The blocks form a free
basis of each split ring, and \eqref{eq:min} identifies their Hom forms.

For a general $S$, the first column is $2e_{S,1}$ or $e_{S,1}+e_{S,2}$,
as $E\otimes S$ is an extension of $S$ by $S$. The same column argument
gives an infinite path in the natural order or a finite component at
$e_{S,1}$. If the latter has size $m$, then $[J_r\otimes S]$, a polynomial
in $M$ applied to $e_{S,1}$, is supported there for every $r$. Its
indecomposable summands all have length at most $m$, and its length is $r$,
so $h(J_r\otimes S,S)\geq r/m$. But adjunction gives
$h(J_r\otimes S,S)=h(J_r,S\otimes S^*)\leq\operatorname{length}(S\otimes S^*)$
by \eqref{eq:min}. Hence the path is infinite. Comparing it with the unit
block gives $[J_r\otimes S]=\sum_{i=1}^r e_{S,i}=[J_{S,r}]$.
Applying the result with reversed tensor product gives the other isomorphism.
\end{proof}

\section{Radical square zero and Taft algebras}\label{sec:hopf}
The categories $\Rep(T_2)$ and $\Rep(T_3)$ distinguish the two Cauchy--Schwarz conditions. In the first, every self-dual object has a law although global Cauchy--Schwarz fails. In the second, even the existence of the law depends on extension data. We begin with the common explanation for the positive examples. A finite-dimensional complex algebra is Morita equivalent to a basic algebra (all simples one-dimensional), with the same Loewy length $\min\{k:J^k=0\}$. Here $J$ is the Jacobson radical, the largest nilpotent ideal, and $J=0$ characterizes semisimplicity.

\begin{proposition}\label{prop:loewy}
Let $\C$ be a finite tensor category whose underlying basic algebra has radical of square zero. Then every self-dual object of $\C$ has a unique, finitely supported, representing measure.
\end{proposition}

\begin{proof}
In a nonsemisimple block the indecomposable projectives $P_i$ with tops $S_i$ are injective with simple socles \cite[Proposition~6.1.3 and Remark~6.1.5]{EGNO}. None is simple, as a simple projective-injective forms its own block. Thus $J^2=0$ makes $JP_i$ nonzero semisimple, equal to the simple socle $S_{\sigma(i)}$. Since the $P_i$ exhaust the indecomposable injectives, $\sigma$ permutes the simples. The projectives and injectives are all uniserial of length 2. The block algebra is therefore Nakayama and its indecomposables are precisely $S_i,P_i$ \cite[Theorem~VI.2.1]{ARS}.

Let $v=\sum_ia_i[S_i]+\sum_ib_i[P_i]$ with $a,b$ real. The top and socle calculations and the projectivity of $P_i$ give
\[
\begin{aligned}
 h(S_i,S_j)&=\delta_{ij}, & h(P_i,S_j)&=\delta_{ij},\\
 h(S_j,P_i)&=\delta_{j,\sigma(i)}, & h(P_i,P_j)&=[P_j:S_i]=\delta_{ij}+\delta_{i,\sigma(j)} .
\end{aligned}
\] Let $T$ be the permutation matrix $Te_i=e_{\sigma(i)}$, and $w=(I+T)b$. Then
\[
 h(v,v)=\|a\|^2+\langle a,b+Tb\rangle+\|b\|^2+\langle b,Tb\rangle
        =\|a\|^2+\langle a,w\rangle+\tfrac12\|w\|^2
        =\|a+w/2\|^2+\tfrac14\|w\|^2\geq0,
\]
using $\frac12\|w\|^2=\|b\|^2+\langle b,Tb\rangle$, since $T$ is orthogonal. Semisimple blocks contribute sums of squares, and different blocks are Hom-orthogonal. So $h(v,v)\geq0$ on the whole split Grothendieck ring, and Theorem~\ref{thm:main} gives existence. As there are finitely many indecomposables, the ranks of the Hankel matrices are bounded by their number, giving finite support and uniqueness.
\end{proof}

\begin{example}[A family of Hopf algebras with one block]\label{ex:exterior}
For $r\geq1$, let $\zeta$ be a primitive $2r$th root of unity, and let $H_r$ be the Hopf algebra with generators $g,x$ and
\[
 g^{2r}=1,\qquad gx=\zeta xg,\qquad x^2=0,
 \qquad \Delta(g)=g\otimes g,\qquad
 \Delta(x)=x\otimes1+g^r\otimes x .
\]
The compatibility of $\Delta$ with $x^2=0$ comes from $g^rx=-xg^r$. This is the rank 1 construction of \cite[\S2]{WHLZ} for the datum $(\langle g\rangle,\chi,g^r,0)$, with $\chi(g)=\zeta^{-1}$ and $x=yg^r$ in their notation. Its basis is $g^j,xg^j$ with $0\leq j<2r$, so $\dim H_r=4r$. The ideal $(x)$ is nonzero with square zero, and the quotient is the semisimple algebra $\CC[\ZZ/2r]$, so $(x)$ is the Jacobson radical. The $2r$ simple modules are the characters with $g$-weights $\zeta^j$, and $x$ joins weight $j$ to weight $j+1$ modulo $2r$, so the basic algebra is the path algebra of the oriented $2r$-cycle, modulo paths of length 2. In particular $\Rep(H_r)$ has a single block, and these categories are pairwise inequivalent, by their numbers of simples. Proposition~\ref{prop:loewy} applies to every self-dual object. Corollary~\ref{cor:finite-global} excludes global Cauchy--Schwarz. For $r=1$ we obtain Sweedler's algebra.
\end{example}

For $m\geq2$ and $q$ a primitive $m$th root of unity, the \emph{Taft algebra} $T_m=H_m(q)$ is the Hopf algebra with generators $g,x$, relations and structure maps
\[
  g^m=1,\quad x^m=0,\quad xg=qgx,\qquad
  \Delta(g)=g\otimes g,\quad \Delta(x)=1\otimes x+x\otimes g,
\]
with $\varepsilon(g)=1$, $\varepsilon(x)=0$, $S(g)=g^{-1}$, $S(x)=-xg^{-1}$. We follow \cite[\S1]{CVZ}, where $x$ is denoted $h$. It has basis $g^ix^j$, $0\leq i,j<m$, so $\dim T_m=m^2$. Here $T_2$ is Sweedler's 4-dimensional Hopf algebra.

For $1\leq\ell\leq m$ and $r\in\ZZ/m$, let $M(\ell,r)$ have basis $v_0,\ldots,v_{\ell-1}$ with
\[
 gv_k=q^{r-k}v_k,\qquad xv_k=v_{k+1}\quad(v_\ell=0).
\]
The operator $x$ forms a single Jordan chain and lowers the $g$-weight by $1$. The module $M(\ell,r)$ is uniserial of length $\ell$, with top of weight $r$ and socle of weight $r-\ell+1$, and these $m^2$ modules are all the indecomposables \cite[Lemmas~2.2--2.3, Corollary~2.4 and Theorem~2.5]{CVZ}. The simples are the $M(1,r)$, the characters $g\mapsto q^r$, $x\mapsto0$, and the projectives are the $P_r=M(m,r)$. An invariant is a vector with $gv=v$ and $xv=0$, so it must be a multiple of the socle vector $v_{\ell-1}$, and
\[
 \dim\Hom(\one,M(\ell,r))=\begin{cases}1&\text{if }r\equiv\ell-1\pmod m,\\0&\text{otherwise}.\end{cases}
\]
Dualizing reverses the chain and negates the weights, so it interchanges top and socle, $M(\ell,r)^*\simeq M(\ell,\ell-1-r)$. On $V\otimes W$ the coproduct makes $g$ act by $g\otimes g$ and $x$ by $1\otimes x+x\otimes g$. We use the tensor product formulas of \cite{CVZ} below.

\begin{proposition}\label{prop:sweedler}
In $\Rep(T_2)$ every self-dual object has a representing measure.
Writing $s$ for the nontrivial simple and $R=P_0\oplus P_1$, a
self-dual object has the form $X=\one^{\oplus a}\oplus s^{\oplus b}
\oplus R^{\oplus c}$ and its measure is
\[
 \tfrac12\delta_{a-b}+\tfrac14\delta_{a+b}
                 +\tfrac14\delta_{a+b+4c}.
\]
\end{proposition}
\begin{proof}
Duality swaps the 2 projectives. The rules
$s^2=1$, $sR=Rs=R$, $R^2=4R$ in the split ring follow from
\cite[Lemma~3.1 and Proposition~3.2]{CVZ}. Its subalgebra on $1,s,R$
has orthogonal idempotents $(1-s)/2$, $(1+s)/2-R/4$, $R/4$.
The invariant dimension functional takes values $1,0,1$ on $1,s,R$,
hence $1/2,1/4,1/4$ on these idempotents. Multiplication by $[X]$
has the corresponding eigenvalues in the formula.
\end{proof}

For $X=R$ this gives $(3\delta_0+\delta_4)/4$. To see directly why
global Cauchy--Schwarz fails in the same category, take $A=P_0$,
$B=P_0\oplus\one$. The top of $P_0$ is $\one$ and its socle is $s$.
Thus $h(A,B)=2$, $h(A,A)=1$, $h(B,B)=3$, and $4>3$.

\begin{proposition}\label{prop:taft-projectives}
Let $Y=\bigoplus_rP_r^{\oplus a_r}$ be a nonzero projective $T_m$-module, and put $k=\sum_ra_r$ and $d=mk=\dim Y$. Then $(c_n(Y))$ is a Hamburger sequence if and only if $ma_{m-1}=k$, and in this case the measure is $(1-m^{-2})\delta_0+m^{-2}\delta_d$.
\end{proposition}

\begin{proof}
For the regular module $R_m=\bigoplus_rP_r$, \cite[Proposition~3.2]{CVZ} gives $P_i\otimes P_j\simeq R_m$. Thus $[Y]^2=k^2[R_m]$, $[Y][R_m]=d[R_m]$ and $[Y]^n=k^2d^{n-2}[R_m]$ for $n\ge2$. Only $P_{m-1}$ has an invariant line, so $c_1=a_{m-1}$ and $c_n=k^2d^{n-2}$ for $n\ge2$. Since $c_4-2dc_3+d^2c_2=0$, the square $t^2(t-d)^2$ vanishes on the support of any representing measure. On $\{0,d\}$ the second moment fixes the mass at $d$ to $k^2/d^2=m^{-2}$. All moments agree precisely when $a_{m-1}=d/m^2=k/m$.
\end{proof}

For every $m\geq3$ the self-dual projective $Y=P_0\oplus P_{m-1}$ violates the condition, since $ma_{m-1}=m\neq2=k$. For a common witness, take $A=\one\oplus Y^{\otimes2}$ and $B=Y$. Then $c_1=1$, $c_2=4$, $c_3=8m$, $c_4=16m^2$, and \eqref{eq:bridge} gives $h(A,B)=c_1+c_3$, $h(A,A)=c_0+2c_2+c_4$ and $h(B,B)=c_2$, so
\[
 h(A,B)^2-h(A,A)h(B,B)=(1+8m)^2-4(9+16m^2)=16m-35>0 .
\]
The categories $\Rep(T_m)$ are pairwise inequivalent, having $m$ simples. Together with the radical-square-zero case, this shows that $m=2$ is the only member of the Taft family for which every self-dual module has a representing measure.

\begin{theorem}\label{thm:taft-failure}
In the representation category of the 9-dimensional Taft algebra $T_3$, the self-dual indecomposable $X=M(2,2)$ has no Hamburger representing measure.
\end{theorem}

\begin{proof}
Duality fixes $X$ since $1-2\equiv2\pmod3$. Specializing \cite[Propositions~3.2 and~3.4]{CVZ},
\[
 X^{\otimes2}\simeq\one\oplus P_1,\qquad
 X\otimes P_r\simeq P_{r+1}\oplus P_{r+2}.
\]
So the unique nonprojective summand of $X^{\otimes n}$ is $\one$ for $n$ even and $X$ for $n$ odd, and these have $1$ and $0$ invariants respectively, $X$ having top of weight 2 and socle of weight 1. Among the projectives only $P_2$ has an invariant line. Iterating the 2 tensor rules yields
\[
\begin{array}{c|rrrrrrr}
 n                         &0&1&2&3&4&5&6\\ \hline
 \text{multiplicity of }P_0&0&0&0&1&1&4&6\\
 \text{multiplicity of }P_1&0&0&1&0&3&2&9\\
 \text{multiplicity of }P_2&0&0&0&1&1&4&6\\ \hline
 c_n                       &1&0&1&1&2&4&7
\end{array}
\]
and the Hankel determinant $\det(c_{i+j})_{0\leq i,j\leq3}$ equals $-1$. In fact the values through degree five already preclude a measure. For $q=t^2-t-1$ one has $L_X(q^2)=0$, so any representing measure would be supported on its zeros. Yet $L_X(t^3q)=c_5-c_4-c_3=1$. The values through degree four agree with the Fibonacci law of Example~\ref{ex:fusion}.
\end{proof}

Nevertheless $M(1,1)\oplus M(1,2)$ has the same ordinary Grothendieck
class as $M(2,2)$. Its powers are semisimple and its law is
$(\delta_2+2\delta_{-1})/3$, the character law for the cyclic group of
order 3. Thus the ordinary class of a self-dual object need not
decide even the existence of a moment measure.

In these examples $h(\one,-)$ reads the unit in the socle. For the length invariant in $\Rep(G)$ with $G$ finite, the corresponding positivity condition is pointwise nonnegativity of the sum of irreducible characters \cite[Remark~2.7(1)]{FFK}. The Schur-product obstruction for fusion bialgebras concerns unitary categorification \cite[Propositions~3.3, 8.3 and Corollary~8.5]{LPW}. The Taft pair exhibits a different dependence, on the object realizing a fixed ordinary class. For projective objects, Theorem~\ref{thm:finite-test} recovers the socle counts from finite data.

\section{Invariants of affine algebraic groups}\label{sec:groups}
All representations of complex affine algebraic groups in this section are finite-dimensional and rational. We distinguish two properties of $\Rep(G)$. The global Hom inequality of Theorem~\ref{thm:ext} tests all pairs of modules. Write $P(G)$ for the existence of a Hamburger law for every self-dual module, equivalently the inequality on sums of its powers (Theorem~\ref{thm:main}). Inflation along a quotient preserves duals and invariant dimensions, so a counterexample for a quotient gives one for $G$.

Write $G=U\rtimes R$ by Mostow's decomposition, with $U=R_u(G^\circ)$ and $R$ linearly reductive, also for disconnected $G$ \cite[Proposition~5.4.1 and its proof]{Conrad}. For unipotent $U$, the characteristic-zero exponential correspondence identifies $\Hom(U,\Ga)$ with $\Lie(U)_{\rm ab}^*$, equivariantly for $R$ \cite[Theorem~14.37]{Milne}.

\subsection{Classifying the global property}
\begin{theorem}\label{thm:global-groups}
Let $G$ be a complex affine algebraic group of finite type, possibly disconnected. Global Cauchy--Schwarz holds in $\Rep(G)$ if and only if $G\simeq R$ or $G\simeq R\times\Ga$ with $R$ linearly reductive.
\end{theorem}

\begin{proof}
Put $V=\Lie(U)_{\mathrm{ab}}$. For each simple $R$-module $W$, inflated to $G$, we have an injection $\Hom_R(V,W)\hookrightarrow\Ext^1_G(\one,W)$. Indeed $f$ integrates to an $R$-equivariant $\bar f:U\to(W,+)$, giving the action $(u,r)(w,t)=(rw+t\bar f(u),t)$ on $W\oplus\CC$. Baer sum adds the $f$'s, and a splitting would give a $U$-fixed lift of $1$, forcing $f=0$.

Under global Cauchy--Schwarz, Theorem~\ref{thm:ext} and complete reducibility give $V=0$ or $\one$. Writing $\mathfrak u=\Lie(U)=\CC y+\gamma_2$ for its lower central series, we find $\gamma_2=[\CC y+\gamma_2,\CC y+\gamma_2]\subset\gamma_3$. Nilpotence now gives $\gamma_2=0$, so $U=1$ or $U=\Ga$ with trivial $R$-action. Conversely $\Rep(R)$ is semisimple, and for $R\times\Ga$ one applies \eqref{eq:min} on each reductive isotypic component.
\end{proof}

\subsection{The tilted torus law for one additive parameter}
Identify $\Ga$ with the upper unitriangular subgroup of $SL_2$, and put
$V_d=\mathrm{Sym}^d(\CC^2)$. Its restriction is a single Jordan block
$J_{d+1}$. By the Jordan decomposition every rational $\Ga$-module $X$
therefore has a lift $\widetilde X$ to $SL_2$, unique up to isomorphism.
Restriction commutes with tensor product, so $\widetilde X^{\otimes n}$
restricts to $X^{\otimes n}$.

\begin{proposition}\label{prop:additive-law}
Let $\chi(\theta)$ be the character of the lift $\widetilde X$ at
$\operatorname{diag}(e^{i\theta},e^{-i\theta})$. The unique representing
measure for $X$ is
\begin{equation}\label{eq:ga-law}
   \chi_*\!\left((1+\cos\theta)\,\frac{d\theta}{2\pi}\right),
   \qquad 0\le\theta<2\pi.
\end{equation}
\end{proposition}
\begin{proof}
Each $V_d$ has one invariant line upon restriction and torus character $z^d+z^{d-2}+\cdots+z^{-d}$.
Integration against the displayed density extracts the constant
coefficient plus half the sum of the coefficients of $z$ and $z^{-1}$.
This gives $1$ for both parities of $d$. Apply the identity to every
irreducible summand of $\widetilde X^{\otimes n}$. The resulting law is
real and supported on $[-\dim X,\dim X]$, which also ensures uniqueness.
\end{proof}

For example $J_2$ has moments $\binom n{\lfloor n/2\rfloor}$
\cite[Proposition~5.3]{DG} and density
$(1+x/2)/(\pi\sqrt{4-x^2})$ on $(-2,2)$. Its even moments agree with
those of the arcsine law. For comparison, Haar integration on $SU(2)$
uses the torus density $1-\cos(2\theta)$ in the same normalization.

\subsection{Recovering the Jordan type from the even moments}
Zhou reconstructs a faithful self-dual representation of the circle from its invariant dimensions by studying the character near its maximum \cite[Theorem~1.1 and Corollary~1.2]{Zhou}. For $\Ga$ the even invariant dimensions give the distribution of the \emph{square} of the lifted character. The tilt in \eqref{eq:ga-law} allows us to distinguish the one-arc and two-arc tails that arise from the parities of the Jordan blocks.

\begin{theorem}\label{thm:reconstruction}
Finite-dimensional rational $\Ga$-modules with the same invariant dimensions in every even tensor degree are isomorphic. Equivalently, the law of $X\otimes X$ determines the Jordan blocks of $X$.
\end{theorem}

\begin{proof}
Write the lifted character as $f(\theta)=\sum_{j=-M}^M m_j e^{ij\theta}$, with $m_j=m_{-j}\geq0$. By Proposition~\ref{prop:additive-law}, the even moments determine the compactly supported law
\[
 \nu_f=(f^2)_*\!\left((1+\cos\theta)\frac{d\theta}{2\pi}\right).
\]
The angular density has full support, so the right endpoint of $\supp\nu_f$ is $d^2$, where $d=f(0)=\dim X$. A point mass forces $f$ constant and hence trivial action, including $X=0$. Assume from now on that $f$ is nonconstant. Then $f''(0)=-\sum_jj^2m_j<0$.

An even-sized block contributes only odd weights, including $\pm1$. An odd-sized block contributes weight $0$ and only even weights, including $\pm2$ when its size is at least $3$. If all block sizes have the same parity, $|f|$ reaches $d$ exactly at $0$ and $\pi$, and $f(\theta+\pi)=\pm f(\theta)$. If both parities occur, weights $0$ and $1$ occur, and $|f|=d$ only at $0$ modulo $2\pi$.

For sufficiently small $\varepsilon>0$, let $u>0$ satisfy $f(u)=d-\varepsilon$. The inverse image of $[(d-\varepsilon)^2,d^2]$ consists of an arc of half-width $u$ around each of the indicated maxima. The arc at $0$ has tilted mass $(u+\sin u)/\pi$, and the arc at $\pi$, when present, has mass $(u-\sin u)/\pi$. Consequently
\[
 \nu_f\bigl([(d-\varepsilon)^2,d^2]\bigr)=
 \begin{cases}
 2u/\pi,&\text{all block sizes have the same parity},\\
 (u+\sin u)/\pi,&\text{both parities occur}.
 \end{cases}
\]
Equal laws in the same case therefore give equal local inverse branches, hence identical trigonometric polynomials. For equality across the two cases, write $f_P$ for the character whose block sizes have the same parity and $f_M$ for the mixed character. Equality of the masses would give
\[
 f_M(u)=f_P\bigl((u+\sin u)/2\bigr).
\]
Both sides are entire, so the identity extends to $u=it$, $t>0$. The left side is $O(e^{M_Mt})$, whereas the positive terms on the right include $m_{M_P}e^{M_P(t+\sinh t)/2}$, where $M_M,M_P>0$ are the largest weights. This is impossible. Finally, the recovered character determines the number $m_k-m_{k+2}$ of blocks of size $k+1$.
\end{proof}

\begin{corollary}[Finite reconstruction in the unit block]\label{cor:finite-reconstruction}
Let $\C$ be a nonsemisimple complex tensor category satisfying global Cauchy--Schwarz. For an object $X$ all of whose composition factors are $\one$, the even invariant dimensions $c_2(X),c_4(X),\ldots$ determine the multiplicity of every $J_r$ in $X$. There is an algorithm that stops after finitely many of these dimensions, without a length bound as input. For each fixed $M$, the number of even moments needed can be bounded effectively and uniformly on the class $c_2(X)\le M$.
\end{corollary}

\begin{proof}
By Theorem~\ref{thm:ext} each $J_r$ is self-dual, so Theorem~\ref{thm:main} applies. Theorem~\ref{thm:universal-jordan} identifies the unit, products and Hom form with those of $\Rep(\Ga)$. Hence $X=\bigoplus_{i=1}^kJ_{r_i}$ has the unique compact tilted-torus law of Proposition~\ref{prop:additive-law}, which distinguishes partitions by Theorem~\ref{thm:reconstruction}.

By \eqref{eq:min}, $c_2(X)=h(X,X)=\sum_{i,j}\min(r_i,r_j)\ge\sum_i r_i$. List all partitions of total at most $c_2$ with the given value of $\sum_{i,j}\min(r_i,r_j)$, allowing the empty partition for $X=0$. For each candidate $\lambda=(r_i)$ put $f_\lambda(z)=\sum_i(z^{r_i-1}+z^{r_i-3}+\cdots+z^{1-r_i})$. Compare successive $c_{2n}$ with $[z^0]f_\lambda^{2n}+[z^1]f_\lambda^{2n}$ by the torus formula. Each false candidate differs at a finite stage by Theorem~\ref{thm:reconstruction}, so this finite test terminates. Comparing all partitions of size at most $M$ until their even-moment vectors are distinct computes the uniform bound.
\end{proof}

For example, $J_1\oplus J_5$ and $J_2^{\oplus2}$ both have $c_1=2$ and $c_2=8$. Their Laurent polynomials are $2+z^4+z^2+z^{-2}+z^{-4}$ and $2(z+z^{-1})$. The formula above gives third moments $38$ and $24$, and fourth moments $196$ and $96$, respectively. For an indecomposable in the unit block the first 2 moments already identify it, since $c_1=1$ and $c_2=r$ characterize $J_r$.

\subsection{Two additive parameters and the affine group}
The preceding law determines every representation of $\Ga$. The next
example shows how even its existence can fail on passing to $\Ga^2$.

\begin{proposition}\label{prop:plane}
The group $\Ga^2$ has a faithful self-dual 6-dimensional module $X$ with $(c_0,c_1,c_2,c_3,c_4)=(1,3,12,57,270)$. This sequence has no representing measure on $\RR$.
\end{proposition}

\begin{proof}
On $\Ga^3$ take the sum of 3 copies of $J_2$, one for each coordinate, and restrict to the plane $(s,t,s+t)$. Each summand preserves its determinant form, so is self-dual, and the first 2 detect both parameters. To compute the effect of restriction, realize $V_a\boxtimes V_b\boxtimes V_c$ as the polynomials in $u,v,w$ of separate degrees at most $a,b,c$, with the 3 additive groups acting by translation of the variables. The polynomials fixed by the plane are the polynomials in $w-u-v$, so the fixed space has dimension $1+\min(a,b,c)$, compared with $1$ for 3 independent translations.

Expand a tensor power of the sum into words in the 3 summands. For independent translations the dimension is $\sum_{i+j+k=n}\binom{n}{i,j,k}b_ib_jb_k$, with $b_n=c_n(J_2)$. Restriction adds 1 invariant for each of the 6 words using all 3 letters at length 3. At length 4, there are 36 words of type $(2,1,1)$, each gaining 1 invariant by $V_1^{\otimes2}=V_2\oplus V_0$. The other words use at most 2 of the parameters, and keep their invariants. This gives
\[
\begin{array}{c|rrrrr}
n&0&1&2&3&4\\ \hline
\Ga^3&1&3&12&51&234\\
\Ga^2&1&3&12&57&270
\end{array}
\]
For $A=X$ and $B=X^{\otimes2}$ inequality \eqref{eq:cs} reads $c_3^2\leq c_2c_4$, while $57^2=3249>3240=12\cdot270$. The verification script computes the $\Ga^2$ row from simultaneous Lie algebra kernels and the $\Ga^3$ row from the multinomial sum.
\end{proof}

Inflating $X$ along a coordinate projection $\Ga^d\twoheadrightarrow\Ga^2$ gives the same violating pair for every $d\geq2$. These categories are pairwise inequivalent, since $\Ext^1(\one,\one)\simeq\Hom(\Ga^d,\Ga)$ has dimension $d$, by the self-extension calculation at the start of Section~\ref{sec:ext}.

\begin{example}[The affine group of the line]\label{ex:affine}
Let $B\subset SL_2$ be the upper triangular Borel subgroup and $D$ its upper left entry character. The self-dual $B$-module $Y=V_3|_B\otimes(D^5\oplus D^{-5})$ descends to $B/\{\pm I\}\simeq\Ga\rtimes\Gm$, with $\Gm$ acting by scaling. The invariants of $V_m|_B\otimes D^j$ form a line if $m=-j$ and vanish otherwise, since $\Ga$ fixes only the highest weight line, of weight $m$. Let $a_{n,m}$ be the multiplicity of $V_m$ in $V_3^{\otimes n}$, with $a_{n,m}=0$ for $m<0$. Choosing $D^5$ in $i$ of the $n$ factors,
\[
 c_n(Y)=\sum_{i=0}^n\binom ni a_{n,5(n-2i)} .
\]
By Clebsch--Gordan \cite[Example~4.9.3]{EGNO}, $V_3^{\otimes2}=V_0\oplus V_2\oplus V_4\oplus V_6$, so $a_{2,0}=1$, $a_{3,5}=3$, $a_{4,0}=4$ and $a_{4,10}=3$, the last one from the pairs of intermediate weights $(4,6),(6,4),(6,6)$. Hence $(c_0,\ldots,c_4)=(1,0,2,9,36)$, and $c_3^2\leq c_2c_4$ fails, since $81>72$.
\end{example}

\begin{theorem}\label{thm:solvable}
For a connected solvable complex affine algebraic group $G$, the property $P(G)$ holds exactly when $G$ is a torus, or the product of a torus with $\Ga$. The trivial torus is allowed.
\end{theorem}

\begin{proof}
These groups satisfy global Cauchy--Schwarz by Theorem~\ref{thm:global-groups}, hence $P(G)$. Conversely, write $G=U\rtimes T$ with $U$ the unipotent radical and $T$ a maximal torus \cite[Theorem~16.33(d)]{Milne}, and $\mathfrak u=\Lie U$. If $T$ acts nontrivially on $\mathfrak u/[\mathfrak u,\mathfrak u]$, choose a $1$-dimensional quotient of nontrivial weight $\lambda$. The exponential equivalence integrates it to a $T$-equivariant surjection $U\to\Ga$, and together with $\lambda\colon T\twoheadrightarrow\Gm$ this gives a quotient of $G$ onto the affine group of Example~\ref{ex:affine}, so $P(G)$ fails. Otherwise $T$ acts trivially on all the quotients $\gamma_j/\gamma_{j+1}$ of the lower central series of $\mathfrak u$, which are spanned by iterated brackets of elements of the first one. By complete reducibility $T$ acts trivially on $\mathfrak u$, hence on $U$, and $G=U\times T$. If $\dim U\geq2$, the argument in the proof of Theorem~\ref{thm:global-groups} shows $\dim\mathfrak u/[\mathfrak u,\mathfrak u]\geq2$, so $G$ surjects onto $\Ga^2$, and Proposition~\ref{prop:plane} shows that $P(G)$ fails.
\end{proof}

\section{Stabilizers of highest weight vectors}\label{sec:stabilizers}
For a connected complex affine group, does $P(G)$ force $G$ to be reductive or the direct product of a reductive group with $\Ga$? The converse follows from Theorem~\ref{thm:global-groups}. To study the open direction, write $G=U\rtimes L$ with nontrivial unipotent radical $U$ and reductive Levi factor $L$. If $L$ acts trivially on $\Lie(U)_{\rm ab}$, the proof of Theorem~\ref{thm:solvable} gives either $G=L\times\Ga$ or a quotient $\Ga^2$. A one-dimensional quotient of $\Lie(U)_{\rm ab}$ with a nontrivial character gives a quotient $\Ga\rtimes\Gm$. In the remaining case choose an irreducible quotient $W$ of dimension at least two on which $L$ acts through $\bar L\subset GL(W)$. The induced surjection $G\to W\rtimes\bar L$ lifts each counterexample for $P(W\rtimes\bar L)$ to $G$. It therefore suffices for the proposed connected classification to find such a counterexample for every faithful irreducible action in dimension at least two.

Several such semidirect products occur as stabilizers of highest weight vectors. Let $H$ be a connected reductive complex group, with a Borel subgroup $B$ and a maximal torus $T\subset B$, and let $v$ be a highest weight vector, of weight $\lambda$, in an irreducible $H$-module. Let $P=U\rtimes L$ be the parabolic subgroup fixing the line $\mathbb Cv$, with unipotent radical $U$ and Levi factor $L\supset T$, and let $L_\lambda$ be the identity component of the kernel of $\lambda$ on $L$. The identity component of the stabilizer of $v$ is then
\[
   G=U\rtimes L_\lambda .
\]
We write $V_\mu$ for the irreducible $H$-module of highest weight $\mu$, and $V^L_\mu$ for the irreducible $L$-module of highest weight $\mu$.

\subsection{Invariants of the restricted modules}
The invariants of $G$ only see the highest weight vectors of $H$ and the Levi factor.

\begin{lemma}\label{lem:levi}
The space $V_\mu^U$ is the irreducible $L$-module $V^L_\mu$, and its restriction to $L_\lambda$ is irreducible. Consequently, for $Y\in\Rep(H)$ and $W\in\Rep(L_\lambda)$, the latter inflated to $G$,
\[
   \dim\,(Y|_G\otimes W)^G=\sum_\mu\,[Y:V_\mu]\,\bigl[W^*:V^L_\mu|_{L_\lambda}\bigr] .
\]
In particular $\dim V_\mu^G$ is $1$ if $\mu$ is a rational multiple of $\lambda$, and is $0$ otherwise.
\end{lemma}

\begin{proof}
The group $L$ normalizes $U$, so it acts on $V_\mu^U$. A vector of $V_\mu^U$ fixed by the unipotent radical of the Borel subgroup $B\cap L$ of $L$ is fixed by the unipotent radical of $B$, so it lies on the highest weight line of $V_\mu$. Thus the $L$-module $V_\mu^U$ has a single line of highest weight vectors, of weight $\mu$, and by complete reducibility it is irreducible. Since $L$ is generated by $[L,L]\subset L_\lambda$ and its central torus, which acts by scalars, the restriction to $L_\lambda$ stays irreducible. As $U$ acts trivially on $W$, the invariants of $V_\mu\otimes W$ are $(V_\mu^U\otimes W)^{L_\lambda}$, which gives the formula. Finally $L_\lambda$ is normal in $L$, so the $L_\lambda$-invariants of $V^L_\mu$ form an $L$-submodule, which is nonzero exactly when $V^L_\mu$ is a character of $L$ trivial on $L_\lambda$. Such a character vanishes on the maximal torus of $[L,L]$, and is a multiple of $\lambda$ on the central torus, so $\mu$ is a rational multiple of $\lambda$, and conversely.
\end{proof}

We apply the lemma to the vector stabilizers $\CC^N\rtimes SL_N\subset SL_{N+1}$ and the complex Euclidean groups $\CC^m\rtimes SO_m\subset SO_{m+2}$, together with their spin covers. The contrast is already visible in dimension three. Many modules of the Euclidean group inherit a law from $Spin_4$, while its affine, dual affine and adjoint modules taken together violate a Hankel inequality.

\subsection{Vector stabilizers}
Let $M=\CC^{N+1}$ with basis $e_1,\ldots,e_{N+1}$. The stabilizer of $e_1$ in $SL_{N+1}$ consists of the block matrices $\left(\begin{smallmatrix}1&u\\0&A\end{smallmatrix}\right)$ with $u\in\CC^N$ a row and $A\in SL_N$, so it is $G_N=\CC^N\rtimes SL_N$, with unipotent radical $U_N=\{A=1\}$ and Levi factor $SL_N$, acting on $V=M/\CC e_1$. For $N=1$ this is $\Ga\subset SL_2$. This is the case $H=SL_{N+1}$, $\lambda=\omega_1$ of Lemma~\ref{lem:levi}, with $L$ the block diagonal $GL_N$ and $L_\lambda=SL_N$. We write $\bar\mu$ for the restriction of $\mu$ to the diagonal torus of $SL_N$, which is the highest weight of $V^L_\mu|_{SL_N}$. The multiples of $\omega_1$ are the highest weights of the modules $\mathrm{Sym}^k(M)$, so $\dim V_\mu^{G_N}$ is $1$ if $V_\mu\simeq\mathrm{Sym}^k(M)$ for some $k\geq0$, and is $0$ otherwise.

For $N=1$ this says that each $V_d$ has $1$ invariant line, as in the proof of Proposition~\ref{prop:additive-law}. For $N=2$ there is again a law for the restricted modules, now carried by the Levi factor.

\begin{proposition}\label{prop:tilted-haar}
Let $Y$ be a self-dual $SL_3$-module, with character $\chi_Y$. The invariant dimensions of $Y|_{G_2}$ are the moments of the pushforward by $\chi_Y$ of the positive measure
\[
     \Bigl(1-\frac{\chi_V(g)}2\Bigr)dg
\]
on $SU(2)\subset SU(3)$, where $dg$ is the Haar measure and $\chi_V$ the character of $V=\CC^2$. In torus coordinates, the density is $\pi^{-1}\sin^2\theta\,(1-\cos\theta)$ on $[0,2\pi)$.
\end{proposition}

\begin{proof}
Write $V(a,b)$ for the irreducible $SL_3$-module with highest weight $a\omega_1+b\omega_2$, so that $V(a,b)^*=V(b,a)$ and $\mathrm{Sym}^k(M)=V(k,0)$. For $Z=\bigoplus m_{ab}V(a,b)$ self-dual we have $m_{ab}=m_{ba}$, and Lemma~\ref{lem:levi} gives $\dim Z^{G_2}=\sum m_{ab}[b=0]=\sum m_{ab}([a=0]+[b=0])/2$. On the other hand, by the branching rule from $GL_3$ to $GL_2$, the restriction of $V(a,b)$ to $SL_2$ is the sum of the modules $V_{\mu_1-\mu_2}$ over the integers $a+b\geq\mu_1\geq b\geq\mu_2\geq0$. It contains $V_0$ once, and $V_1$ exactly $[a>0]+[b>0]$ times. Hence
\[
 \int_{SU(2)}\chi_Z\Bigl(1-\frac{\chi_V}2\Bigr)dg=\sum_{a,b}m_{ab}\Bigl(1-\frac{[a>0]+[b>0]}2\Bigr)=\dim Z^{G_2}.
\]
Apply this to $Z=Y^{\otimes n}$. The density is nonnegative since $\chi_V\leq2$.
\end{proof}

For instance $Y=M\oplus M^*$ has invariant dimensions $1,1,4,12,48,192,832,\ldots$, the moments of $2+4\cos\theta$ under this density. This positivity does not extend to the modules which involve the Levi factor separately, and the answer to the general question is negative.

\begin{theorem}\label{thm:levi}
The property $P(G)$ fails for $G=\CC^2\rtimes SL_2$ and for $G=\CC^3\rtimes SL_3$. More precisely:
\begin{enumerate}
\item The self-dual $18$-dimensional $G_2$-module $X=M\otimes M^*\otimes V$ has
$(c_0,\ldots,c_4)=(1,1,4,46,356)$,
and $c_3^2\leq c_2c_4$ fails, since $2116>1424$.
\item The self-dual $6$-dimensional $G_3$-module $X=\Lambda^2(\CC^4)$ has
$(c_0,\ldots,c_6)=(1,0,1,1,3,10,21)$,
and $L(p^2)=-25$ for $p=t^3-6t^2+3t+5$.
\end{enumerate}
\end{theorem}

\begin{proof}
(1) Let $m^{(n)}_b$ be the number of summands $V(a,b)$, $a\geq0$, in $(M\otimes M^*)^{\otimes n}$. By Lemma~\ref{lem:levi}, $c_n(X)=\sum_bm^{(n)}_b[V^{\otimes n}:V_b]$. From $M\otimes M^*=V(0,0)\oplus V(1,1)$ we get $c_1=1$. The square is $2V(0,0)\oplus4V(1,1)\oplus V(2,2)\oplus V(3,0)\oplus V(0,3)$, so $m^{(2)}_0=3$, $m^{(2)}_2=1$ and $c_2=4$. Iterating the Pieri rule gives $m^{(3)}_1=19$, $m^{(3)}_3=8$ and $(m^{(4)}_0,m^{(4)}_2,m^{(4)}_4)=(67,66,24)$. With $V^{\otimes3}=2V_1\oplus V_3$ and $V^{\otimes4}=2V_0\oplus3V_2\oplus V_4$ this gives $c_3=46$ and $c_4=356$. The Weyl-denominator computation in \texttt{independent\_levi\_check.py} verifies all five values independently.

(2) By Lemma~\ref{lem:levi}, $c_n(X)$ is the number of summands $\mathrm{Sym}^k(M)$, $k\geq0$, in $\Lambda^2(M)^{\otimes n}$. By the Pieri rule this counts the ways of building a Young diagram of shape $(k+a,a,a,a)$ by adding $n$ times $2$ boxes in different rows. For instance $c_6=16+5$, with $16$ invariants of $SL_4$ and $5$ copies of $\mathrm{Sym}^4(M)$. Orthogonalizing $1,t,t^2,t^3$ for the resulting functional $L$ gives $p_2=t^2-t-1$ of norm $1$, and then $p_3=p$, with $L(p^2)=c_6-6^2-3^2-1^2=-25$. By Theorem~\ref{thm:main}, $P(G_3)$ fails.
\end{proof}

In case (2), the generating function
\[
 \sum_{k\geq0}\chi_{\mathrm{Sym}^k(M)}(g)r^k=\det(1-rg)^{-1}\qquad(0<r<1)
\]
has limit $(1-i)^{-4}=-1/4$ at the central element $g=iI$ of $SU(4)$, where the character of $\Lambda^2(M)$ reaches its minimum $-6$. This negative boundary value gives a heuristic explanation of the obstruction above. Turning it into a proof, for instance through a signed spectral measure, would require a boundary-value analysis at the matrices having eigenvalue $1$. For the doubled module $X\otimes X$ the obstruction persists, but it only appears at Hankel size 43 (Theorem~\ref{thm:exterior-doubled}). For $N=2$, Proposition~\ref{prop:tilted-haar} gives a positive law for every restricted self-dual $SL_3$-module, while the counterexample uses the Levi factor separately.

The failures of Theorem~\ref{thm:levi} persist for all $N$, and a single module works. Let $U$ be the natural module of the Levi factor $SL_N$, inflated to $G_N$, and consider in the same way the affine group $\CC^N\rtimes GL_N\subset GL_{N+1}$, with its natural modules $M$ and $U$.

\begin{theorem}\label{thm:vector-uniform}
For $N\geq2$, the self-dual module $X=M\oplus M^*\oplus U\oplus U^*$ has no law, both for $G_N=\CC^N\rtimes SL_N$ and for the affine group $\CC^N\rtimes GL_N$. \begin{enumerate}
\item For $\CC^N\rtimes GL_N$ with $N\geq4$, and for $G_N$ with $N\geq9$, the invariant dimensions $(c_0,\ldots,c_8)$ are $(1,1,7,22,153,806,5770,41378,313005)$, with Hankel determinants $1,6,399,336106,-4915544623$.
\item For the smaller $N$, the first negative Hankel determinant has size $5$, except for $G_2$ and $G_4$, where it has size $6$.
\end{enumerate}
Thus $P(G)$ fails for all the groups $\CC^N\rtimes SL_N$ with $N\geq2$, and $\CC^N\rtimes GL_N$ with $N\geq1$.
\end{theorem}

\begin{proof}
For the affine group $G'=\CC^N\rtimes GL_N$, Lemma~\ref{lem:levi} applies with $H=GL_{N+1}$, $\lambda=\varepsilon_1$ and $L_\lambda=GL_N$. Expand $X^{\otimes n}$ into modules $Y|_{G'}\otimes W$, with $Y=M^{\otimes a}\otimes M^{*\otimes b}$, $W=U^{\otimes c}\otimes U^{*\otimes d}$ and $a+b+c+d=n$. For $N\geq n$ the summands of $Y$ are the rational modules $V_{[\alpha,\beta]}$ of $GL_{N+1}$, of highest weight $(\alpha_1,\alpha_2,\ldots,0,\ldots,0,\ldots,-\beta_2,-\beta_1)$, with $\alpha\vdash a-k$ and $\beta\vdash b-k$, and their multiplicities do not depend on $N$ \cite{Koike}. The Levi module $V^L_{[\alpha,\beta]}$ restricted to $GL_N$ is the rational module $[\hat\alpha,\beta]$, where $\hat\alpha$ is $\alpha$ without its first row, and its multiplicity in $W^*$ does not depend on $N$ either. So $c_n(X)$ does not depend on $N\geq n$.

For $G_N\subset G'$, the quotient $G'/G_N\simeq GL_1$ acts on the invariants of $G_N$ through the powers of the determinant, and
\[
   (X^{\otimes n})^{G_N}=\bigoplus_{k\in\ZZ}\,(X^{\otimes n}\otimes\det{}^{-k})^{G'} .
\]
The central scalar matrices of $GL_N$ act on $X^{\otimes n}$ with weights in $[-n,n]$, and on $\det^k$ with weight $kN$, so only $k=0$ contributes when $N>n$. Hence the invariant dimensions for $G_N$ and $G'$ agree when $N>n$.

For clarity, the finite calculation can be reproduced by the following Pieri recurrence. Keep a multiplicity for each pair of dominant integral weights $(\alpha,\beta)$ of lengths $N+1,N$, respectively, starting with one copy of $(0,0)$. At each step add or subtract $1$ at a single coordinate of either weight, retaining each dominant result with its multiplicity. These are precisely the four summands of $X$. For $G'$ sum the resulting multiplicities over the pairs satisfying
\[
       \alpha_{i+1}+\beta_{N+1-i}=0\qquad(1\le i\le N).
\]
Indeed $\alpha_2,\ldots,\alpha_{N+1}$ is the restricted Levi highest weight, and $-\beta_N,\ldots,-\beta_1$ is the dual weight. For $G_N$ one sums over pairs for which these $N$ sums are equal. Applying this recurrence through degree ten for $2\le N\le9$ gives the determinants in the statement, including the degree-eight stable tuple for $4\le N<8$ in the affine case. The exact weight calculation and its output are supplied in \texttt{vector\_stabilizers.py} and \texttt{vector-stabilizers-output.txt}. The case $N=1$ of the affine groups is Example~\ref{ex:affine}.
\end{proof}

\subsection{The complex Euclidean groups}\label{sec:euclidean}
Let $E=\CC^{m+2}$ with a nondegenerate quadratic form, and choose isotropic vectors $e,f\in E$ with $\langle e,f\rangle=1$. Then $E=\CC e\oplus W\oplus\CC f$, with $W=\{e,f\}^\perp\simeq\CC^m$, and the stabilizer of $e$ in $SO(E)$ is the complex Euclidean group
\[
   E_m=W\rtimes SO(W)=\CC^m\rtimes SO_m .
\]
Translation by $w\in W$ fixes $e$, takes $x\in W$ to $x-\langle x,w\rangle e$, and takes $f$ to $f+w-\tfrac12\langle w,w\rangle e$. For $m\ge2$, the full inverse image of $E_m$ in $Spin_{m+2}$ is the connected double cover
\[
        \widetilde E_m=\CC^m\rtimes Spin_m\longrightarrow E_m .
\]
In the parabolic stabilizing $\CC e$, the Levi is $(\Gm\times Spin_m)/\langle(-1,-1)\rangle$. The two factors arise from the orthogonal summands $\CC e+\CC f$ and $W$, and their lifts meet in the central subgroup $\{\pm1\}$. The first factor acts on $e$ by its square (choosing its orientation accordingly). The Levi of the vector stabilizer is therefore $Spin_m$. For $m=3$ the two Levis are $GL_2$ and $SL_2$, respectively. For $m=1$ the ordinary stabilizer is $E_1=\Ga$; its inverse image has identity component $\Ga$ and kernel $Spin_1=\mu_2$.

For $m\ge2$ take $H=Spin_{m+2}$ and $\lambda=\varepsilon_1$ in Lemma~\ref{lem:levi}. A highest weight $\mu=(\mu_1,\mu_2,\ldots)$ restricts on the $Spin_m$ Levi to $\hat\mu=(\mu_2,\mu_3,\ldots)$. The stabilizer of the anisotropic vector $e+f$ in $H$ is $K_m=Spin_{m+1}$.

\begin{proposition}\label{prop:spherical}
For $m\geq2$ and every $Spin_{m+2}$-module $Y$,
\[
   \dim Y^{\widetilde E_m}=\dim Y^{K_m}=\sum_{k\geq0}\,[Y:V_{k\varepsilon_1}] .
\]
The law of $Y|_{\widetilde E_m}$ is the distribution of $\chi_Y$ on $Spin(m+1)$ under Haar probability measure. This also gives the law on $E_m$ whenever the covering kernel acts trivially on $Y$.
\end{proposition}

\begin{proof}
By Lemma~\ref{lem:levi}, the module $V_\mu$ has a $\widetilde E_m$-invariant line exactly when $\mu$ is a rational multiple of $\varepsilon_1$. For $m\geq2$ this means $\mu=k\varepsilon_1$ with $k\in\NN$, since the half-integral weights of $Spin_{m+2}$ have no coordinate $0$. On the other hand, by the interlacing branching rule from $Spin_{m+2}$ to $Spin_{m+1}$ \cite{GW}, the restriction of $V_\mu$ contains the trivial module once if $\mu=k\varepsilon_1$, and otherwise has no invariants. Applying this to $Y^{\otimes n}$ gives the required moments by Haar integration. The restricted character is real. One of the consecutive spin groups has odd vector dimension and hence all its irreducible modules are self-dual. This gives self-duality either of $Y$ or of its restriction to $K_m$.
\end{proof}

The comparison is the classical horospherical degeneration of the complex sphere to its isotropic cone \cite[\S1.2, Theorem~5 and Example~2]{Timashev}. In the harmonic polynomials $V_{k\varepsilon_1}$, the anisotropic stabilizer fixes a zonal harmonic, and the isotropic stabilizer fixes $x\mapsto\langle x,e\rangle^k$. This is also the group contraction of In\"on\"u and Wigner \cite{IW}. For the full module category of $E_m$, the existence of these laws leaves room for counterexamples. Write $E$ for the restricted vector module of $SO_{m+2}$ and $W$ for the vector module inflated from $SO_m$.

\begin{theorem}\label{thm:euclidean}
The property $P(E_m)$ fails for every $m\geq2$. More precisely:
\begin{enumerate}
\item For $m\geq5$ the self-dual module $X=E\oplus W$ has $(c_0,\ldots,c_4)=(1,1,3,10,33)$, and
\[
   L\bigl((2t^2-7t+1)^2\bigr)=-2 .
\]
\item For $m=4$ the module $X=E\oplus W$ has $(c_0,\ldots,c_6)=(1,1,3,10,34,197,825)$, and its Hankel determinant of size $4$ is $-13341$. For $m=2$ it has $(c_0,\ldots,c_4)=(1,1,4,17,60)$, and its Hankel determinant of size $3$ is $-37$.
\item For $m=3$ embed $E_3$ in $G_3=\CC^3\rtimes SL_3\subset SL_4$, as $\CC^3\rtimes SO_3$. Then the self-dual 14-dimensional module $X=(M\oplus M^*\oplus\Lambda^2M)|_{E_3}$ has
\[
   (c_0,\ldots,c_8)=(1,1,8,49,380,3233,29520,284623,2859780),
\]
and $L(p^2)=-5682$ for $p=t^4-20t^3+112t^2-107t-178$.
\end{enumerate}
\end{theorem}

\begin{proof}
(1) Write $d(a,b)=\dim(E^{\otimes a}\otimes W^{\otimes b})^{E_m}$, so that $c_n(X)=\sum_a\binom na d(a,n-a)$. By Lemma~\ref{lem:levi},
\[
   d(a,b)=\sum_\lambda m_a(\lambda)\,m'_b(\hat\lambda),
\]
where $m_a(\lambda)$ is the multiplicity in $E^{\otimes a}$ of the irreducible $SO_{m+2}$-module $[\lambda]$ attached to a partition $\lambda$, the partition $\hat\lambda=(\lambda_2,\lambda_3,\ldots)$ is $\lambda$ without its first row, and $m'_b(\nu)$ is the multiplicity of $[\nu]$ in $W^{\otimes b}$, for $SO_m$. Since $W^{\otimes b}$ is self-dual, dual irreducibles occur there with the same multiplicity. For $a\leq4$ and $m+2\geq7$, Brauer's theory gives $m_a(\lambda)$ as the number of walks of length $a$ from $\emptyset$ to $\lambda$ in Young's lattice, each step adding or removing one box \cite{Brauer,GW}:
\begin{center}\small
\begin{tabular}{@{}c|l@{}}
$a$ & $m_a(\lambda)$\\ \hline
$1$ & $(1)$:1\\
$2$ & $\emptyset$:1, $(2)$:1, $(1,1)$:1\\
$3$ & $(1)$:3, $(3)$:1, $(2,1)$:2, $(1,1,1)$:1\\
$4$ & $\emptyset$:3, $(2)$:6, $(1,1)$:6, $(4)$:1, $(3,1)$:3, $(2,2)$:2, $(2,1,1)$:3, $(1^4)$:1
\end{tabular}
\end{center}
Only the partitions with $\hat\lambda\in\{\emptyset,(1)\}$ enter below, that is $\lambda=(k)$ and $\lambda=(k,1)$, and for $m+2\geq7$ these stay irreducible on $SO_{m+2}$, and pairwise nonisomorphic. The coincidences in this range, namely $[1^4]\simeq[1^3]$ for $m=5$ and the splitting of $[1^4]$ for $m=6$, only concern partitions with $\hat\lambda=(1,1)$ or $(1,1,1)$. On the side of $W$, only the following values are needed, and for $m\geq5$ they are those of the orthogonal group: $m'_b(\emptyset)=1,0,1,0,3$ for $b=0,\ldots,4$ (pairings), $m'_1((1))=1$, and $m'_2((1))=0$. We obtain
\begin{center}\small
\begin{tabular}{@{}c|ccccc|c@{}}
$n$ & $d(n,0)$ & $d(n-1,1)$ & $d(n-2,2)$ & $d(n-3,3)$ & $d(n-4,4)$ & $c_n$\\ \hline
$1$ & $1$ & $0$ & & & & $1$\\
$2$ & $2$ & $0$ & $1$ & & & $3$\\
$3$ & $4$ & $1$ & $1$ & $0$ & & $10$\\
$4$ & $10$ & $2$ & $2$ & $0$ & $3$ & $33$
\end{tabular}
\end{center}
For instance $d(4,0)$ counts the $\lambda$ with $\hat\lambda=\emptyset$, that is $3+6+1$, and $d(3,1)=m_3((2,1))\,m'_1((1))=2$. Then $c_4=10+4\cdot2+6\cdot2+4\cdot0+3=33$, and similarly in the lower degrees. The Hankel matrix $(c_{i+j})_{0\leq i,j\leq2}$ has determinant $-1$, and $2t^2-7t+1$ is twice the corresponding monic orthogonal polynomial, which gives $-2$. Equivalently, in the setting of Theorem~\ref{thm:main}, the modules $A=\one\oplus(X^{\otimes2})^{\oplus2}$ and $B=X^{\oplus7}$ have $h(A,A)=145$, $h(A,B)=147$ and $h(B,B)=147$.

(2) For $m\leq4$ the stable values above change. For $m=4$ the determinant gives $m'_4(\emptyset)=4$, and for $m=3$ we have $m'_2((1))=1$, since $\Lambda^2\CC^3\simeq\CC^3$, and $m'_3(\emptyset)=1$. The exact characteristic-zero weight calculations in \texttt{euclidean\_groups.py} give the stated sequences and determinants. For $2\le m\le7$, the additional Lie-algebra kernel calculations in \texttt{independent\_euclidean\_groups\_check.py} give matching upper bounds through the degrees used here. Indeed reducing an integral matrix modulo $2^{31}-1$ can only decrease its rank, hence can only increase its nullity. For $m=2$ the failure also follows from Theorem~\ref{thm:solvable}.

(3) Identify $M=\CC e_1\oplus\CC^3$, with $SO_3=SO(\CC^3)$ acting on the second summand. By Lemma~\ref{lem:levi} for $G_3$, the space $V_\mu^{U_3}$ is the irreducible $SL_3$-module of highest weight $\bar\mu$, and $E_3=U_3\rtimes SO_3$. By the Cartan--Helgason theorem for the symmetric pair $(SL_3,SO_3)$ \cite{GW}, an irreducible $SL_3$-module has an $SO_3$-invariant line exactly when its highest weight is twice a dominant weight, and has no invariants otherwise. Thus $c_n(X)$ is the number of summands $V_\mu$ of $(M\oplus M^*\oplus\Lambda^2M)^{\otimes n}$ such that $\bar\mu$ is even, and \texttt{euclidean\_groups.py} computes these numbers with the Pieri rule. The Hankel determinants of sizes $1,\ldots,5$ are $1$, $7$, $531$, $390368$ and $-2855167756$. The independent modular Lie-algebra computation in \texttt{independent\_euclidean\_check.py} gives matching upper bounds through degree eight by the same rank comparison.
\end{proof}

The modules in (3) are natural. Identify $M$ with the hyperplane $e^\perp\subset\CC^5$, on which the quadratic form has radical $\CC e$. Then $M^*\simeq\CC^5/\CC e$ is the affine module $\CC^3\oplus\CC$, on which $(w,g)$ acts by the matrix $\left(\begin{smallmatrix}g&w\\0&1\end{smallmatrix}\right)$, and $u\wedge v\mapsto\langle\cdot,v\rangle u-\langle\cdot,u\rangle v$ identifies $\Lambda^2M$ with the adjoint module $\Lie(E_3)\subset\mathfrak{gl}(M)$. Thus $X$ is the sum of the affine module of the Euclidean group of 3-space, its dual, and its adjoint module. In contrast, for $m=3$ the module $E\oplus W$ of (1) and (2) has a law, with $(c_0,\ldots,c_6)=(1,1,3,11,48,232,1207)$, by the following result.

\subsection{The spin cover in dimension three}
For $\widetilde E_3=\CC^3\rtimes SL_2\subset Sp_4=Spin_5$ and $K=K_3=Spin_4=SL_2\times SL_2$, the compact-group formula extends to modules that also involve the Levi factor. Each representation of $\widetilde E_3$ on which the central covering involution acts trivially descends to $E_3$ with the same invariant dimensions.

\begin{theorem}\label{thm:contraction}
For $Y\in\Rep(Sp_4)$ and $W\in\Rep(SL_2)$, inflated to $\widetilde E_3$,
\[
   \dim\,(Y|_{\widetilde E_3}\otimes W)^{\widetilde E_3}=\dim\,\bigl(Y|_K\otimes(W\boxtimes\CC)\bigr)^K .
\]
Consequently every $\widetilde E_3$-module of the form $X=\bigoplus_iY_i|_{\widetilde E_3}\otimes W_i$ has a law, namely the distribution of the function
\[
   \chi(g,h)=\sum_i\chi_{Y_i}(g,h)\,\chi_{W_i}(g)
\]
on $SU(2)\times SU(2)\subset Sp(2)$, with respect to the Haar measure.
\end{theorem}

\begin{proof}
It is enough to take $Y=V_\mu$ and $W=V_k$ irreducible, with $\mu=(\mu_1,\mu_2)$ in the coordinates of $Spin_5$, where $\mu_1\geq\mu_2\geq0$ are both integers or both half-integers. By Lemma~\ref{lem:levi}, the left hand side is $[V_k:V_{2\mu_2}]=\delta_{k,2\mu_2}$, because the Levi factor $Spin_3$ sees the coordinate $\mu_2$. On the other hand, the restriction of $V_\mu$ to $Spin_4$ is the sum of the modules of highest weight $(c,d)$, with $c,d\equiv\mu_1$ modulo $\ZZ$ and
\[
   \mu_1\geq c\geq\mu_2\geq|d| ,
\]
each once \cite{GW}, and $(c,d)$ corresponds to $V_{c+d}\boxtimes V_{c-d}$. This contains $V_k\boxtimes\CC$ only when $c=d=k/2$, and then $c\geq\mu_2\geq|d|$ forces $\mu_2=k/2$. Thus the right hand side is $\delta_{k,2\mu_2}$ too. For the second assertion, $X^{\otimes n}$ is again of this form, the transported modules multiply in the same way, and all the modules involved are self-dual.
\end{proof}

The formula transports $W$ to the first $SU(2)$ factor. For the four-dimensional spinor $S$ of $Spin_5$, the covering involution acts by $-1$, so $S$ is a genuine module of $\widetilde E_3$. The five-dimensional vector module $E$ descends to $E_3$. With $s_1=\chi_{V_1}(g)$ and $s_2=\chi_{V_1}(h)$, which are independent and semicircular on $[-2,2]$, their laws are those of
\[
   \chi_S=s_1+s_2,\qquad \chi_E=1+s_1s_2 .
\]
For $S$ this gives the moments
\[
   c_{2n}=\sum_k\binom{2n}{2k}C_kC_{n-k}=C_nC_{n+1} ,
\]
the number of walks of length $2n$ in the quarter plane with steps N, S, E, W returning to the origin \cite{GKS}. Equivalently, $\chi_S^2$ has the law of $xy$, with $x$ free Poisson of parameter $1$ and $y$ independent, with the size-biased law $y\,d\pi(y)$. For $E$ the moments are $1,1,2,4,11,31,101,337,1205,4405,\ldots$ Finally the module $E\oplus W$ has the law of
\[
   1+s_1s_2+(s_1^2-1)=s_1(s_1+s_2) ,
\]
which is also the law of $S\otimes V_1$. Thus $E\oplus W$ and the module of Theorem~\ref{thm:euclidean}(3) give examples on the same ordinary group $E_3$ with opposite answers to the Hamburger moment problem.

Subsection~\ref{sec:further} and the computational supplement record other negative squares for restrictions to Euclidean groups and to vector stabilizers, together with the degrees tested (\texttt{euclidean\_search.py} and \texttt{levi.py}). For the family in Theorem~\ref{thm:contraction}, the compact-group formula proves existence in every degree.

\subsection{Further families and computations}\label{sec:further}
The programs \texttt{linear.py} and \texttt{orthogonal.py} define further stabilizers. For each, \texttt{families.py} tests self-dual modules $Y|_G\otimes W$ built from small irreducibles of $H$ and $L_\lambda$, and sums of 2 of them, up to dimension 400. The target is degree 12 (Hankel size 7), with a second target of degree 40 (size 21) for families with no detected obstruction. Time and size limits truncate individual records: for example, only 11 of the 151 modules for $\CC^4\rtimes SL_4$ reach degree 12, and only 29 of the 63 for $\CC^3\rtimes Spin_3$ reach degree 40. The supplement records each attained degree and verifies each negative square (\texttt{summarize\_searches.py}).

The \texttt{euclidean.py} transcript records 1491 of 1719 announced candidates for $\widetilde E_3\subset Sp_4$, all through degree 12, with no second pass or final summary. The \texttt{euclidean\_search.py} transcript records 465 modules $Y|_{E_3}\otimes W$, with $Y\in\Rep(SL_4)$ and $W\in\Rep(SO_3)$, through degrees 10--12. It computes restriction from $SL_3$ to $SO_3$ by Gelfand--Tsetlin patterns. In the table, \emph{Failures} counts modules with an exact negative Hankel direction in the archived run, and \emph{Size} is the smallest failing size found in each listed family. A record without failure need not have reached the target or have a law.
\begin{center}\small
\begin{tabular}{@{}llrrc@{}}
\toprule
Group $G$ & $H$ & Modules & Failures & Size\\
\midrule
$\Ga$ & $Spin_3=SL_2$ & 3 & 0 & --\\
$\CC^3\rtimes Spin_3$ & $Spin_5$ & 63 & 0 & --\\
$\CC^m\rtimes Spin_m$, $m=2,4,5,6,7,8$ & $Spin_{m+2}$ & 925 & 61 & 3, 4\\
$\Lambda^2(\CC^n)\rtimes SL_n$, $n=5,6$ & $Spin_{2n}$ & 59 & 15 & 4\\
$\CC^N\rtimes SL_N$, $2\leq N\leq7$ & $SL_{N+1}$ & 576 & 199 & 3, 4\\
$\CC^N\rtimes GL_N$, $1\leq N\leq4$ & $GL_{N+1}$ & 1632 & 495 & 3\\
$\mathrm{Sym}^2(\CC^2)\rtimes SL_2$ & $Sp_4$ & 63 & 0 & --\\
$\mathrm{Sym}^2(\CC^n)\rtimes SL_n$, $n=3,4$ & $Sp_{2n}$ & 206 & 48 & 3\\
$\mathrm{Sym}^2(\CC^2)\rtimes GL_2$ & $Sp_4$ & 189 & 85 & 3\\
$(\CC^k\otimes\CC^l)\rtimes(SL_k\times SL_l)$, $kl=6,8,9$ & $SL_{k+l}$ & 939 & 214 & 3\\
$\widetilde E_3=\CC^3\rtimes SL_2$, by \texttt{euclidean.py} & $Sp_4$ & 1491 & 0 & --\\
$E_3=\CC^3\rtimes SO_3$, by \texttt{euclidean\_search.py} & $SL_4$ & 465 & 33 & 4\\
\bottomrule
\end{tabular}
\end{center}
The first row is the connected stabilizer: the full spin preimage in dimension one is $\Ga\times\mu_2$, with $\mu_2$ acting by $-1$ on the spinor (Section~\ref{sec:euclidean}). The Grassmannian rows have $(k,l)=(2,3),(2,4),(3,3)$; $\mathrm{Sym}^2(\CC^2)\rtimes GL_2$ is the Siegel parabolic of $Sp_4$. Proposition~\ref{prop:additive-law} and Theorem~\ref{thm:contraction} prove laws in every degree for the tested additive and transported $\widetilde E_3$ families. The latter occurs under both names $\CC^3\rtimes Spin_3$ and $\mathrm{Sym}^2(\CC^2)\rtimes SL_2$; other modules of the same cover need not have laws.

Every other listed family has a detected failure at size 3, except $m=4$, the half-spin groups, $\CC^7\rtimes SL_7$ and $E_3\subset SL_4$, whose smallest sizes found are 4. For $\CC^m\rtimes Spin_m$ with $m\neq1,3$, the first witness found is $E\oplus W$, as in Theorem~\ref{thm:euclidean}. Restricted witnesses at size 4 include the half-spin module $S^+$ of $Spin_{12}$ on $\Lambda^2(\CC^6)\rtimes SL_6$, the sum $S^+\oplus S^-$ for $Spin_{10}$, and $\Lambda^2M\oplus\Lambda^2M^*$ for $\CC^N\rtimes SL_N$, $N=6,7$. On $E_3\subset SL_4$, $\Lambda^2M\otimes(\CC\oplus W)$, of dimension 24, fails at size 4; the module of Theorem~\ref{thm:euclidean}(3) fails at size 5. The smallest dimension found is 9, for $\Lambda^2M\oplus W$, failing at size 7.

These witnesses vary with the rank; Theorem~\ref{thm:vector-uniform} instead gives a single construction for every $N$, with a negative determinant of size 5 or 6. Finally, \texttt{contraction.py} compares the invariant dimensions and Haar integrals of Theorem~\ref{thm:contraction} for 51 modules by exact integer character computations.


\section{Projective covers and Cartan matrices}\label{sec:cartan}

Let $\C$ be a finite tensor category over $\CC$, with simples $S_0=\one,S_1,\ldots,S_{r-1}$ and projective covers $P_0=P(\one),P_1,\ldots,P_{r-1}$. The cover $P_i$ is the indecomposable projective with simple top $S_i$. Exactness of $\Hom(P_i,-)$ gives $h(P_i,Y)=[Y:S_i]$. In particular, the Cartan matrix $C_{ij}=[P_j:S_i]$ records the dimensions $h(P_i,P_j)$. Write $e_0$ for the unit coordinate. Unimodularity means that the simple socle of $P_0$ is $\one$ \cite[\S\S6.1, 6.5]{EGNO}.

For a projective object the ordinary Grothendieck ring gives a finite test for Hamburger positivity. For the unit cover in fusion type, the residues of its rational moment transform first force unimodularity. The shifted sequence then has a positive spectral measure whose mass at zero gives the obstruction in Theorem C.

\begin{theorem}[Finite test for a projective object]\label{thm:finite-test}
Let $Q$ be projective, let $m_i$ be the multiplicity of $P_i$ in $Q^*$,
and let $N$ denote left multiplication by $[Q]$ on $\Gr(\C)$ in the
simple basis, with coordinate vectors as columns. Then $c_0(Q)=1$ and
\[
 c_n(Q)=m^{\mathsf T} N^{n-1}e_0\quad(n\ge1),\qquad
 G_Q(z):=\sum_{n\ge0}\frac{c_n(Q)}{z^{n+1}}
   =\frac{1+m^{\mathsf T}(zI-N)^{-1}e_0}{z}.
\]
The following are equivalent.
\begin{enumerate}
\item The invariant dimensions of $Q$ are Hamburger moments.
\item $H_r(Q)=(c_{i+j}(Q))_{0\le i,j\le r}$ is positive semidefinite.
\item All poles of $G_Q$ in lowest terms are real and simple with positive residues.
\end{enumerate}
The residues give the unique representing measure. If positivity fails, some $f\in\mathbb Z[t]$ of degree at most $r$ satisfies $L_Q(f^2)<0$, where $L_Q(t^n)=c_n(Q)$. If $Q\simeq Q^*$,
these conditions are also equivalent to \eqref{eq:cs} for all $A,B$
that are sums of powers of $Q$ of degrees at most $r$. In particular,
failure is witnessed by 2 such objects. In this case $[Q]=Cm$, so
the input for the test consists of the ordinary multiplication table,
the Cartan matrix and $m$.
\end{theorem}
\begin{proof}
Duals of projectives are projective \cite[Proposition~6.1.3]{EGNO}.
The exact functor $\Hom(P_i,-)$ has dimension $[Y:S_i]$ on $Y$, since
its dimension on $S_j$ is $\delta_{ij}$. For $n\ge1$, rigidity identifies
$\Hom(\one,Q^{\otimes n})$ with $\Hom(Q^*,Q^{\otimes(n-1)})$.
This proves the moment formula and summing the series gives $G_Q$.

Let $L=L_Q$ and $q=t\det(tI-N)$. Cayley--Hamilton gives $L(t^j q)=0$ for $j\ge0$. Reducing $f$ modulo $q$ preserves $L(f^2)$ and leaves a remainder of degree at most $r$ (compare \cite[\S3]{CF}). Thus $H_r\succeq0$ tests all squares. Hamburger's theorem gives (1), and $L(q^2)=0$ puts every representing measure on the real zeros of $q$. Partial fractions of $G_Q$, whose leading term is $1/z$, now give (3), its converse and uniqueness. If $H_r$ has a negative direction, diagonalizing its rational symmetric form by congruence gives such a direction with rational entries. Clearing denominators yields the stated $f$. For $Q\simeq Q^*$, \eqref{eq:bridge} identifies $H_r$ with the Hom matrix on $\one,Q,\ldots,Q^{\otimes r}$. Splitting $f$ into positive and negative coefficient parts gives the two sums violating \eqref{eq:cs}, as in Theorem~\ref{thm:main}.
\end{proof}
\subsection{Categories of fusion type}
For a fusion ring $R$ write $\tau$ for its unit coefficient. Its simple basis is orthonormal for $\langle x,y\rangle=\tau(xy^*)$, and $N_{x^*}=N_x^*$ in the notation of Section~\ref{sec:fusion}.

\begin{definition}\label{def:fusion-type}
A finite tensor category $\C$ is \emph{of fusion type} if its Grothendieck ring $\Gr(\C)$ is a fusion ring in the basis of simple classes. Equivalently, $[S\otimes T:\one]\leq1$ for all simples $S,T$, with equality only if $T\simeq S^*$. A \emph{categorification} of a fusion ring $R$ is a finite tensor category $\C$ together with an isomorphism $\Gr(\C)\simeq R$ which maps the simple classes onto the basis.
\end{definition}

For the equivalence note that $[S\otimes S^*:\one]\geq1$, by coevaluation. Thus the fusion ring axiom for the unit coefficient holds precisely under the stated condition, the involution being induced by duality. The same condition applied to the evaluation $S\otimes{}^*S\to\one$ identifies the right dual with $S^*$, so duality is an involution on simple classes. Exactness and rigidity give the remaining axioms. Fusion categories are of fusion type, and so are the categories $\Rep(T_m)$ and $\Rep(H_r)$ of Section~\ref{sec:hopf}, whose Grothendieck rings are group rings. In contrast, the small quantum group $u_q(\mathfrak{sl}_2)$ at a root of unity of odd order $p\geq3$ is not of fusion type: for its Steinberg module $S$, of dimension $p$, the product $S\otimes S^*$ contains the projective cover of $\one$, in which $\one$ occurs twice.

\begin{proposition}\label{prop:fusion-type}
Let $\C$ be a finite tensor category of fusion type, let $P$ be the projective cover of $\one$, and put $a=[P]\in\Gr(\C)$. Let $S,T$ be simple.
\begin{enumerate}
\item The left and right duals of $S$ are isomorphic, so $S^{**}\simeq S$, and $\one$ is a direct summand of $S\otimes S^*$.
\item $P\otimes S$ and $S\otimes P$ are both projective covers of $S$. Hence $a$ is central, and the Cartan matrix is the matrix $N_a$ of multiplication by $a$.
\item If $\C$ is unimodular, then $P\simeq P^*$, so $a=a^*$ and the Cartan matrix is symmetric.
\item For any isomorphism $\phi:S\to S^{**}$ the scalar $\mathrm{ev}_{S^*}(\phi\otimes\mathrm{id}_{S^*})\mathrm{coev}_S$ is nonzero.
\item If the invariant dimensions of $P$ are Hamburger moments, then $\C$ is unimodular.
\end{enumerate}
\end{proposition}

\begin{proof}
(1) Both $S^*$ and the right dual ${}^*S$ are simple, and $\one$ embeds in $S\otimes S^*$ and in ${}^*S\otimes S$. By the unit coefficient axiom both have class $[S]^*$, so they are isomorphic. Now $\one$ is a subobject of $S\otimes S^*$ by coevaluation, and a quotient by $\Hom(S\otimes S^*,\one)\simeq\Hom(S,S^{**})$. Since its multiplicity is $1$, the composite $\one\to S\otimes S^*\to\one$ cannot vanish, since a zero composite would put one copy in the kernel and another in the quotient.

(2) The object $P\otimes S$ is projective, and $h(P\otimes S,T)=h(P,T\otimes S^*)=[T\otimes S^*:\one]=\delta_{S,T}$ by Theorem~\ref{thm:frobenius}. Hence $P\otimes S$ is the projective cover of $S$. The same holds for $S\otimes P$, using right duals and (1). Thus $a$ commutes with the basis, and $C_{ij}=[P\otimes S_j:S_i]$ is the coefficient of $[S_i]$ in $a[S_j]$.

(3) The dual of an indecomposable projective is an indecomposable projective \cite[Proposition~6.1.3]{EGNO}, and $P^*$ has $\one$ as a quotient, since $P$ has $\one$ as a subobject.

(4) The map $\mathrm{ev}_{S^*}(\phi\otimes\mathrm{id})$ is a nonzero element of the line $\Hom(S\otimes S^*,\one)$, so by the proof of (1) its composite with the coevaluation is nonzero.

(5) Let $D$ be the simple top of $P^*$ and $d=[D]$. By (2), $P^*\simeq P\otimes D$, whence $a^*=ad$. The central element $a$ is normal in the finite-dimensional $*$-algebra $R_{\CC}$, where $R=\Gr(\C)$. Write $e_\lambda$ for its spectral idempotents. Each is a self-adjoint central projection, and $w_\lambda=\tau(e_\lambda)=\tau(e_\lambda e_\lambda^*)>0$ by faithfulness of $\tau$. For $\lambda\ne0$, the identity $a^*=ad$ gives
\[
 d e_\lambda=\frac{\bar\lambda}{\lambda}e_\lambda,\qquad
 \tau(d^*e_\lambda)=\frac{\lambda}{\bar\lambda}w_\lambda .
\]
By Theorem~\ref{thm:finite-test}, $c_{n+1}(P)=\tau(d^*a^n)$ for $n\ge0$. Thus
\[
 G_P(z)=\frac{1}{z}\left(1+\sum_\lambda
                    \frac{\tau(d^*e_\lambda)}{z-\lambda}\right).
\]
Every nonzero spectral value $\lambda$ is a pole, with residue $w_\lambda/\bar\lambda$. A Hamburger measure forces these poles to be real with positive residues, so every nonzero $\lambda$ is positive. There is at least one such pole, since $\FPdim(P)>0$ is an eigenvalue of $N_a$. The only other possible pole is zero. The representing measure therefore has strictly positive first moment $h(\one,P)$. The socle of $P$ is simple, so it is isomorphic to $\one$, giving unimodularity.
\end{proof}

The equality $a=a^*$ alone need not imply unimodularity. In the Sweedler category $a=1+s$ in $\ZZ[C_2]$, although $P$ and $P^*$ have different simple tops.

\subsection{The Hom form on powers of the unit cover}
Assume in this subsection that $\C$ is unimodular and of fusion type, with $\Gr(\C)=R$. Put $P=P_0$ and $a=[P]$ as above, and denote by
$\nu_a$ the spectral probability measure of left multiplication by $a$
at $e_0$ in the regular representation of $R$. Write $L_P(t^n)=c_n(P)$.

\begin{theorem}[Cartan criterion]\label{thm:cartan}
The Cartan matrix is the symmetric matrix $C=N_a$, and $a$ is central.
For $n\ge1$ we have $c_n(P)=\tau(a^{n-1})=\int t^{n-1}\,d\nu_a(t)$.
The following are equivalent.
\begin{enumerate}
\item The Hom dimensions satisfy \eqref{eq:cs} on finite sums of powers of $P$.
\item $P$ has a Hamburger measure.
\item Every self-dual projective has a Hamburger measure.
\item $C-e_0e_0^{\mathsf T}$ is positive semidefinite.
\item $C$ is positive definite and $(C^{-1})_{00}\le1$.
\end{enumerate}
Write $w_\lambda=\nu_a(\{\lambda\})$, with $w_0=0$ if $0$ is absent from the spectrum. On all real polynomials one has
\begin{equation}\label{eq:cartan-hermite}
 L_P(f)=\left(1-\sum_{\lambda\ne0}\frac{w_\lambda}{\lambda}\right)f(0)
       +\sum_{\lambda\ne0}\frac{w_\lambda}{\lambda}f(\lambda)+w_0 f'(0).
\end{equation}
Under the equivalent positivity conditions the unique measure of $P$ is $\mu_P=t^{-1}\nu_a+(1-\int t^{-1}d\nu_a)\delta_0$, supported on $[0,\infty)$.
\end{theorem}
\begin{proof}
By Proposition~\ref{prop:fusion-type}, $a$ is central, $C=N_a$, and $C$ is symmetric. The transfer formula of Theorem~\ref{thm:finite-test}, with $Q=P$ and $m=e_0$, gives the shifted moment identity. Applying it to the polynomial $(f(t)-f(0))/t$, whose value at zero is $f'(0)$, gives \eqref{eq:cartan-hermite}.

The equivalence (1)$\Leftrightarrow$(2) is Theorem~\ref{thm:main}.
To compute its quadratic form, write $f=u+tg$ with $u\in\RR$ and
$g\in\RR[t]$, and set $v=g(C)e_0$. The moment identity gives
\begin{equation}\label{eq:cartan-square}
 L_P(f^2)=u^2+2u\langle v,e_0\rangle+\langle Cv,v\rangle
 =(u+\langle v,e_0\rangle)^2+
        \langle(C-e_0e_0^{\mathsf T})v,v\rangle .
\end{equation}
Thus (2) tests (4) on $K=\RR[C]e_0$. Each spectral projection of $C$ is multiplication by a polynomial in $a$, hence cannot vanish on $e_0=1$ unless it is zero. Every eigenvalue is therefore visible on $K$. Positivity on $K$ makes $C\succeq0$ everywhere, and on $K^\perp$ the rank-one term vanishes. This proves (2)$\Leftrightarrow$(4). Under (4), $\ker C\perp e_0$, so visibility gives $C>0$. Conjugating by $C^{-1/2}$ now gives exactly (5).

Under (5), $\nu_a$ has strictly positive support and
$\int t^{-1}\,d\nu_a=(C^{-1})_{00}$. The stated formula for $\mu_P$
has mass one and the correct higher moments. For any self-dual projective
$Q$, write $Q=P\otimes V$ for a semisimple $V$ using the cover calculation
above. Uniqueness of its simple multiplicities and $P\simeq P^*$ give
$v=[V]=v^*$. For $n\ge1$ the transfer formula yields
$c_n(Q)=\tau(a^{n-1}v^n)$. Since $a$ is central, the scalar spectral
measure of the self-adjoint $N_{av}$ at $C^{-1/2}e_0$ has these
positive-degree moments. Adding $(1-(C^{-1})_{00})\delta_0$ gives
a probability measure. This proves (3), which implies (2) by taking $Q=P$.
\end{proof}

\begin{corollary}[Semisimplicity and the unit cover]\label{cor:cartan-failure}
Let $\C$ be a finite complex tensor category of fusion type, and put $P=P(\one)$. If the identity and double-dual functors are naturally isomorphic as additive functors, then $P$ has a Hamburger measure if and only if $\C$ is semisimple. Under the additional hypothesis of unimodularity, these conditions are also equivalent to tensor-power Cauchy--Schwarz for $\C$.

In any unimodular category of fusion type with singular Cartan matrix, two sums of powers of $P$ of degrees at most the number of distinct Cartan eigenvalues violate \eqref{eq:cs}.
\end{corollary}
\begin{proof}[Proof of Corollary~\ref{cor:cartan-failure} and Theorem C]
First suppose $\C$ is unimodular and $C$ is singular, with $s$ distinct eigenvalues. Its minimal polynomial over $\mathbb Q$ is
$t b(t)$ with $b(0)\ne0$ and $\deg b=s-1$, since $C$ is symmetric with
integer entries. Put $g=b/b(0)$ and $w=\langle g(C)e_0,e_0\rangle$.
The matrix $g(C)$ projects onto $\ker C$, so $w=\nu_a(\{0\})>0$.
The polynomial $tg-w$ takes the value $-w$ at every Cartan eigenvalue and has derivative $1$ at zero. Formula~\eqref{eq:cartan-hermite} gives
\[
                          L_P((tg-w)^2)=-w^2 .
\]
The coefficients are rational. Clearing denominators and taking the
positive and negative coefficient parts therefore gives 2 actual sums
of powers of $P$ as in Theorem~\ref{thm:finite-test}.
Now assume the stated natural isomorphism. If $P$ has a Hamburger measure, Proposition~\ref{prop:fusion-type}(5) gives unimodularity. The preceding negative square rules out a singular Cartan matrix. By \cite[Theorem~6.6.1]{EGNO} the Cartan matrix of a nonsemisimple category with this natural isomorphism is singular. Hence $\C$ is semisimple. Conversely, in a semisimple category $P=\one$ and its law is $\delta_1$. Every self-dual object then has a law by \eqref{eq:gram}. In a unimodular category $P$ is itself self-dual, so tensor-power Cauchy--Schwarz implies its Hamburger positivity by Theorem~\ref{thm:main}.
\end{proof}

In the unimodular case the shift has a direct measure interpretation. If $\mu_P$ exists, then $t\,d\mu_P=d\nu_a$ by uniqueness on compact support, so the shifted measure has no atom at zero. Formula~\eqref{eq:cartan-hermite} explains what happens to the obstruction upon doubling. Substituting $f(t)=q(t^2)$ kills the derivative at zero, and the evaluations at $\lambda$ and $-\lambda$ combine. Proposition~\ref{prop:doubling-cartan} computes the resulting law when its masses are nonnegative.

There is a separate numerical observation. In every unimodular finite tensor category, $P^*\simeq P$ and $c_2(P)=h(P,P)=[P:\one]$. A nonsimple $P$ contains distinct unit factors at the top and the socle. Thus $c_2(P)=1$ is equivalent to $P=\one$, hence to semisimplicity. More generally $c_1(P)=h(\one,P)$ is $1$ precisely when the simple socle of $P$ is $\one$, that is, when $\C$ is unimodular. Consequently $(c_1(P),c_2(P))=(1,1)$ characterizes semisimplicity in every finite tensor category over $\CC$. The value $c_2(P)\ge2$ is compatible with positive moment data.

For example, consider the quaternion bosonization constructed in Example~\ref{ex:q-eight}. Its unit cover has class $a=2\cdot1+[V]$ in the representation ring of $Q_8$, where $V$ is the two-dimensional simple. The character of $a$ takes values $4,0,2$ on $1,-1,Q_8\setminus\{1,-1\}$, respectively. Thus
\begin{equation}\label{eq:quaternion-law}
 \nu_a=\tfrac18(\delta_0+6\delta_2+\delta_4),\qquad
 (c_0(P),\ldots,c_6(P))=(1,1,2,5,14,44,152).
\end{equation}
The first five entries agree with the rate-one free Poisson law. The Hankel determinants of sizes two, three and four are $1,1,-3$. More explicitly, take $g=(t-2)(t-4)/8$ and $w=1/8$ in the proof of Corollary~\ref{cor:cartan-failure}. Clearing the denominator yields
\[
          f(t)=t(t-2)(t-4)-1,\qquad L_P(f^2)=-1 .
\]

\begin{example}[A witness of least possible degree]\label{ex:exterior-witness}
Let $V$ have even dimension $m\ge2$ and take the bosonization
$H=\bigwedge V\rtimes\CC[C_2]$, with generator $g$ acting by $-1$ on
$V$ and $\Delta(v)=v\otimes1+g\otimes v$. The ideal of positive
exterior degree is nilpotent and the quotient is $\CC[C_2]$. Thus the
2 simples have classes $1,s$ with $s^2=1$. The projective cover
$P=H(1+g)/2$ of the unit has successive exterior degrees as a filtration.
Even and odd degrees contribute trivial and sign factors respectively.
Its socle is the top exterior power, trivial since $m$ is even.
Duality therefore fixes $P$ and the category is unimodular. With $d=2^m$,
\[
 a=(d/2)(1+s),\qquad \nu_a=(\delta_0+\delta_d)/2,\qquad
 (c_0,\ldots,c_4)=(1,1,d/2,d^2/2,d^3/2).
\]
The determinants of $H_1(P)$ and $H_2(P)$ are $d/2-1>0$ and $-d^3/8$.
Hence the Hom inequality holds on all sums of $\one,P$. At degree 2
take $A=P$ and $B=\one\oplus P^{\otimes2}$. Their Hom dimensions are
$h(A,A)=d/2$, $h(A,B)=1+d^2/2$, and $h(B,B)=1+d+d^3/2$, giving
\[
                h(A,B)^2-h(A,A)h(B,B)=1+\frac{d^2-d}{2}>0.
\]
\end{example}

There are two simples in this example, so it attains the degree bound $r$ in Theorem~\ref{thm:finite-test}. The Sweedler category has the same ordinary ring $\ZZ[C_2]$ and has measures for all its self-dual objects by Proposition~\ref{prop:sweedler}. Its unit cover illustrates the nonunimodular case of Theorem C. Here $a=1+s$ and $d=s$, so $c_{n+1}(P)=\tau(s(1+s)^n)$ gives $(c_0,\ldots,c_4)=(1,0,1,2,4)$ and
\[
                         L_P((t^2-2t-1)^2)=-1.
\]
This $P$ is not self-dual, so its negative square is compatible with tensor-power Cauchy--Schwarz for the category.

\subsection{Recovering Cartan spectral data from the moments}
Assume again that $\C$ is unimodular and of fusion type. The shifted sequence of $P$ always has a positive measure $\nu_a$, even
when the original sequence fails Hamburger positivity. This gives a
finite inverse construction, analogous to the orthogonalization in
Section~\ref{sec:hankel}.

\begin{proposition}\label{prop:cartan-recovery}
Under the hypotheses of Theorem~\ref{thm:cartan}, the integers
$c_1(P),\ldots,c_{2r}(P)$ determine the distinct eigenvalues of the
Cartan matrix, their spectral masses at $e_0$, and the minimal polynomial
of the matrix.
\end{proposition}
\begin{proof}
Put $m_j=c_{j+1}(P)$ and $s=\operatorname{rank}(m_{i+j})_{0\le i,j<r}$.
It is a Gram matrix in $L^2(\nu_a)$. By visibility and Vandermonde, its rank $s$ counts the distinct eigenvalues and its leading $s\times s$ block is invertible. Solve $\sum_{j=0}^{s-1}u_jm_{i+j}=-m_{i+s}$ for $0\le i<s$. The product over the distinct eigenvalues of $(t-\lambda)$ solves this system as a monic polynomial, hence equals $t^s+\sum_{j<s}u_jt^j$. It is the minimal polynomial by symmetry. Vandermonde on its roots gives the masses from $m_0,\ldots,m_{s-1}$. The largest index required is at most $2r-1$.
\end{proof}

For a given table, write $a=\sum_i a_i b_i$ in the simple basis with $b_0=1$. As $a=a^*$, the first shifted moments give
\begin{equation}\label{eq:cartan-first}
                   a_0=c_2(P),\qquad \sum_i a_i^2=c_3(P).
\end{equation}
There are finitely many integral candidates with coefficients $\geq0$. By the shifted moment formula, a central self-dual candidate agrees with the full sequence exactly when its regular matrix has the reconstructed spectral data at $e_0$.

\begin{corollary}\label{cor:rank-two-recovery}
If $R$ has rank two with simple basis $1,b$, the two integers $c_2(P)$ and $c_3(P)$ determine the full Cartan matrix:
\[
 C=c_2(P)I+\sqrt{c_3(P)-c_2(P)^2}\,N_b .
\]
\end{corollary}
\begin{proof}
The second coordinate in \eqref{eq:cartan-first} is an integer $\geq0$, hence is the indicated square root. Apply $C=N_a$.
\end{proof}
For Example~\ref{ex:exterior-witness} this recovers $C=\tfrac d2\left(\begin{smallmatrix}1&1\\1&1\end{smallmatrix}\right)$ even though $P$ has no Hamburger measure.

\subsection{Third and fourth moments detect semisimple products}\label{sec:chevalley}
Low moments can also compare a nonsemisimple category to its ordinary multiplication table. Recall that $\C$ has the \emph{Chevalley property} if the tensor product of any 2 simples is semisimple \cite[\S4.12]{EGNO}. In this case the semisimple objects themselves realize that table in a fusion category.

\begin{proposition}[Moment characterization of the Chevalley property]\label{prop:chevalley}\label{prop:chevalley-moments}
Let $\C$ be a finite complex tensor category. Put $U=\bigoplus_iS_i$, one copy of each simple, and $b_i=[U^{\otimes2}:S_i]$. Then
\[
 c_3(U)\le\sum_i b_i,\qquad c_4(U)\le\sum_i b_i^2.
\]
Equality in either inequality holds precisely when $\C$ has the Chevalley property. If $\C$ is of fusion type and $u=[U]$, the 2 upper bounds are $\tau(u^3)$ and $\tau(u^4)$, respectively.
\end{proposition}
\begin{proof}
Duality permutes the simples, so $U\simeq U^*$. Write $\ell$ for composition length. For any finite-length object $M$, put $m_i=[M:S_i]$ and $s_i=h(S_i,M)$. Schur's lemma gives $\sum_i s_i=\ell(\operatorname{soc}M)\le\ell(M)=\sum_i m_i$, with equality exactly when $M$ is semisimple. Rigidity therefore gives
\[
 c_3(U)=h(U,U^{\otimes2})=\ell(\operatorname{soc}U^{\otimes2})
 \le\ell(U^{\otimes2})=\sum_i b_i.
\]
Every $S_i\otimes S_j$ is a direct summand of $U^{\otimes2}$, so equality is equivalent to the Chevalley property.

Applying left exactness of $\Hom(-,M)$ to a composition series gives $h(M,M)\le\sum_i m_i s_i\le\sum_i m_i^2$. The last inequality is strict unless $M$ is semisimple, in which case both are equalities. Apply this to $M=U^{\otimes2}$ and use $c_4(U)=h(M,M)$. In fusion type, $u=u^*$ and orthonormality of the simple basis give $\sum_i b_i=\tau(u^*u^2)=\tau(u^3)$ and $\sum_i b_i^2=\tau((u^2)^*u^2)=\tau(u^4)$.
\end{proof}

In Example~\ref{ex:exterior-witness}, $U=\one\oplus s$ gives $c_3(U)=2+2=4$ and $c_4(U)=2^2+2^2=8$, while the unit cover has a negative Hankel determinant. The 2 measurements distinguish semisimple products from semisimplicity of the category.

Now suppose a finite category realizes the rank-2 rule $x^2=1+mx$ with $m\ge2$, and write $X$ for the simple of class $x$. The unit is a summand of $X^{\otimes2}$: it is both a subobject and a quotient by rigidity, and it has composition multiplicity one. Thus $X^{\otimes2}=\one\oplus M$ where $[M]=mx$. Here $c_3(X)=h(X,M)\le m$. Equality would split $M$ into $m$ copies of $X$, giving the Chevalley property. Ostrik's classification \cite[Main Theorem]{O2} excludes a fusion category with this table. Consequently any such realization would have
\[
             c_3(X)<m,\qquad c_4(X)<1+m^2 .
\]
The second inequality follows by applying the proof above to $X^{\otimes2}$. For $m=2$ the object $M$ must be a nonsplit self-extension of $X$. Its socle and top are $X$, and its endomorphism space has basis the identity and the composite of its quotient map with its socle embedding. Indeed an endomorphism acting trivially on the top has image in the socle and factors through the simple quotient. Thus $(c_0,\ldots,c_4)=(1,0,1,1,3)$, the same prefix as for the two-dimensional simple of $\Rep(S_3)$. These constraints require neither braiding nor pivotality and leave the existence of a finite realization open.

\begin{question}\label{q:fusion-chevalley}
Does every finite complex tensor category of fusion type have the Chevalley property? Equivalently, for $U=\bigoplus_iS_i$ and $u=[U]$, is $c_3(U)=\tau(u^3)$ automatic?
\end{question}
This asks whether the semisimple objects in every finite categorification of a fusion ring $R$ already give a fusion categorification of $R$. The cubic deficit counts the composition factors of $U^{\otimes2}/\operatorname{soc}(U^{\otimes2})$. A positive answer would exclude all finite categorifications of rings without fusion categorification, including $x^2=1+mx$ for $m\ge2$.

In a possible counterexample let $M$ be an indecomposable nonsimple summand of $S\otimes T$ for simples $S,T$. If $\one$ occurs in this product, then $T\simeq S^*$ and its unique copy splits off by Proposition~\ref{prop:fusion-type}~(1). Hence $\one$ does not occur in $M$. Choose an indecomposable nonsimple subquotient of $M$ of least length and take a simple quotient of it. Its kernel is semisimple by minimality. Pushing out along a projection of the kernel onto a suitable simple summand gives a nonsplit extension $0\to Z\to E\to Z'\to0$ with $Z,Z'\ne\one$. Exact rigid adjunction gives $\Ext^1(\one,Z\otimes Z'^*)\ne0$. The long exact sequence along a composition series then gives a constituent $W$ of $Z\otimes Z'^*$ with $\Ext^1(\one,W)\ne0$. Here $W\ne\one$ by \cite[Theorem~4.4.1]{EGNO}.

\subsection{Zero divisors and the field criterion}
In every finite tensor category of fusion type the Cartan matrix is the matrix of multiplication by $a=[P]$. This gives a ring-theoretic consequence of Cartan singularity.

\begin{theorem}\label{thm:field}
Let $\C$ be a finite tensor category over $\CC$ of fusion type, with Grothendieck ring $R$, such that the identity and the double dual are isomorphic as additive functors. If $\C$ is not semisimple, then the class $a=[P(\one)]$ is a zero divisor in $R$. In particular, if $R\otimes\mathbb Q$ is a field, then $\C$ is semisimple.
\end{theorem}

\begin{proof}
If $\C$ is not semisimple, its Cartan matrix is singular by \cite[Theorem~6.6.1]{EGNO}, as in the proof of Corollary~\ref{cor:cartan-failure}. Since $C=N_a$ has integer entries, there is a nonzero $v\in R$ with $av=0$. In a field this forces $a=0$, which is impossible since $\tau(a)=[P:\one]\geq1$.
\end{proof}

The assumption on the double dual holds when $\C$ is pivotal, and also when $\C$ is braided, by the Drinfeld isomorphism recalled in Section~\ref{sec:categorification}. The theorem applies to every fusion ring of rank 2 other than $\ZZ[C_2]$, since the ring with $x^2=1+mx$, $m\geq1$, has rational form $\mathbb Q(\sqrt{m^2+4})$. It also applies to the ring $\widetilde K_l$ of Section~\ref{sec:verlinde} when $2l+3$ is prime.

\section{Braidings, twists and categorification}\label{sec:categorification}

We prove Theorem D by structural arguments about extensions, transparent objects and the Radford isomorphism. These exclusions of nonsemisimple realizations complement the analytic criterion of Theorem C. They will also apply when the moment data in Theorem~\ref{thm:small-recognition} identify an even Verlinde ring in a finite category.

Etingof observed that the near-group rules $gX=Xg=X$, $X^2=X+\sum_{g\in G}g$ with $|G|>2$ already exclude a finite realization: the invertible factors split off $X^2$, leaving the simple $X$, so the semisimple objects would provide a fusion categorification of a ring known to admit none \cite{MOEtingof}. For the braided case below we also use his observation on pivotal factorizable categories of fusion type.

Throughout, $\C$ is a finite tensor category over $\CC$. Recall that $\C$ has the Chevalley property if tensor products of simples are semisimple (Section~\ref{sec:chevalley}). The semisimple objects then form a fusion category $\mathcal D\subset\C$, with $\Gr(\mathcal D)=\Gr(\C)$, so $\C$ is of fusion type. Recall also that the \emph{Drinfeld center} $\mathcal Z(\mathcal D)$ has as objects the pairs $(W,\sigma)$, where $W\in\mathcal D$ and $\sigma$ is a half-braiding, a natural family of isomorphisms $\sigma_Y:Y\otimes W\to W\otimes Y$ with $\sigma_{Y\otimes Y'}=(\sigma_Y\otimes\mathrm{id})(\mathrm{id}\otimes\sigma_{Y'})$. With this orientation the center braiding from $(W,\sigma)\otimes(V,\tau)$ to $(V,\tau)\otimes(W,\sigma)$ is $\tau_W$. The center is a braided fusion category \cite[\S\S7.13 and~8.5]{EGNO}.

\subsection{The extension object}
We use the Yoneda description of $\Ext^1$ from Section~\ref{sec:ext}, writing $g_*\xi$ for the pushout of an extension $\xi$ along a morphism $g$ of its subobject, and $f^*\xi$ for the pullback along a morphism $f$ to its quotient. Since tensoring is exact, an extension $\xi\in\Ext^1(A,B)$ and an object $Y$ give $\xi\otimes Y\in\Ext^1(A\otimes Y,B\otimes Y)$ and $Y\otimes\xi$.

\begin{proposition}\label{prop:ext-object}
Let $\C$ have the Chevalley property, with semisimple part $\mathcal D$.
\begin{enumerate}
\item There are an object $W\in\mathcal D$ and an extension $\omega\in\Ext^1_\C(\one,W)$ such that $g\mapsto g_*\omega$ is a bijection $\Hom_{\mathcal D}(W,Z)\to\Ext^1_\C(\one,Z)$ for every $Z\in\mathcal D$.
\item For $Y\in\mathcal D$ there is a unique isomorphism $\sigma_Y:Y\otimes W\to W\otimes Y$ with $(\sigma_Y)_*(Y\otimes\omega)=\omega\otimes Y$, and $(W,\sigma)$ is an object of $\mathcal Z(\mathcal D)$.
\item We have $\Hom_{\mathcal D}(W,\one)=0$, and $\C$ is semisimple if and only if $W=0$.
\end{enumerate}
\end{proposition}

\begin{proof}
(1) The functor $Z\mapsto\Ext^1_\C(\one,Z)$ is linear on the semisimple category $\mathcal D$, with finite-dimensional values, so it is represented by $W=\bigoplus_S\Ext^1_\C(\one,S)^*\otimes S$, the sum being over the simples. The extension $\omega$ is the universal element.

(2) For $Y,Z\in\mathcal D$ consider the maps
\begin{align*}
 a_Y&:\Hom(W\otimes Y,Z)\to\Ext^1(Y,Z),& g&\mapsto g_*(\omega\otimes Y),\\
 b_Y&:\Hom(Y\otimes W,Z)\to\Ext^1(Y,Z),& g&\mapsto g_*(Y\otimes\omega).
\end{align*}
Both are bijective. Indeed, the exact functors $-\otimes Y$ and $-\otimes Y^*$ are adjoint, which gives an isomorphism $\Ext^1(Y,Z)\simeq\Ext^1(\one,Z\otimes Y^*)$, with inverse $\eta\mapsto(\mathrm{id}_Z\otimes\mathrm{ev}_Y)_*(\eta\otimes Y)$. If $g$ corresponds to $g'\in\Hom(W,Z\otimes Y^*)$ under \eqref{eq:rigidity}, then the triangle identity gives $(\mathrm{id}_Z\otimes\mathrm{ev}_Y)_*(g'_*\omega\otimes Y)=g_*(\omega\otimes Y)$. So $a_Y$ is the composite of 3 bijections, and similarly for $b_Y$, using right duals. Now $b_Y^{-1}a_Y$ is a bijection $\Hom(W\otimes Y,Z)\to\Hom(Y\otimes W,Z)$, natural in $Z$. By the Yoneda lemma it is the composition with an isomorphism $\sigma_Y$, which is characterized by the stated equality, the case $g=\mathrm{id}$.

For $f:Y\to Y'$ the morphisms $(\mathrm{id}_W\otimes f)\sigma_Y$ and $\sigma_{Y'}(f\otimes\mathrm{id}_W)$ both push $Y\otimes\omega$ to $f^*(\omega\otimes Y')$, so they are equal by injectivity of $b_Y$. In the same way $(\sigma_Y\otimes\mathrm{id})(\mathrm{id}\otimes\sigma_{Y'})$ pushes $Y\otimes Y'\otimes\omega$ to $\omega\otimes Y\otimes Y'$, so it equals $\sigma_{Y\otimes Y'}$.

(3) By (1) this space is $\Ext^1_\C(\one,\one)$, which vanishes for a finite tensor category over $\CC$ \cite[Theorem~4.4.1]{EGNO}. If $W=0$ then $\Ext^1(\one,S)=0$ for every simple $S$. By induction on the length $\Ext^1(\one,-)=0$, so $\one$ is projective and $\C$ is semisimple \cite[Corollary~4.2.13]{EGNO}. The converse is clear.
\end{proof}

For the Sweedler algebra, $W$ is the sign character, and for the bosonization of $\bigwedge V$ in Example~\ref{ex:exterior-witness} it is the sum of $\dim V$ copies of the sign character.

\subsection{Braided categories with semisimple products}
A \emph{braiding} on $\C$ is a natural family of isomorphisms $c_{X,Y}:X\otimes Y\to Y\otimes X$ satisfying the 2 hexagon identities \cite[\S8.1]{EGNO}. An object $X$ is \emph{transparent} if $c_{Y,X}c_{X,Y}=\mathrm{id}$ for every $Y$, and the transparent objects form the \emph{M\"uger center} $\C'$. This is a symmetric tensor category. A braided category is called \emph{nondegenerate} if $\C'$ consists of the multiples of $\one$, which is equivalent to factorizability \cite[Theorem~1.1]{Shimizu}. At the other extreme, the symmetric fusion categories are, by Deligne's theorem, the categories $\Rep(G,z)$ of representations of a finite group $G$, with the usual flip modified by the sign $-1$ on pairs of simples on which a fixed central element $z$ with $z^2=1$ acts by $-1$. We call these simples \emph{odd}. The category is \emph{Tannakian} if there are none, so that it is $\Rep(G)$ with the usual flip \cite[\S9.9]{EGNO}.

A braided category has a natural, in general nonmonoidal, \emph{Drinfeld isomorphism} $u_X:X\to X^{**}$, which satisfies $u_{X\otimes Y}=(u_X\otimes u_Y)(c_{Y,X}c_{X,Y})^{-1}$ \cite[\S8.9]{EGNO}. In particular the identity and the double dual are isomorphic as additive functors, and Theorem~\ref{thm:field} applies.

\begin{lemma}\label{lem:tannakian}
Let $\C$ be a braided finite tensor category over $\CC$, and let $\mathcal E\subset\C$ be a Tannakian fusion subcategory. Then $\Ext^1_\C(\one,S)=0$ for every simple object $S$ of $\mathcal E$.
\end{lemma}

\begin{proof}
Let $\mathcal E\simeq\Rep(G)$, and let $A\in\mathcal E$ be the algebra of functions on $G$. It is a commutative algebra in $\C$, and its modules in $\C$ form a finite tensor category $\C_G$, with unit $A$, the de-equivariantization \cite[\S4.3]{NP}. The free module functor $X\mapsto A\otimes X$ is exact and left adjoint to the exact forgetful functor, so $\Ext^1_\C(\one,A)\simeq\Ext^1_{\C_G}(A,A)$. This vanishes by \cite[Theorem~4.4.1]{EGNO}. Since $S$ is a direct summand of the regular representation $A$, the claim follows.
\end{proof}

\begin{theorem}\label{thm:braided-chevalley}
Let $\C$ be a braided finite tensor category over $\CC$ with the Chevalley property, and let $\mathcal D$ be its braided fusion subcategory of semisimple objects. Then the extension object $W$ is transparent in $\mathcal D$, and all its simple summands are odd objects of the M\"uger center $\mathcal D'$. In particular, if $\mathcal D'$ is Tannakian, for instance if $\mathcal D$ is nondegenerate, then $\C$ is semisimple.
\end{theorem}

\begin{proof}
Let $0\to W\to E\to\one\to0$ represent $\omega$, and let $Y\in\mathcal D$. By naturality the braiding $c_{Y,E}$ is an isomorphism from the extension $Y\otimes\omega$ to the extension $\omega\otimes Y$, which induces $c_{Y,W}$ on the subobjects and $c_{Y,\one}=\mathrm{id}$ on the quotients. Thus $(c_{Y,W})_*(Y\otimes\omega)=\omega\otimes Y$, and $\sigma_Y=c_{Y,W}$ by Proposition~\ref{prop:ext-object}. The same argument with the inverse braiding gives $\sigma_Y=c_{W,Y}^{-1}$. Hence $c_{W,Y}c_{Y,W}=\mathrm{id}$, and $W\in\mathcal D'$. Write $\mathcal D'=\Rep(G,z)$. Its simples which are not odd form the Tannakian category $\Rep(G/\langle z\rangle)$, and they do not occur in $W$, by Lemma~\ref{lem:tannakian} and the definition of $W$. If there are no odd simples then $W=0$.
\end{proof}

The theorem is sharp. For $\bigwedge V\rtimes\CC[C_2]$ with its super flip braiding, $\mathcal D=\mathcal D'=\Rep(C_2,z)$ is the category of super vector spaces, and $W$ is a sum of copies of the odd line.

\subsection{Braided categories of fusion type}
Without the Chevalley property the extension object is not available, and we use instead the structure theory of factorizable categories.

\begin{lemma}[The coend pairing and the right partial trace]\label{lem:coend-trace}
Let $\C$ be a braided finite tensor category, $F=\int^Y Y^*\otimes Y$ its coend, $\iota_Y$ the dinatural maps and $\omega_F$ its Hopf pairing. For $\phi:X\to X^{**}$ put
\[
 \chi_{X,\phi}=\iota_{X^*}(\phi\otimes\mathrm{id}_{X^*})\mathrm{coev}_X,
 \qquad \Omega(f)=\omega_F(f\otimes\mathrm{id}_F).
\]
Under $\Hom(F,\one)\cong\End(\mathrm{Id}_{\C})$, the functional $\Omega(\chi_{X,\phi})$ corresponds on $M$ to the right partial trace of $h=c_{X,M}c_{M,X}$,
\[
 \Phi(X,\phi)_M =
 (\mathrm{id}_M\otimes\mathrm{ev}_{X^*})
 (\mathrm{id}_M\otimes\phi\otimes\mathrm{id}_{X^*})
 (h\otimes\mathrm{id}_{X^*})
 (\mathrm{id}_M\otimes\mathrm{coev}_X).
\]
\end{lemma}
\begin{proof}
The defining pairing on $\iota_A\otimes\iota_B$ is
$(\mathrm{ev}_A\otimes\mathrm{ev}_B)
(\mathrm{id}_{A^*}\otimes c_{B^*,A}c_{A,B^*}\otimes\mathrm{id}_B)$
\cite[(3.4) and (5.1)]{Shimizu}. Take $A=X^*$, $B=M$. The middle endomorphism on $X^*\otimes M^*$ is $h^*$, since duality reverses the order of the crossings. With $\kappa=(\phi\otimes\mathrm{id}_{X^*})\mathrm{coev}_X$ this gives
\[
 \Omega(\chi_{X,\phi})\iota_M=
 (\mathrm{ev}_{X^*}\otimes\mathrm{ev}_M)
 (\mathrm{id}_{X^{**}}\otimes h^*\otimes\mathrm{id}_M)
 (\kappa\otimes\mathrm{id}_{M^*\otimes M}).
\]
Transposing $h^*$ across the evaluations by
$\mathrm{ev}_{M\otimes X}(h^*\otimes\mathrm{id})
=\mathrm{ev}_{M\otimes X}(\mathrm{id}\otimes h)$
and cancelling the resulting evaluation--coevaluation pair gives
$\mathrm{ev}_M(\mathrm{id}_{M^*}\otimes\Phi(X,\phi)_M)$. This is the defining correspondence for $\Hom(F,\one)\cong\End(\mathrm{Id}_{\C})$.
\end{proof}

\begin{lemma}[Descent of sphericalized traces]\label{lem:trace-descent}
Let $\C$ be a braided unimodular finite tensor category of fusion type, with $r$ simples and Radford isomorphism $\delta:\mathrm{Id}\to(-)^{****}$. Let $\widetilde\C$ be its sphericalization and $\epsilon=(\one,-1)$ its sign line. Then
\[
 Q=(\Gr(\widetilde\C)\otimes\CC)/(\epsilon+1)
\]
is a semisimple algebra of dimension $r$, and the partial traces of Lemma~\ref{lem:coend-trace} define an algebra map
$\Phi:Q\to\End(\mathrm{Id}_{\C})$. If $\C$ is nondegenerate, this map is injective.
\end{lemma}
\begin{proof}
Unimodularity supplies $\delta$ \cite[Theorem~7.19.1]{EGNO}. The sphericalization consists of pairs $(X,\phi)$ with $\phi:X\to X^{**}$ satisfying $\phi^{**}\phi=\delta_X$; morphisms commute with the $\phi$'s and tensor products use $\phi\otimes\psi$. It is a pivotal finite tensor category \cite[Definition~5.1 and Proposition~5.5]{CJNP}. By Proposition~\ref{prop:fusion-type}(1), normalize an isomorphism $\phi_i:S_i\to S_i^{**}$ for each simple $S_i$ so that $(S_i,\pm\phi_i)$ are objects. These are all the simples: if $(X,\phi)$ is simple and $H=\Hom_\C(S_i,X)\ne0$, then $f\mapsto\phi^{-1}f^{**}\phi_i$ is an involution of $H$, and an eigenvector gives a nonzero morphism $(S_i,\pm\phi_i)\to(X,\phi)$.

The underlying product of 2 simples has unit multiplicity at most one. Coevaluation supplies the unit constituent for a simple and its dual; tensoring the dual by $\epsilon$ replaces that constituent by $\epsilon$. Thus $\widetilde\C$ is of fusion type. The central sign line exchanges the two lifts of every simple, so $Q$ has dimension $r$, with basis the images of $(S_i,\phi_i)$. It is a direct factor of the semisimple fusion algebra, cut out by the central idempotent $(1-\epsilon)/2$.

The braiding lifts to $\widetilde\C$. For $A=(X,\phi)$ write $\widehat\Phi(A)$ for the right partial trace of the double braiding upstairs and $U:\widetilde\C\to\C$ for forgetful. The right-trace versions of \cite[Lemmas~2.1 and~2.4]{GR}, obtained by duality, give additivity on exact sequences and multiplicativity in $A$. For $M\in\C$, equivariant induction has $UI(M)=M\oplus M^{**}$. Let $i_M,p_M$ be the first inclusion and projection. The underlying formula of Lemma~\ref{lem:coend-trace} is natural even for morphisms in $\C$, so
\[
 U\bigl(\widehat\Phi(A)_{I(M)}\bigr)i_M=i_M\Phi(X,\phi)_M,\qquad
 \Phi(X,\phi)_M=p_MU\bigl(\widehat\Phi(A)_{I(M)}\bigr)i_M.
\]
The first identity makes the summand invariant. Restriction therefore preserves compositions as well as sums, descending the algebra identities to every $M$. Since $\Phi(\epsilon)=-\mathrm{id}$, the map factors through $Q$.

If $\C$ is nondegenerate, its coend map $\Omega$ is injective \cite[Theorem~1.1]{Shimizu}. The $r$ characters $\chi_{S_i,\phi_i}$ are independent: $\bigoplus_SS^*\otimes S\to F$ is a monomorphism, and each character is nonzero in the summand indexed by $S_i^*$ \cite[Lemma~5.5 and the proof of Lemma~5.6]{Shimizu}. Lemma~\ref{lem:coend-trace} identifies their images under $\Omega$ with $\Phi(S_i,\phi_i)$, proving injectivity.
\end{proof}

\begin{theorem}\label{thm:nondegenerate}
Let $\C$ be a braided finite tensor category over $\CC$ of fusion type. If $\C$ is nondegenerate, or more generally if its M\"uger center $\C'$ is Tannakian, then $\C$ is semisimple.
\end{theorem}
\begin{proof}
Suppose first that $\C$ is nondegenerate, with $r$ simples. It is unimodular \cite[Proposition~8.10.10]{EGNO}. For pivotal $\C$, the semisimplicity conclusion is Etingof's observation on \cite[Corollary~8.7]{GR}, \cite{MOEtingof}; Lemma~\ref{lem:trace-descent} removes pivotality. Its injection $Q\to\End(\mathrm{Id}_\C)$ makes $Q$ commutative, since natural endomorphisms of the identity commute. The image therefore contains $r$ nonzero orthogonal idempotents. Hence $\C$ has at least $r$ blocks, each with one simple. Since $\Ext^1(\one,\one)=0$ \cite[Theorem~4.4.1]{EGNO}, the unit block consists of multiples of $\one$. The unit is projective, so $\C$ is semisimple.

Now suppose $\C'=\Rep(G)$. Let $F:\C\to\C_G$ be de-equivariantization. Equivalently, $\C\simeq(\C_G)^G$ and $F$ forgets the equivariant structure \cite[\S4.3]{NP}. For a simple $S$ of $\C$, the socle of $FS$ is a nonzero equivariant subobject because each $g\in G$ preserves simples. Thus $FS$ is semisimple. Conversely, for a simple $V$ of $\C_G$ the induced equivariant object has underlying object $\bigoplus_{g\in G}gV$. This object is semisimple also in the equivariantization: averaging a retraction of any subobject makes the retraction equivariant. Consequently $V$ occurs in $FS$ for some simple $S$. If $\one$ occurs in $FS$, its isotypic component is equivariant and hence equals $FS$; these $S$ are precisely the irreducibles $\pi$ of $\Rep(G)$, with $F\pi=(\dim\pi)\one$.

Let $A=\bigoplus_\pi(\dim\pi)\pi$ be the regular object of $\Rep(G)$. For simples $S,T$ in $\C$, exactness, fusion-ring reciprocity and induction adjunction give
\[
 [FS\otimes(FT)^*:\one]
 =\sum_\pi(\dim\pi)[S\otimes T^*:\pi]
 =[A\otimes S:T]=\dim\Hom_{\C_G}(FT,FS).
\]
Indeed induction of $FS$ is $A\otimes S$, which is semisimple by the same averaging argument. Expanding both ends over the simple summands $U,V$ of $FS,FT$, respectively, gives a weighted sum of the nonnegative integers $[U\otimes V^*:\one]-\delta_{UV}$ equal to zero. Every simple occurs in such an image, so each difference is zero. Thus $\C_G$ is of fusion type. Its M\"uger center is trivial by \cite[Theorem~7.5]{NP}, hence it is semisimple by the first case. Equivariant averaging gives semisimplicity of $\C$.
\end{proof}

\begin{theorem}\label{thm:braided-exclusion}
Every braided categorification over $\CC$ of a fusion ring of rank 3 is semisimple. The same holds in rank 2, except for the ring $\ZZ[C_2]$.
\end{theorem}

\begin{proof}
Let $\C$ be such a categorification, and assume that it is not semisimple. By Deligne's theorem \cite[\S9.11]{EGNO} the symmetric finite tensor category $\C'$ is the representation category of a finite supergroup, whose semisimple objects form a category $\Rep(G,z)$. If $z=1$, its parity action forces the odd Lie algebra of the supergroup to vanish, and in characteristic zero the center would be Tannakian. Hence $z\ne1$ by Theorem~\ref{thm:nondegenerate}. Since $G$ has at most 3 conjugacy classes and a central involution, $G=C_2$. Indeed, an abelian group with a central involution and at most 3 classes has order 2, and in the nonabelian case with 3 classes $|G|-2$ would be a class size dividing $|G|$, forcing $|G|\le4$. Thus $\C$ has a transparent odd line $f$, with $f\otimes f\simeq\one$. In rank 2 the ring is $\ZZ[C_2]$. In rank 3 let $S$ be the remaining simple, so that $f\otimes S\simeq S$. For a simple $T$ let $d(T)$ be the scalar of Proposition~\ref{prop:fusion-type}~(4) for $\phi=u_T$, which depends only on the isomorphism class of $T$. By transparency $u_{f\otimes S}=u_f\otimes u_S$, so $d(f\otimes S)=d(f)d(S)$, and $d(f)=-1$ for an odd line. Hence $d(S)=-d(S)$, contradicting $d(S)\neq0$.
\end{proof}

In rank 2 Example~\ref{ex:exterior-witness} realizes the exception with its super flip braiding. Thus the Ising and $\Rep(S_3)$ fusion tables force semisimplicity in the braided case. For $x^2=1+mx$ with $m\ge2$, the theorem and Ostrik's classification \cite[Main Theorem]{O2} exclude every braided finite realization. At rank 3 the category $\Rep(T_3)$ realizes $\ZZ[C_3]$ nonsemisimply without a braiding (Section~\ref{sec:hopf}). Finiteness also matters: $\Rep(S_3\times\Ga)$ is symmetric and nonsemisimple, with the same ordinary Grothendieck ring in the simple basis as $\Rep(S_3)$ (Example~\ref{ex:nil}). The proof shows more generally that a nonsemisimple braided category of fusion type, of any rank, has a nontrivial transparent fusion subcategory $\Rep(G,z)$ with $z\neq1$, and that a transparent odd line permutes the simples without fixed points.

A $\ZZ/2$-grading of a fusion ring with basis $(b_i)$ assigns degrees $\epsilon_i\in\{0,1\}$ such that $\epsilon_0=0$ and $N_{ij}^k>0$ implies $\epsilon_i+\epsilon_j=\epsilon_k$ modulo $2$. The grading is nontrivial if some degree is $1$.

For a basis element $b$ let $\Gamma_b$ be the undirected graph on the \emph{whole} basis of $R$, with an edge $S$--$T$ whenever $N_{bS}^{T}>0$. Keep loops. Frobenius reciprocity gives $N_{b^*}=N_b^{\mathsf T}$, so $\Gamma_{b^*}=\Gamma_b$.

\begin{theorem}[Parity and bipartite fusion graphs]\label{thm:graded-subring}\label{thm:bipartite}
If $R$ admits a nonsemisimple braided finite tensor categorification over $\CC$, it contains a based fusion subring with integral Frobenius--Perron dimensions and a nontrivial $\ZZ/2$-grading. For every odd basis element $b$ of this subring, $\Gamma_b$ is bipartite. In particular, if no nonunit basis element of $R$ has both integral Frobenius--Perron dimension and bipartite fusion graph, then every braided finite tensor categorification of $R$ over $\CC$ is semisimple.
\end{theorem}
\begin{proof}
As in the proof of Theorem~\ref{thm:braided-exclusion}, the semisimple objects in the M\"uger center form $\mathcal E=\Rep(G,z)$ with $z\ne1$ \cite[Theorem~9.11.9 and Exercise~9.11.13(i)--(ii)]{EGNO}. Its simples remain simple in $\C$, since transparency is closed under subquotients. The action of $z$ gives the asserted grading, which is nontrivial since the regular representation is faithful. The positive dimension character $f_S=\FPdim(S)$ restricts to ordinary dimension on $\mathcal E$ by uniqueness of the Frobenius--Perron character.

For any $X\in\C$ put $d_X=\mathrm{ev}_{X^*}(u_X\otimes\mathrm{id}_{X^*})\mathrm{coev}_X$. Naturality of the Drinfeld isomorphism and additivity of the quantum trace on short exact sequences \cite[Proposition~4.7.5]{EGNO} give $d_X=\sum_T[X:T]d_T$. For a simple $S$, Proposition~\ref{prop:fusion-type}(4) gives $d_S\ne0$. Let $b$ be an odd simple of $\mathcal E$. The signed flip gives $d_b=-f_b$. Transparency gives $u_{b\otimes S}=u_b\otimes u_S$ and hence $d_{b\otimes S}=d_b d_S$ by multiplicativity of the trace \cite[Proposition~4.7.3(3)]{EGNO}. Put $r_S=d_S/f_S$. Additivity of both traces and multiplicativity of $f$ give
\[
 -r_S=\sum_T\frac{N_{bS}^{T}f_T}{f_b f_S}\,r_T,
 \qquad \sum_T\frac{N_{bS}^{T}f_T}{f_b f_S}=1 .
\]
Choose a vertex $S$ maximizing $|r_S|$ in a connected component of $\Gamma_b$. All positive terms of this convex combination must have $r_T=-r_S$, by equality in the triangle inequality. The same argument for the odd dual $b^*$ applies to every incoming edge at $S$. Each neighbor again has maximal modulus, so we can repeat the argument along any path in the component. The two distinct values $r_S,-r_S$ give its two colors, and every edge joins opposite colors.
\end{proof}

The graph test also detects every ring having no nontrivially graded subring. Indeed, if $\Gamma_b$ is bipartite for a nonunit $b$, color the unit vertex even. The component of the unit consists of the basis elements in words in $b,b^*$, since reversing an edge amounts to multiplication by $b^*$. A constituent of a word of length $n$ has color $n\bmod 2$. Concatenating words shows that these colors grade the based subring generated by $b,b^*$, and $b$ is odd.

Call a fusion ring \emph{fusion-simple} if its only based fusion subrings are $\ZZ1$ and the whole ring.

\begin{corollary}[Fusion-simple rings]\label{thm:fusion-simple}
Let $R$ be a fusion-simple ring. Every braided finite tensor categorification of $R$ over $\CC$ is semisimple unless $R\simeq\ZZ[C_2]$. The exceptional ring admits nonsemisimple symmetric categorifications.
\end{corollary}
\begin{proof}
Take $\mathcal E=\Rep(G,z)$ as in the proof of Theorem~\ref{thm:graded-subring}. Its even subcategory $\mathcal E_0=\Rep(G/\langle z\rangle)$ is proper in $\mathcal E$. If $G/\langle z\rangle$ is nontrivial, $\Gr(\mathcal E_0)$ is a nontrivial proper based fusion subring of $R$, a contradiction. Thus $G=C_2$, and fusion simplicity applied to $\Gr(\mathcal E)=\ZZ[C_2]$ gives $R=\ZZ[C_2]$. The symmetric exterior-algebra examples supply the exception.
\end{proof}

A transparent odd simple can have dimension greater than one, even when every invertible simple is even.

\begin{example}[An odd extension object of dimension 2]\label{ex:q-eight}\label{ex:quaternion-witness}
Let $G=Q_8$, let $z=-1\in G$, and let $V$ be its irreducible 2-dimensional representation. Explicitly, generators $a,b$ act by
\[
 \rho(a)=\begin{pmatrix}i&0\\0&-i\end{pmatrix},\qquad
 \rho(b)=\begin{pmatrix}0&1\\-1&0\end{pmatrix}.
\]
Then $z$ acts by $-1$ and $\det V=1$. Form the 32-dimensional Hopf algebra $H=\bigwedge V\rtimes\CC[G]$, with $gvg^{-1}=\rho(g)v$ and
\begin{align*}
 \Delta(g)&=g\otimes g,& \Delta(v)&=v\otimes1+z\otimes v,\\
 \varepsilon(g)&=1,& \varepsilon(v)&=0,\\
 S(g)&=g^{-1},& S(v)&=-zv.
\end{align*}
Its triangular structure
\[
 R_H=\tfrac12(1\otimes1+1\otimes z+z\otimes1-z\otimes z)
\]
gives the super flip braiding. The ideal of positive exterior degree is the radical, with quotient $\CC[G]$. Hence $\Rep(H)$ is a nonsemisimple symmetric category of fusion type with the Chevalley property. We have $S^2=\operatorname{Ad}(z)$, so $z$ is a pivot. The unit cover $P=He$, where $e=|G|^{-1}\sum_{g\in G}g$, has exterior-degree factors $\one,V,\one$. Its socle is trivial, so $H$ is unimodular, and the extension object of Proposition~\ref{prop:ext-object} is $W=V$. Every linear character of $Q_8$ kills $z$, so this category has an odd 2-dimensional simple and has no odd line.

The exterior-degree factors give $[P]=2\cdot1+[V]$, as used in \eqref{eq:quaternion-law}. Thus the same category illustrates both the transparent odd extension object here and the failure of Hamburger positivity for the unit cover in Section~8.
\end{example}

\paragraph{Application to the multiplicity-free census.}\label{sec:parity-census}
Among the 353 multiplicity-free tables of ranks at most nine in \cite{Census,VSrings}, 303 are commutative. Table~\ref{tab:braided-census} gives the results for these 303 rings. For the subring test, enumerate all subsets containing the unit that are closed under duality and product constituents. On each subset solve $\epsilon_0=0$ and $\epsilon_i+\epsilon_j+\epsilon_k=0$ over $\mathbb F_2$ whenever $N_{ij}^k>0$. The subring has a nontrivial grading exactly when there is a nonzero solution. This test excludes a nonsemisimple braided finite categorification for 39 rings, or 41 on adjoining the other two rank-three tables from Theorem~\ref{thm:braided-exclusion}. We count the rank-one ring as fusion-simple. Corollary~\ref{thm:fusion-simple} applies to 18 of the 19 fusion-simple rings.

\begin{table}[htbp]\centering\small
\caption{Tests for the 303 commutative multiplicity-free rings. The last two rows count rings failing the respective necessary conditions for a nonsemisimple braided finite categorification.}\label{tab:braided-census}
\begin{tabular}{@{}lrrrrrrrrrr@{}}
\toprule
Rank & 1&2&3&4&5&6&7&8&9&Total\\
\midrule
Commutative rings &1&2&4&10&16&37&40&77&116&303\\
Fusion-simple rings &1&2&2&1&3&3&3&2&2&19\\
No parity-graded subring &1&1&2&3&3&8&5&6&10&39\\
No integral bipartite candidate &1&1&4&5&14&22&36&32&96&211\\
\bottomrule
\end{tabular}
\end{table}

For the graph test, inspect each nonunit $b$. Two-color $\Gamma_b$ on all components, rejecting any loop or odd cycle, and test whether the Perron root of $N_b$ is integral. For an exact test let $p(t)=\det(tI-N_b)$. An integer $k$ from $1$ to the maximum row sum is the Perron root precisely when $p(k)=0$ and all coefficients of $p(t+k)$ are nonnegative. Indeed, when $k$ is the Perron root, every real eigenvalue is at most $k$, and a nonreal conjugate pair $\lambda,\bar\lambda$ contributes $(t+k-\operatorname{Re}\lambda)^2+(\operatorname{Im}\lambda)^2$. Each shifted factor has nonnegative coefficients. Conversely a monic shifted polynomial with nonnegative coefficients has no positive real root, while the spectral radius of $N_b$ is a real eigenvalue. The program \texttt{census.py} computes both tests. A separate implementation, \texttt{check\_bipartite\_census.py}, verifies the graph counts on the same input.

Two rank-four tables illustrate the gain. In the Tambara--Yamagami ring of $C_3$ (source line 130), the three invertibles are even and the odd simple $X$ has dimension $\sqrt3$. Each nonunit invertible fixes $X$, giving a loop. In the table on source line 8759 the nontrivial invertible $g$ fixes both other simples $x,y$, and $xy=x+y$. The graph for $g$ has loops at $x,y$, that for $x$ has a loop at $y$, and that for $y$ has a loop at $x$. Thus both tables fail the graph test. Altogether, for 211 tables any braided finite categorification over $\CC$, if it exists, must be semisimple. The remaining 92 have at least one integral bipartite candidate. Their realizability is undecided by this test.

\subsection{The Radford isomorphism and the even Verlinde rings}\label{sec:verlinde}
We return to categories without braiding. Let $\C$ be a unimodular finite tensor category over $\CC$ with the Chevalley property, let $\mathcal D$ denote its fusion subcategory of semisimple objects, and let $\delta$ be the Radford isomorphism of $\C$, as in Lemma~\ref{lem:trace-descent}. Assume that the fusion category $\mathcal D$ has a spherical structure $j:\mathrm{Id}_{\mathcal D}\to(-)^{**}$ with $j^{**}j=\delta$ on $\mathcal D$. Then $\mathcal Z(\mathcal D)$ is a modular category, and every object $(W,\sigma)$ has a twist $\theta_{(W,\sigma)}\in\mathrm{Aut}(W)$, given by $\theta_{(W,\sigma)}=u_{(W,\sigma)}^{-1}j_W$, which acts by a scalar on each simple summand \cite[\S\S8.10 and~8.20]{EGNO}.

\begin{lemma}[Dualization of the universal extension]\label{lem:extension-dual}
For the extension object $(W,\sigma)$ and its universal class $\omega$, one has
\[
 \omega^{**}=(\nu_W)_*\omega,\qquad
 \nu_W=(\mathrm{ev}_W\otimes\mathrm{id}_{W^{**}})
 (\sigma_{W^*}^{-1}\otimes\mathrm{id}_{W^{**}})
 (\mathrm{id}_W\otimes\mathrm{coev}_{W^*}).
\]
Here $\nu_W:W\to W^{**}$ is the Drinfeld isomorphism for the inverse center braiding.
\end{lemma}
\begin{proof}
For $\xi\in\Ext^1(\one,Z)$ and $\eta\in\Ext^1(Y,\one)$ put
\[
 \xi^\flat=(\mathrm{ev}_Z)_*(Z^*\otimes\xi),\qquad
 \eta^\sharp=\mathrm{coev}_Y^*(\eta\otimes Y^*).
\]
With the usual dual exact sequences, $\xi^\flat=-\xi^*$ and $\eta^\sharp=-\eta^*$. For the first sign, represent $\xi$ by $0\to Z\xrightarrow{i}E\xrightarrow{p}\one\to0$ and form $K=(Z^*\otimes E)\times_{Z^*}E^*$. The exact tensor-product square gives an epimorphism
$(i^*\otimes\mathrm{id}_E,\mathrm{id}_{E^*}\otimes p):E^*\otimes E\to K$
with kernel $(p^*\otimes i)(Z)$, on which $\mathrm{ev}_E$ vanishes. Evaluation descends to $K$ and restricts to $(\mathrm{ev}_Z,\mathrm{id}_\one)$ on its kernel $(Z^*\otimes Z)\oplus\one$ over $Z^*$. This splits the pushout representing $\xi^\flat+\xi^*$, proving the sign. The dual argument gives the second sign, so $\omega^{**}=(\omega^\flat)^\sharp$.

Now $W^*\otimes\omega=(\sigma_{W^*}^{-1})_*(\omega\otimes W^*)$ by Proposition~\ref{prop:ext-object}, and pulling $\omega\otimes W^*\otimes W^{**}$ back along $\mathrm{coev}_{W^*}$ gives $(\mathrm{id}_W\otimes\mathrm{coev}_{W^*})_*\omega$. Combining these identities gives $(\omega^\flat)^\sharp=(\nu_W)_*\omega$. The displayed formula for $\nu_W$ is precisely the Drinfeld formula for the inverse braiding with our convention for $\sigma$.
\end{proof}

\begin{lemma}[Separation of signs]\label{lem:sign-separation}
Under the hypotheses of this subsection, suppose that
$(j_Z^{-1})_*(\xi^{**})=\xi$ for all $Z\in\mathcal D$ and $\xi\in\Ext^1_\C(\one,Z)$.
Then no block of the sphericalization contains both a positive lift $S^+=(S,j_S)$ and a negative lift $T^-=(T,-j_T)$. If one embedding of the field of spherical dimensions makes all simple dimensions positive, then $\C$ is semisimple.
\end{lemma}
\begin{proof}
An extension of $\one^+$ by $Z^-$ has underlying class $\xi$ with $\xi^{**}=(-j_Z)_*\xi$, by equivariance at the ends. The hypothesis gives $\xi=-\xi$, hence $\xi=0$. Averaging an underlying splitting over the order-two action makes it equivariant. Exact tensor adjunction gives
\[
 \Ext^1_{\widetilde\C}(S^+,T^-)
 \simeq\Ext^1_{\widetilde\C}
   \bigl(\one^+,T^-\otimes(S^+)^*\bigr)=0:
\]
the second argument splits into negative lifts by the Chevalley property and monoidality of $j$. Tensoring by the sign line gives vanishing in the opposite direction. Thus the projective cover $Q$ of $\one^+$ has only positive-lift composition factors.

Pivotal dimension is additive \cite[Proposition~4.7.5]{EGNO}, so
\[
 \dim(Q)=\sum_{S\ {\rm simple}}[Q:S^+]\,\dim_j(S).
\]
The common embedding makes this sum positive, since $Q$ has quotient $\one^+$. The nonzero pivotal trace of $\mathrm{id}_Q$ factors through the projective $Q\otimes Q^*$ and splits off $\one^+$. Forgetful preserves projectives: its right adjoint, equivariant coinduction, is exact, with underlying functor $M\mapsto M\oplus M^{**}$. Thus $\one$ is projective in $\C$, proving semisimplicity.
\end{proof}

\begin{theorem}\label{thm:radford}
Let $(W,\sigma)$ represent $\Ext^1_{\C}(\one,-)|_{\mathcal D}$, with universal class $\omega$ as in Proposition~\ref{prop:ext-object}. For $Z\in\mathcal D$ define $\vartheta_Z(\xi)=(j_Z^{-1})_*(\xi^{**})$, using $\one^{**}=\one$. If $\theta$ is the twist of $(W,\sigma)$, then
\[
 \vartheta_Z(g_*\omega)=(g\theta)_*\omega
 \quad\text{for every }g:W\to Z,\qquad \theta^2=\mathrm{id}_W.
\]
Each simple summand of $(W,\sigma)$ has twist $1$ or $-1$, and its underlying object contains no copy of $\one$. If $\theta=\mathrm{id}_W$ and one embedding of the field of spherical dimensions makes all simple dimensions of $\mathcal D$ positive, then $\C$ is semisimple.
\end{theorem}
\begin{proof}
Naturality of $\delta$ gives $\xi^{****}=(\delta_Z)_*\xi$, hence
\[
 \vartheta_Z^2(\xi)
 =\bigl(j_Z^{-1}(j_Z^{**})^{-1}\bigr)_*\xi^{****}
 =(\delta_Z^{-1})_*(\delta_Z)_*\xi=\xi.
\]
By Lemma~\ref{lem:extension-dual} and naturality of $j$,
$\vartheta_Z(g_*\omega)=(g\rho)_*\omega$ with $\rho=j_W^{-1}\nu_W$.
The inverse center braiding has twist $\theta^{-1}$, so $\nu_W=j_W\theta$ and $\rho=\theta$. The representing bijection of Proposition~\ref{prop:ext-object} gives $\theta^2=1$; its part (3) excludes the underlying unit. If $\theta=1$, every $\vartheta_Z$ is the identity, and Lemma~\ref{lem:sign-separation} applies.
\end{proof}

When the fusion category $\mathcal D$ is itself modular, $\mathcal Z(\mathcal D)\simeq\mathcal D\boxtimes\mathcal D^{\mathrm{rev}}$, and the simple objects of the center are the pairs $Y\boxtimes Y'$ of simples of $\mathcal D$, with underlying object $Y\otimes Y'$ and twist $\theta_Y\theta_{Y'}^{-1}$ \cite[Proposition~8.20.12]{EGNO}. The theorem then says that the extension object is a sum of pairs with $\theta_Y^2=\theta_{Y'}^2$ and $Y'\not\simeq Y^*$. If the squares of the twists of the simples of $\mathcal D$ are distinct, there are no such pairs.

\begin{corollary}[Odd-order center twists]\label{cor:odd-center-twists}
Let $\C$ be unimodular and finite with the Chevalley property, and let its semisimple part $\mathcal D$ have a spherical structure $j$ with $j^{**}j=\delta|_{\mathcal D}$. If every simple twist of $\mathcal Z(\mathcal D)$ has odd order and one embedding of the field of spherical dimensions makes all simple dimensions positive, then $\C$ is semisimple.
\end{corollary}
\begin{proof}
Theorem~\ref{thm:radford} gives $\theta_W^2=1$. Odd order forces $\theta_W=1$, and its dimension criterion applies.
\end{proof}

At odd Verlinde levels this criterion applies even when twists coincide. For $l\geq1$ let $\widetilde K_l$ be the fusion ring with basis $X_0=1,X_1,\ldots,X_l$ and
\[
 X_1X_i=X_{i-1}+X_i+X_{i+1}\quad(1\leq i<l),\qquad X_1X_l=X_{l-1}+X_l .
\]
This is the even part of the Verlinde ring of $\mathfrak{sl}_2$ at level $2l+1$, with $X_i$ the class of the simple of highest weight $2i$ \cite[Example~4.10.6]{EGNO}. The ring $\widetilde K_1$ is the Fibonacci ring, and $\widetilde K_2$ is the ring $(m,k,l,n)=(1,1,1,0)$ of Proposition~\ref{prop:small-tables}.

\begin{theorem}\label{thm:verlinde}
For every $l\geq 1$, every finite tensor category over $\CC$ with Grothendieck ring $\widetilde K_l$ in the simple basis is semisimple.
\end{theorem}

\begin{proof}
Let $\C$ be a categorification of $\widetilde K_l$. All its simples are self-dual. Put $Y=X_l$. Set $U_{2a}=X_a$ for $0\leq 2a\leq l+1$ and $U_{2a+1}=X_{l-a}$ for $0\leq 2a+1\leq l+1$. The table gives $[Y\otimes U_j]=[U_{j-1}]+[U_{j+1}]$ for $1\leq j\leq l$, where $U_0,\ldots,U_l$ are distinct and $U_{l+1}=U_l$. Starting with $Y\otimes U_0=U_1$, we show inductively that $Y\otimes U_j\simeq U_{j-1}\oplus U_{j+1}$ for $1\le j\le l$. The preceding decomposition gives $h(U_{j-1},Y\otimes U_j)=h(Y\otimes U_{j-1},U_j)=1$ and $h(Y\otimes U_j,U_{j-1})=1$ by Theorem~\ref{thm:frobenius}. So $U_{j-1}$ is both a subobject and a quotient of $Y\otimes U_j$, and has multiplicity one. The composite from this simple subobject to the isomorphic simple quotient is therefore nonzero, so $U_{j-1}$ splits off. Thus all the powers of $Y$ are semisimple, and every simple is a summand of such a power. Products of two simples are consequently summands of powers of $Y$, proving the Chevalley property. The category $\C$ is unimodular since $\one$ is its only invertible object.

By \cite[Remark~A.4(iii)]{EOsemi}, the fusion category $\mathcal D$ is the even part of a Temperley--Lieb realization at a root of unity $q$ for which $\zeta=q^2$ has order $N=2l+3$. For its spherical structure $j$ and a choice of modular braiding, the formulas of \cite[\S\S1--2]{EMfield} give
\[
  \dim_j(X_a)=[2a+1]_q=\sum_{h=-a}^a\zeta^h,\qquad
  \theta_{X_a}=\zeta^{a(a+1)}.
\]
 The ratio of $j^{**}j$ and $\delta|_{\mathcal D}$ is a monoidal automorphism of the identity. Its scalar on $X_1$ is $1$ since $X_1$ is a summand of $X_1^{\otimes 2}$, and every simple is a summand of a power of $X_1$. Hence $j^{**}j=\delta|_{\mathcal D}$.

Every simple twist in $\mathcal Z(\mathcal D)\simeq\mathcal D\boxtimes\mathcal D^{\rm rev}$ is a quotient of two $N$th roots of unity, hence has odd order. A single automorphism of $\mathbb Q(\zeta)$ sends $\zeta$ to $e^{2\pi i/N}$ and all $\dim_j(X_a)$ to
\[
 \frac{\sin((2a+1)\pi/N)}{\sin(\pi/N)}>0\qquad(0\leq a\leq l).
\]
Corollary~\ref{cor:odd-center-twists} applies.
\end{proof}

The even ring at level $1$ has rank one. In that case $\Ext^1_{\C}(\one,\one)=0$ \cite[Theorem~4.4.1]{EGNO}, so induction on length gives semisimplicity as well. At level $13$ the labels $0,5$ have equal twists, as do $3,6$. The identity on first extensions still separates the two signs in the sphericalization.

\begin{corollary}\label{cor:verlinde-nichols}
Let $\mathcal D$ be an even Temperley--Lieb fusion category with based ring $\widetilde K_l$, $l\geq1$. For every nonzero object $V$ in $\mathcal Z(\mathcal D)$ the Nichols algebra $\mathcal B(V)$ has infinitely many nonzero homogeneous components. The same holds for $V$ in $\mathcal D$ with its modular braiding.
\end{corollary}
\begin{proof}
Suppose that $B=\mathcal B(V)$ has finitely many nonzero degrees. Each degree is a quotient of the corresponding tensor power of $V$, so $B$ is a Hopf algebra object in $\mathcal Z(\mathcal D)$. Its antipode is invertible. On degree $n$ of the braided tensor algebra the antipode is $(-1)^n$ times the braid reversing the $n$ factors. This map descends to a surjection on the finite-length Nichols quotient in each degree and is therefore invertible there.

Take right $B$-modules internal to $\mathcal D$. With the half-braiding $Y\otimes B\to B\otimes Y$ of Proposition~\ref{prop:ext-object}, the coproduct gives the action on $M\otimes N$ by moving the first copy of $B$ past $N$ and applying the two actions. The central Hopf-monad construction makes this module category rigid \cite[\S5.3, Proposition~5.4 and Theorem~5.7]{BLV}, in the right-module convention. Kernels and cokernels are computed in $\mathcal D$. The free modules $S\otimes B$, for simple $S$ of $\mathcal D$, form a finite projective generating family by the free-forgetful adjunction and exactness of the forgetful functor. With its simple counit module $\one$ this is a finite tensor category.

The augmentation ideal $I=B_{>0}$ is nilpotent. For a simple right module $M$, the submodule $MI$ is zero, since $MI=M$ would imply $MI^n=M$ for every $n$. The simples are precisely the simples of $\mathcal D$ with counit action, and their products have counit action. Hence the based Grothendieck ring is $\widetilde K_l$. The regular module satisfies $BI=I\ne0$ because $B_1=V\ne0$. This contradicts Theorem~\ref{thm:verlinde}, since $I$ acts trivially on every sum of simples. Finally, the braided inclusion $\mathcal D\to\mathcal Z(\mathcal D)$ is exact and preserves the quantum symmetrizers defining the homogeneous Nichols quotients. This proves the assertion for $\mathcal D$.
\end{proof}

The proof constructs a finite realization from a hypothetical finite central Nichols algebra. In the other direction, for an ordinary finite-dimensional Hopf algebra $H$ with the Chevalley property, the radical-graded bosonization $R\#(H/\operatorname{rad} H)$ gives a finite Nichols quotient $\mathcal B(R_1)$ \cite{AS}. A construction of this kind from the universal extension object for a general finite tensor category with the Chevalley property remains open.

For the unitary $PSU(2)_k$ category at odd level, Li bounded the length of a generator of a finite Nichols algebra by one \cite[Theorem~6.4]{LiNichols}. Corollary~\ref{cor:verlinde-nichols} also rules out a single nonzero simple generator. At level $13$ it applies in particular to $X_5\boxtimes X_0$ and $X_3\boxtimes X_6$. Infinite-dimensionality leaves open the nilpotence of a specified homogeneous element. The following exact calculation gives a lower bound on the powers of a cup.

\begin{proposition}\label{prop:nichols-cup}
In the even part of $SU(2)_{13}$, let $Q\in\Hom(\one,\mathcal B^2(X_5))$ be the degree-two class of a nonzero invariant cup. Then $Q^4\ne0$.
\end{proposition}
\begin{proof}
Set $q=e^{\pi i/15}$ and $v=5$. In the left-associated fusion-path basis of $\Hom(\one,X_5^{\otimes 2j})$, the path $e=(5,0,\ldots,5,0)$ represents $j$ adjacent cups up to a nonzero scalar. The quotient onto the degree-$2j$ Nichols component has kernel equal to the kernel of the quantum symmetrizer $\mathfrak S_{2j}$. We compute one matrix entry of that symmetrizer.

In a block for recoupling $a,v,v$ to $b$, factor the usual orthogonal Racah matrix as $F_{xy}=\pm A_x T_{xy}B_y$, with $T$ the alternating quantum-factorial sum and
\[
 A_x^2=[2x+1]\Delta(a,v,x)^2\Delta(x,v,b)^2,\quad
 B_y^2=[2y+1]\Delta(v,v,y)^2\Delta(a,y,b)^2 .
\]
Here $[i]=(q^i-q^{-i})/(q-q^{-1})$ and $\Delta$ is the quantum triangle coefficient. The common sign cancels in the braid matrix. Rescale each path by the product of its vertex factors $\Delta$ and internal edge factors $\sqrt{[2x+1]}$. The braid entry with column $x$ and row $x'$ then becomes
$A_x^2\sum_y T_{x'y}B_y^2T_{xy}R_{vv}^y$, with
$R_{vv}^y=(-1)^{2v-y}q^{y(y+1)-2v(v+1)}$.
It is rational in $q$, and the diagonal of the symmetrizer is invariant under the rescaling. The denominators are products of $[1],\ldots,[14]$, all units at the specialization $q\mapsto3$ of order $30$ in $\mathbb F_{31}$. Multiplying the braid matrices by the recursion
$\mathfrak S_n=(\mathfrak S_{n-1}\otimes1)(1+T_{n-1}+T_{n-1}T_{n-2}+\cdots+T_{n-1}\cdots T_1)$
gives the following exact entries. The script \texttt{verify\_nichols.py} reproduces the calculation.
\[
\begin{array}{c|rrrr}
2j&2&4&6&8\\
\dim\Hom(\one,X_5^{\otimes 2j})&1&4&34&364\\
(\mathfrak S_{2j})_{ee}\pmod {31}&2&30&23&26
\end{array}
\]
The last entry is nonzero in $\mathbb F_{31}$, so its preimage in the cyclotomic field is nonzero. Thus the class of the four cups survives. The rational-gauge matrices for the other Galois realizations are obtained by applying the field automorphism entrywise. Vanishing of each minor is invariant under this operation, so all algebraic ranks, and the nonvanishing of the fourth power of this cup, are shared by those realizations.
\end{proof}

The numerical estimates of graded Frobenius--Perron dimensions for $X_5$ and $X_3\boxtimes X_6$ at level $13$, together with the rank thresholds, are given in the computational supplement (\texttt{nichols\_verlinde.py} and its output).

\section{Hankel determinants and tensor products}\label{sec:hankel}

Orthogonalizing powers via \eqref{eq:bridge} recovers simple summands or differences of Jordan blocks, bounds quantum-group orders and recognizes fusion rules.

\subsection{Hankel determinants and tensor-product recurrences}

For $c_0=1$ put $\Delta_0=1$ and
\[
 \Delta_k=\det(c_{i+j})_{0\leq i,j<k}\quad(k\geq1),\qquad
 \Delta_1=1,\qquad\Delta_2=c_2-c_1^2.
\]
The last quantity is the variance when a law exists. We use the standard Gram--Schmidt formulas \cite[Chapter~I]{Chi}.

\begin{theorem}\label{thm:jacobi}
Let $\mu$ be a probability measure on $\RR$ with all moments finite and $m$ support points, allowing $m=\infty$. Its unique monic orthogonal polynomials are $p_0=1$ and, for $1\le n<m$,
\[
 p_n(t)=\frac{1}{\Delta_n}\det\begin{pmatrix}c_0&c_1&\cdots&c_n\\ \vdots&\vdots&&\vdots\\ c_{n-1}&c_n&\cdots&c_{2n-1}\\ 1&t&\cdots&t^n\end{pmatrix}.
\]
Write $p_{-1}=0$ and $\lambda_0=0$. For $0\le n<m$ put $\beta_n=\int t p_n^2\,d\mu/\|p_n\|^2$ and, when $n\ge1$, $\lambda_n=\|p_n\|^2/\|p_{n-1}\|^2>0$. Then $tp_n=p_{n+1}+\beta_np_n+\lambda_np_{n-1}$ for $n+1<m$ and $\|p_n\|^2=\lambda_1\cdots\lambda_n=\Delta_{n+1}/\Delta_n$. At the last index of a finite law the recurrence holds in $L^2(\mu)$ with the next vector zero. Hence
\[
 \Delta_k=\prod_{i=0}^{k-1}\|p_i\|^2=\prod_{i=1}^{k-1}\lambda_i^{k-i}\qquad(k\leq m),
\]
and $\Delta_k=0$ for $k>m$.
\end{theorem}

A positive definite Gram matrix on $1,t,\ldots,t^n$ determines $p_0,\ldots,p_n$, their norms, $\lambda_1,\ldots,\lambda_n$ and $\beta_0,\ldots,\beta_{n-1}$; $\beta_n$ also requires $c_{2n+1}$. For self-dual $X$ with $x=[X]$, \eqref{eq:bridge} realizes the construction in the split ring with its Hom form, giving $h(p_n(x),p_n(x))=\Delta_{n+1}(X)/\Delta_n(X)$. Translation preserves determinants by unitriangular change of basis.

\begin{proposition}\label{prop:binet}
Assume global Cauchy--Schwarz and $X\simeq X^*$. Let $U$ have rows $u_{X^{\otimes i}}$, $0\le i<k$, whose $(S,r)$-coordinate counts summands in the $S$-block of length at least $r$. For semisimple $\C$ these are the simple multiplicity vectors. Then
\[
 \Delta_k(X)=\sum_T\det(U_T)^2,
\]
where $U_T$ ranges over the $k\times k$ column minors. In particular $\Delta_k(X)$ is a nonnegative integer.
\end{proposition}

\begin{proof}
By \eqref{eq:min} we have $h(A,B)=\langle u_A,u_B\rangle$, so $(c_{i+j})_{i,j<k}=UU^{\mathsf T}$, and the Cauchy--Binet formula gives the result. Only finitely many columns of $U$ are nonzero.
\end{proof}

For a Hamburger sequence, $\Delta_k=0$ forces $\Delta_{k+1}=0$. In contrast, the Taft module $M(2,2)$ gives $(\Delta_1,\ldots,\Delta_4)=(1,1,0,-1)$. The 6-dimensional $\Ga^2$-module in Section~\ref{sec:groups} gives $(1,3,-63)$, and Example~\ref{ex:affine} gives $(1,2,-17)$.

Set $S_{-1}=0$. Unit determinants force the following tensor product rules.

\begin{theorem}\label{thm:unit}
Let $X\simeq X^*$ and $K\ge1$. Assume $X^{\otimes j}$ is semisimple for
$0\le j<K$. The following conditions are equivalent.
\begin{enumerate}
\item $\Delta_1(X)=\cdots=\Delta_K(X)=1$.
\item There are pairwise nonisomorphic simples $S_0=\one,S_1,\ldots,S_{K-1}$
and integers $\beta_n\ge0$ such that
\[
 X\otimes S_n\simeq S_{n+1}\oplus S_n^{\oplus\beta_n}\oplus S_{n-1}
 \qquad(0\le n\le K-2).
\]
\end{enumerate}
Then $p_n(x)=[S_n]$ for $0\le n<K$ and $\lambda_n=1$ for $1\le n<K$.
\end{theorem}
\begin{proof}
The recurrences in (2) express each $[S_n]$ as a monic integral polynomial
of degree $n$ in $x$. These vectors are orthonormal, so the unitriangular
change of basis from powers gives (1).

Conversely, start with $S_0=\one$. If $[S_j]=p_j(x)$ have been constructed
through $j=n\le K-2$, consider $X\otimes S_n$, a summand of the semisimple
power $X^{\otimes(n+1)}$. Remove from its multiplicity vector the coordinates
of $S_0,\ldots,S_n$. The remainder is the monic Gram--Schmidt vector
$p_{n+1}(x)$, of squared norm $\Delta_{n+2}/\Delta_{n+1}=1$.
Its coordinates are nonnegative integers, so it consists of one new simple
$S_{n+1}$. For $j<n$ the removed multiplicity is
$h(S_j,X\otimes S_n)=h(X\otimes S_j,S_n)=\delta_{j,n-1}$ by reciprocity
and the earlier recurrences. The remaining diagonal multiplicity is
$\beta_n=h(S_n,X\otimes S_n)$. Every squared norm is one, giving the
stated $\lambda_n$.
\end{proof}

\begin{example}\label{ex:unit}
For the fundamental representations of $O_N^+$ ($N\ge2$) and $S_N^+$
($N\ge4$), all $\lambda_n=1$. The diagonal parameters are respectively
$(0,0,\ldots)$ and $(1,2,2,\ldots)$, giving the semicircle and the
rate-one free Poisson law. For the adjoint representation $V_2$ of
$SU(2)$ the diagonal is $(0,1,1,\ldots)$, giving the Riordan moments
$1,0,1,1,3,6,15,36,91,\ldots$. In the Fibonacci ring the path ends after
two vertices, while for the two-dimensional simple of $\Rep(S_3)$ it
ends after three, with the sign representation as its last vertex.
\end{example}

For $\Rep(\Ga)$ write $e_m=[J_m]-[J_{m-1}]$ as in Lemma~\ref{lem:layer-positive}.
The Jordan product rule of Theorem~\ref{thm:universal-jordan} gives an explicit
orthogonalization for $J_2$ and $J_3$.

\begin{theorem}\label{thm:hankel-ga}
In $\Rep(\Ga)$, the Hankel determinants and orthogonal polynomials of the blocks of sizes $2$ and $3$ are as follows.
\begin{enumerate}
\item For $X=J_2$ we have $\Delta_k=1$ for all $k$, with $p_n(x)=[J_{n+1}]-[J_n]=e_{n+1}$, and Jacobi parameters $\beta_0=1$, $\beta_n=0$ and $\lambda_n=1$ for $n\geq1$.
\item For $X=J_3$ we have $\Delta_k=2^{k-1}$ for $k\ge1$, with $p_0=1$ and $p_n(x)=[J_{2n+1}]-[J_{2n-1}]=e_{2n}+e_{2n+1}$ for $n\geq1$, and Jacobi parameters $\beta_n=1$, $\lambda_1=2$, $\lambda_n=1$ for $n\geq2$. In fact $\mu_{J_3}$ is the law of $1+2\cos\varphi$, with $\varphi$ uniform, and $c_n(J_3)$ are the central trinomial coefficients.
\end{enumerate}
\end{theorem}

\begin{proof}
For $J_2$, set $q_n=e_{n+1}$ for $n\ge0$. The Jordan product rule gives
$xq_0=q_1+q_0$ and $xq_n=q_{n+1}+q_{n-1}$ for $n\ge1$.
These orthonormal vectors are therefore the evaluations $p_n(x)$ of the
monic orthogonal polynomials.

For $J_3$, put $q_0=1$ and $q_n=e_{2n}+e_{2n+1}$ for $n\ge1$.
The same rule gives $xq_0=q_1+q_0$, $xq_1=q_2+q_1+2q_0$, and
$xq_n=q_{n+1}+q_n+q_{n-1}$ for $n\ge2$. The $q_n$ are orthogonal
with squared norms $1,2,2,\ldots$, giving the parameters and determinants
by Theorem~\ref{thm:jacobi}. The lifted character of $J_3$ is
$1+2\cos(2\theta)$. Under $\theta\mapsto 2\theta$ the tilted density
of Proposition~\ref{prop:additive-law} pushes forward to uniform measure,
as the cosine terms cancel on the two preimages. Its $n$th moment is the
constant coefficient of $(z^{-1}+1+z)^n$.
\end{proof}

\subsection{Hankel determinants in the layer basis}
Under global Cauchy--Schwarz the orthonormal basis
$e_{S,r}=[J_{S,r}]-[J_{S,r-1}]$ of Lemma~\ref{lem:layer-positive}
allows the same integral orthogonalization for arbitrary powers.

\begin{theorem}\label{thm:layer-path}
For a self-dual $X$ in a global Cauchy--Schwarz category and $K\geq1$,
the condition $\Delta_1=\cdots=\Delta_K=1$ is equivalent to the existence
of distinct basis vectors $v_0=e_{\one,1},v_1,\ldots,v_{K-1}$ and
integers $\beta_n\geq0$ such that, writing $v_{-1}=0$ and $x=[X]$,
\[
 xv_n=v_{n+1}+\beta_n v_n+v_{n-1}\qquad (0\leq n\leq K-2).
\]
In that case $v_n=p_n(x)$ for $0\leq n<K$.
\end{theorem}
\begin{proof}
If the determinants are $1$, the determinantal formula of
Theorem~\ref{thm:jacobi}, with $\Delta_0=1$, makes $p_n(x)$ an integral
vector of squared norm $1$ for $n<K$. It equals a signed basis vector,
and the coordinates used for different $n$ are distinct by orthogonality.
In the monic recurrence the coefficient of the next polynomial is $1$
and, for $1\leq n\leq K-2$, the coefficient of the previous one is the
ratio of 2 unit norms. Lemma~\ref{lem:layer-positive} forces consecutive
signs to agree and gives $\beta_n\geq0$ integral. Since $p_0(x)=e_{\one,1}$,
all signs are positive. Conversely the recurrence gives each $v_n$ as
a monic polynomial of degree $n$ in $x$. The unitriangular change from powers
to these orthonormal vectors gives all the stated determinants and identifies
the polynomials.
\end{proof}

Multiplication by the class of an object is essential in Lemma~\ref{lem:layer-positive}.
In $\Rep(\Ga)$, for example, the virtual class $e_2$ satisfies $e_2^2=e_3-e_2+e_1$.

\begin{corollary}\label{cor:variance-one}
For a self-dual object $X$ which is not semisimple in a global Cauchy--Schwarz
category, $\Delta_2(X)=1$ if and only if
$X\simeq J_{\one,2}\oplus\one^{\oplus a}$ for some $a\geq0$.
In this case $\Delta_k(X)=1$ for every $k$, and $\mu_X$ is the translate
by $a$ of the law of $J_2$ in $\Rep(\Ga)$.
\end{corollary}
\begin{proof}
$\Delta_2$ is the sum of squares of the coordinates of $[X]$ away from
$e_{\one,1}$. A nonsemisimple summand contributes a full initial segment of
at least 2 coordinates in its block. Exactly one unit square away from
$e_{\one,1}$ is thus possible precisely for the displayed objects.
Theorem~\ref{thm:universal-jordan} identifies their moment sequences, and
the measure of $J_2$ is compactly supported, hence determinate. Translation
preserves the Hankel determinants.
\end{proof}

\paragraph{An order bound from a finite string of moments.}
Let $K\subset S_N^+$ be finite, $N\ge4$, with fundamental representation $A$ and $D=\dim C(K)$. Since $A$ tensor-generates, its law has mass $1/D$ at $N$ by Proposition~\ref{thm:fp}. Agreement through degree $2m$ with the rate-one free Poisson law gives
\[
 D\ge\sum_{j=0}^m d_j^2,\qquad
 d_0=1,\quad d_1=N-1,\quad d_{j+1}=(N-2)d_j-d_{j-1}.
\]
Indeed the first $m+1$ free Poisson orthonormal polynomials remain orthonormal for $A$ and take values $d_j$ at $N$. Apply Bessel's inequality to $\sqrt D\,\mathbf1_{\{N\}}$ \cite[Theorem~9.2]{Simon}. If $S_N\subsetneq K$, then $c_4(A)=14$: it lies between $14$ and $15$, and equality with $15$ puts the flip in $\End(A^{\otimes2})$, forcing commutativity. Here $m=2$ gives $1+(N-1)^2+(N^2-3N+1)^2<2N!$, so the bound is illustrative: the proper Hopf inclusion $\CC[S_N]\subset C(K)^*$ already gives $D\ge2N!$ by Nichols--Zoeller freeness \cite{NZ}.

\subsection{Recognizing fusion rules from invariant dimensions}\label{sec:recognition}
For a marked basis element $x$ write $s_j=\tau(x^j)$, where $\tau$ takes the unit coefficient. These are the invariant dimensions of the corresponding simple in a fusion category. The distinction between $s_j$ and $c_j(X)$ in a general tensor category is the reason for the hypotheses in Theorem~\ref{thm:small-recognition} below.

\begin{proposition}[Recovery at known rank]\label{prop:small-tables}
At known rank 2 the based ring and its nonunit basis element $x$ are determined by $s_3$. At known rank 3, every nonpointed fusion ring has a self-dual basis element $x$ whose powers contain every basis element as a constituent and for which $s_3,s_4,s_5$ determine the marked based ring. Both cutoffs are sharp, even for rings admitting fusion categorifications.
\end{proposition}
\begin{proof}
In rank 2 $x^2=1+mx$, so $m=s_3$. The pointed and Fibonacci rings agree through degree 2. By \cite[Proposition~3.1]{O3}, the pointed rank-3 ring is $\ZZ[C_3]$. Every other rank-3 ring has self-dual basis $1,x,y$ and table
\[
 x^2=1+mx+ky,\qquad xy=kx+ly,\qquad y^2=1+lx+ny,
 \qquad k^2+l^2=kn+lm+1.
\]
At least one of $k,l$ is positive, so choose the marking to make $k>0$. The multiplicity vectors of $x^2$ and $x^3$ give
\[
 m=s_3,\qquad k=\sqrt{s_4-1-m^2},\qquad
 l=\frac{s_5-2m-m^3-2mk^2}{k^2},\qquad
 n=\frac{k^2+l^2-lm-1}{k}.
\]
Indeed $s_4=\langle x^2,x^2\rangle$ and $s_5=\langle x^2,x^3\rangle$. The 2-dimensional simple of $\Rep(S_3)$ and $V_2$ of the even level-5 $\mathfrak{sl}_2$ category give $1,0,1,1,3$ through degree 4 and fifth moments $5$ and $6$, respectively. This proves sharpness in rank three.
\end{proof}

Here generation means generation as a fusion subring, by taking constituents. For example, $(m,k,l,n)=(2,2,3,3)$ gives a fusion ring in which both $x$ and $y$ generate in this sense. Yet $\ZZ[x]=\ZZ1+\ZZ x+2\ZZ y$ and $\ZZ[y]=\ZZ1+3\ZZ x+\ZZ y$. The square relations give $2y\in\ZZ[x]$, $3x\in\ZZ[y]$, and the indicated lattices are stable under multiplication by $x$ and $y$, respectively.

For a self-dual simple $X$ in a finite realization with a fusion ring in the simple basis, the reconstruction formulas apply to invariant dimensions when the needed powers are semisimple. Indeed $c_{a+b}(X)=h(X^{\otimes a},X^{\otimes b})=s_{a+b}$ whenever both powers split. For $a+b\le5$ it suffices to know this for powers through degree 3. Theorem~\ref{thm:small-recognition} below uses 1 further moment to verify closure in rank 3.

The marked ring $(m,k,l,n)=(1,1,1,0)$ of the proof is the ring $\widetilde K_2$ of Section~\ref{sec:verlinde}. More generally, in the ring $\widetilde K_l$ multiplication by $X_1$ is given by the Jacobi matrix $J_l$ of size $l+1$ with diagonal $(0,1,\ldots,1)$ and adjacent diagonals equal to $1$, so that $\tau(X_1^n)=(J_l^n)_{00}$. The first sequences are
\[
\begin{array}{c|l}
 l & \tau(X_1^n),\ 0\leq n\leq 2l+2\\ \hline
 1 &1,0,1,1,2\\
 2 &1,0,1,1,3,6,14\\
 3 &1,0,1,1,3,6,15,36,90 .
\end{array}
\]
They agree through degree $2l+1$ with the Riordan moments of the adjoint representation of $SU(2)$ (Example~\ref{ex:unit}).

\begin{theorem}\label{thm:small-recognition}
Let $X\simeq X^*$ be an object of a complex tensor category and let $l\geq1$. Assume that $X^{\otimes j}$ is semisimple for $0\leq j\leq l+1$, and that $c_n(X)=(J_l^n)_{00}$ for $0\leq n\leq2l+2$. Then all the powers of $X$ are semisimple, and their summands and finite sums form a fusion subcategory $\mathcal D$ with Grothendieck ring $\widetilde K_l$, in which $[X]=X_1$. If the ambient category is finite and every simple occurs as a composition factor of a power of $X$, then $\C=\mathcal D$ for every $l\ge1$.
\end{theorem}

\begin{proof}
The spectral measure of $J_l$ at the first basis vector has $l+1$ atoms, and its Jacobi parameters are the entries of $J_l$. So the given moments have $\Delta_1=\cdots=\Delta_{l+1}=1$ and $\Delta_{l+2}=0$. Theorem~\ref{thm:unit} with $K=l+1$ gives nonisomorphic simples $S_0=\one,S_1,\ldots,S_l$ with $[S_n]=p_n(x)$ and $X\otimes S_n\simeq S_{n+1}\oplus S_n^{\oplus\beta_n}\oplus S_{n-1}$ for $n<l$, where $\beta_0=0$ and $\beta_n=1$ otherwise. At the terminal vertex, the monic polynomial $q=(t-1)p_l-p_{l-1}$ annihilates $J_l$, by its last row and cyclicity of the first basis vector. Moment agreement through degree $2l+2$ gives $h(q(x),q(x))=0$. All the terms of $q(x)=[X\otimes S_l]-[S_l]-[S_{l-1}]$ come from semisimple objects, since $X\otimes S_l$ is a summand of $X^{\otimes(l+1)}$. The squared norm is the sum of the squares of the simple multiplicity differences, so $X\otimes S_l\simeq S_l\oplus S_{l-1}$. Thus the sums of the $S_n$ are stable under tensoring with $X$, and by induction every power of $X$ is such a sum. The $S_n$ are self-dual, their classes being polynomials in $x$, and their products are summands of powers of $X$. Hence $\mathcal D$ is a fusion subcategory, and $[S_n]\mapsto X_n$ identifies its Grothendieck ring with $\widetilde K_l$. Under the last assumption the $S_n$ exhaust the ambient simples, so $\Gr(\C)=\widetilde K_l$ in the simple basis. Theorem~\ref{thm:verlinde} makes $\C$ semisimple, whence $\C=\mathcal D$.
\end{proof}

The cutoff $2l+2$ is sharp even among fusion categories. The weighted paths
for $J_l$ and $J_{l+1}$ have identical returns to vertex $0$ through
length $2l+1$. At length $2l+2$ the latter has exactly one additional
return, going directly to the new endpoint $l+1$ and back. These two
marked rings have different ranks and both admit fusion realizations.

The semisimplicity hypothesis matters, even for $l=1$. The self-dual Taft module $M(2,2)$ of Theorem~\ref{thm:taft-failure} has the Fibonacci moments $c_0,\ldots,c_4$. In degree 5 it gives $4$, and the Fibonacci simple gives $3$.

\section{Doubling and the Stieltjes problem}\label{sec:doubling}

For self-dual $X$, the invariants of $(X\otimes X)^{\otimes k}$ count
$\dim\End(X^{\otimes k})$. For $O_N^+$ ($N\ge2$) and $S_N^+$ ($N\ge4$),
the fattening bijection $NC_2(2k)\simeq NC(k)$ identifies the even moments
of the semicircle law with the moments of the free Poisson law of rate one
\cite{BCS}. For its planar algebra interpretation see \cite{Banica,Jones}.
For a general object the corresponding alternating product is $X\otimes X^*$.

\subsection{The Stieltjes criterion}
Let $X$ be an object with $X^{**}\simeq X$. This holds for every object of a pivotal or braided category, and for the simples of a category of fusion type. Then $Z=X\otimes X^*$ is self-dual, and we call
\[
        d_k(X)=c_k(X\otimes X^*)=\dim\Hom\bigl(\one,(X\otimes X^*)^{\otimes k}\bigr)
\]
the \emph{doubled sequence} of $X$. Theorem~\ref{thm:main} applies to $Z$. The point is that the shifted Hom dimensions also control the support of the measure.

\begin{theorem}\label{thm:stieltjes}
Let $X^{**}\simeq X$ and $Z=X\otimes X^*$. The following are equivalent.
\begin{enumerate}
\item The doubled sequence $(d_k(X))_{k\geq0}$ has a representing measure on $[0,\infty)$.
\item The inequality \eqref{eq:cs} holds for all $A,B$ which are finite direct sums of objects $Z^{\otimes i}$, and for all $A,B$ which are finite direct sums of objects $Z^{\otimes i}\otimes X$, $i\geq0$.
\end{enumerate}
\end{theorem}

\begin{proof}
Put $L(t^k)=d_k(X)$. As $Z$ is self-dual, $h(Z^{\otimes i},Z^{\otimes j})=d_{i+j}(X)$, and \eqref{eq:rigidity} gives
\[
 h(Z^{\otimes i}\otimes X,Z^{\otimes j}\otimes X)=h(Z^{\otimes i},Z^{\otimes j}\otimes X\otimes X^*)=d_{i+j+1}(X).
\]
So for $a,b\in\ZZ_{\geq0}[t]$ the sums $a(Z)$, $b(Z)$ have $h(a(Z),b(Z))=L(ab)$ as in \eqref{eq:bridge}, and the sums $a(Z)\otimes X$, $b(Z)\otimes X$ have Hom dimension $L(tab)$. By Stieltjes' theorem, recalled in Section~\ref{sec:moments}, (1) means that $L(p^2)\geq0$ and $L(tp^2)\geq0$ for every real polynomial $p$. The sign decomposition in the proof of Theorem~\ref{thm:main}, applied to each of the 2 families, shows that this is equivalent to (2).
\end{proof}

\begin{proposition}\label{prop:doubling}
Let $X\simeq X^*$, so that $d_k(X)=c_{2k}(X)=\dim\End(X^{\otimes k})$.
\begin{enumerate}
\item If $X$ has a law $\mu$, then $X\otimes X$ has the law $\mathrm{sq}_*\mu$ on $[0,\infty)$, where $\mathrm{sq}(t)=t^2$.
\item If $c_{2k+1}(X)=0$ for every $k$, then $X$ has a law if and only if $X\otimes X$ has a law on $[0,\infty)$.
\end{enumerate}
\end{proposition}

\begin{proof}
(1) is clear. For (2), order the monomials $1,t,\ldots,t^{2d+1}$ by parity. Since the odd moments vanish, the Hankel matrix of $X$ becomes block diagonal, with blocks $(c_{2i+2j})$ and $(c_{2i+2j+2})$. These are the Hankel matrix and the shifted Hankel matrix of the doubled sequence, and we conclude by the theorems of Hamburger and Stieltjes.
\end{proof}

The two diagonal blocks test sums of even powers and sums of odd powers, respectively. If all odd moments vanish, every Stieltjes law $\nu$ of $X\otimes X$ lifts to the symmetric law $\frac12(\sqrt{\,\cdot\,})_*\nu+\frac12(-\sqrt{\,\cdot\,})_*\nu$ of $X$.

When $X$ is not self-dual there are 2 doublings. The alternating one, considered here, is the one of subfactor theory. The other one, $\dim\End(X^{\otimes k})=\dim\Hom(\one,X^{\otimes k}\otimes X^{*\otimes k})$, is part of the mixed moments of Section~\ref{sec:complex}. For the fundamental representation of $U_N^+$, $N\geq2$, the doubled sequence consists of the Catalan numbers, as for $O_N^+$, while $\End(X^{\otimes k})=\CC$ for every $k$ \cite{BanU}.

\subsection{What doubling forgets}
We go through the examples of the previous sections. For the additive group, doubling removes the tilt of Proposition~\ref{prop:additive-law} when all block sizes have the same parity.

\begin{proposition}\label{prop:doubling-ga}
Let $X$ be a $\Ga$-module whose Jordan block sizes all have the same parity, with lifted character $\chi(\theta)$. Then the law of $X\otimes X$ is $(\chi^2)_*(d\theta/2\pi)$. Thus the doubled sequence of $X$ is that of the module of the torus $\Gm$ having the same weights.
\end{proposition}

\begin{proof}
For even block sizes all weights are odd and $\chi(\theta+\pi)=-\chi(\theta)$. For odd block sizes all weights are even and $\chi(\theta+\pi)=\chi(\theta)$. In both cases $\chi^2$ has period $\pi$, so the term $\cos\theta$ of the density in \eqref{eq:ga-law} cancels between $\theta$ and $\theta+\pi$.
\end{proof}

For example, $J_2$ and the $\Gm$-module of weights $\pm1$ both have doubled moments $\binom{2k}k$. Their odd moments are $\binom{2k+1}k$ and $0$, respectively. The full sequence for $J_3$ agrees with that of the torus module of weights $-2,0,2$, since its lifted character has period $\pi$. Within $\Rep(\Ga)$ the doubled sequence determines the Jordan type by Theorem~\ref{thm:reconstruction}.

For a self-dual unit cover the derivative term in \eqref{eq:cartan-hermite} vanishes on polynomials in $t^2$. This gives an exact test for the doubled law.

\begin{proposition}\label{prop:doubling-cartan}
Let $\C$ be a unimodular finite tensor category of fusion type. With $P,a,\nu_a$
as in Theorem~\ref{thm:cartan}, let $w_\lambda=\nu_a(\{\lambda\})$ and
$D=\{|\lambda|:\lambda\in\supp\nu_a,\ \lambda\ne0\}$.
Put $b_s=(w_s-w_{-s})/s$ for $s\in D$. The doubled sequence of $P$
has a Hamburger law precisely when $b_s\ge0$ for every $s\in D$ and
$\sum_{s\in D}b_s\le1$. That law is Stieltjes and equals
\[
 \sum_{s\in D}b_s\delta_{s^2}
 +\Bigl(1-\sum_{s\in D}b_s\Bigr)\delta_0 .
\]
\end{proposition}
\begin{proof}
Apply \eqref{eq:cartan-hermite} to $f(t)=q(t^2)$. The term $w_0f'(0)$ vanishes, and evaluation at the two eigenvalues $s,-s$ gives the single coefficient $(w_s-w_{-s})/s$ at $s^2$, with an absent eigenvalue assigned weight zero. The resulting signed measure has total mass one. The square of $t\prod_{s\in D}(t-s^2)$ has integral zero against it, so every positive representing measure is supported on the displayed nodes. Interpolation forces the stated masses.
\end{proof}

In Example~\ref{ex:exterior-witness} we have $\nu_a=(\delta_0+\delta_d)/2$, so $P\otimes P$ has the law $(1-\frac1{2d})\delta_0+\frac1{2d}\delta_{d^2}$, while $P$ has no law. Thus positivity of the doubled sequence alone can hold in a nonsemisimple category, even though its first moment $d_1(P)=c_2(P)$ still detects semisimplicity under unimodularity. For the quaternion example of Example~\ref{ex:quaternion-witness}, the same formula gives the law $\frac{19}{32}\delta_0+\frac38\delta_4+\frac1{32}\delta_{16}$ for $P\otimes P$, although $P$ has the negative Hankel determinant computed there.

Doubling does not always repair the failure of positivity. The Taft module $X=M(2,2)$ of Theorem~\ref{thm:taft-failure} has $(d_0,\ldots,d_4)=(1,1,2,7,28)$, and the Hankel determinant of size 3 of this sequence is $-1$. So $X\otimes X$ has no law either. The same happens for the module of the affine group in Example~\ref{ex:affine}, where the first negative Hankel determinant of the doubled sequence has size 6. The exterior square of the four-dimensional standard module has a much later obstruction.

\begin{theorem}\label{thm:exterior-doubled}
Let $X=\Lambda^2(\CC^4)$ restricted to $\CC^3\rtimes SL_3$, and let
$d_k=c_{2k}(X)$. A probability measure with exactly $42$ atoms in $(0,36)$ represents $d_0,\ldots,d_{83}$. Adjoining $d_{84}=c_{168}(X)$ precludes a positive measure on $\RR$. The leading Hankel determinants of $(d_k)$ and of $(d_{k+1})$
are positive through size $42$ and negative at size $43$.
\end{theorem}

\begin{proof}
By Lemma~\ref{lem:levi}, $c_n(X)$ counts the summands
$\mathrm{Sym}^k(\CC^4)$ of $X^{\otimes n}$ with multiplicity. Under
$SL_4=\mathrm{Spin}_6$, $X$ is the vector representation and $\CC^4$ a half-spin
representation. Its $k$th symmetric power has highest weight
$\tfrac k2(1,1,1)$ in a choice of coordinates. All weights of $X$
are integral, so only $k=2j$ occurs. The minuscule tensor product rule
therefore counts walks of length $n$ from $0$ to $(j,j,j)$ with steps
$\pm e_i$ in $x_1\ge x_2\ge|x_3|$.

Here is an integer formula for checking the needed counts.
Write $U_n(v)$ for the unrestricted walk count, $\rho=(2,1,0)$ and
$W=W(D_3)$ for the permutations with an even number of sign changes.
Weyl alternation \cite{GrabinerMagyar} gives
\[
 c_n(X)=\sum_{j=0}^{\lfloor n/3\rfloor}
        \sum_{w\in W}\det(w)\,U_n((j,j,j)+\rho-w\rho).
\]
For $B(k,s)=\binom{k}{(k+s)/2}$, interpreted as zero for inadmissible
lower indices, the unrestricted count is
\[
 U_n(x,y,z)=\sum_{k=0}^{n}\binom nk B(k,z) B(n-k,x+y) B(n-k,x-y).
\]
Indeed choose the $k$ steps in the third direction. In the first two
directions the coordinates $x+y$ and $x-y$ then take independent
unit steps. The program \texttt{verify\_exterior.py} uses these formulas
and integer determinants. With indices starting at zero, put
$D_m=(d_{i+j})_{i,j<m}$ and $D_m^+=(d_{i+j+1})_{i,j<m}$.
All leading minors of $D_{42},D_{42}^+,36D_{42}-D_{42}^+$ are positive.
The determinants of $D_{43}$ and $D_{43}^+$ are negative, with
$1603$ and $1641$ digits respectively. The independent chamber walk
program \texttt{exterior\_square\_walks.py} gives the same $173$ values
through degree $172$ by CRT. The a priori bound $c_n(X)\le6^n$ is
strictly below the product of its nine primes for $n\le172$. A third check, \texttt{verify\_exterior\_pieri.py}, iterates the vertical two-box Pieri rule on partitions of length at most four, summing the multiplicities of shapes $(a+j,a,a,a)$ at each degree. It reproduces all $173$ counts using nonnegative integers.

Give the polynomials of degree at most $41$ the inner product with
matrix $D_{42}$. The operator $T=D_{42}^{-1}D_{42}^+$ is self-adjoint
for this inner product, and $0<T<36I$. It sends $t^j$ to $t^{j+1}$ for
$j<41$, since $\langle Tt^j,t^i\rangle=d_{i+j+1}=\langle t^{j+1},t^i\rangle$ for each $0\le i\le41$. The spectral measure
of $T$ at $1$ has exactly $42$ atoms in $(0,36)$, since $1,T1,\ldots,T^{41}1$ span the space. It represents degrees through $82$ by their inner products, and degree $83$ by $\langle Tt^{41},t^{41}\rangle=d_{83}$. For $L(t^n)=d_n$, the Schur complement gives a compact negative square. Put $H=D_{42}$ and $v=(d_{42},\ldots,d_{83})^{\mathsf T}$. Then
\begin{equation}\label{eq:exterior-schur}
\begin{gathered}
 p(t)=t^{42}-(1,t,\ldots,t^{41})H^{-1}v,\\
 L(p^2)=d_{84}-v^{\mathsf T}H^{-1}v
       =\frac{\det D_{43}}{\det D_{42}}<0.
\end{gathered}
\end{equation}
Thus $d_{84}=c_{168}(X)$ is the first impossible degree for the doubled sequence; $\det D_{43}^+$ also uses $d_{85}$. Clearing denominators in $p$ and splitting its coefficients gives the two sums of powers of $X^{\otimes2}$ in Theorem~\ref{thm:main}.
\end{proof}

The values also satisfy the order-six polynomial recurrence in
\texttt{exterior\_square\_walks.py} at every tested starting index
$0\le n\le166$. Its extension to every $n$ is conjectural.

\begin{proposition}[Arbitrarily long positive truncations]\label{prop:long-truncations}
For every $k\ge1$ there are a finite complex tensor category $\C$ and a self-dual $Y\in\C$ whose doubled moments through degree $2k-1$ have a probability measure on $(0,\infty)$, although the full sequence has no Hamburger measure. The category and object may depend on $k$.
\end{proposition}
\begin{proof}
In $\Rep(T_3)$ take $Q=P_0\oplus P_2$ and $Y=\one\oplus Q$. Proposition~\ref{prop:taft-projectives} gives $c_1(Q)=1$ and $c_j(Q)=6^j/9$ for $j\ge2$. Expanding $(1+t)^{2n}$ gives
\[
 d_n=d_n(Y)=\frac{49^n+8+6n}{9}.
\]
For $m\ge1$, the external tensor power $Y_m=Y^{\boxtimes m}$ in $\Rep(T_3^{\otimes m})$ has doubled moments $d_n^m$, since the invariant spaces factor. Duality swaps $P_0$ and $P_2$, so every $Y_m$ is self-dual.

Each $d_n^m$ is a sum of polynomials in $n$ times $(49^a)^n$, $0\le a\le m$, with coefficient $(8+6n)^m/9^m$ at $a=0$. Its generating function is therefore rational with a pole of order $m+1$ at $1$. A polynomial $p$ annihilating this exponential-polynomial sequence gives $L(t^np)=0$ for every $n\ge0$, hence $L(p^2)=0$. A Hamburger measure would then have finite support in the real zero set of $p$, so its generating function would have only simple poles. Thus no doubled sequence of $Y_m$ has a Hamburger law.

On the other hand, putting $u_n=9d_n$, one has for $n\ge1$
\[
 u_{n-1}u_{n+1}-u_n^2=576(24n+7)49^{n-1}-36>0.
\]
Strict log-convexity gives $d_{i+j+s}^2<d_{2i+s}d_{2j+s}$ for $i\ne j$ and $s=0,1$. For fixed $k$, diagonal congruence normalizes $(d_{i+j+s}^m)_{i,j<k}$ to have diagonal $1$ and off-diagonal entries tending to $0$ as $m\to\infty$. Both matrices are positive definite for large $m$. As in Theorem~\ref{thm:exterior-doubled}, their positive multiplication operator supplies a probability measure on $(0,\infty)$ through degree $2k-1$.
\end{proof}

The construction varies $m$, hence the category and object, with the prescribed truncation. It excludes a universal finite test over these categories and is compatible with the rank-dependent projective bound of Theorem~\ref{thm:finite-test}.

\subsection{Jacobi parameters and fattening}
We now compare the orthogonal polynomials of the 2 laws, in the situation of Proposition~\ref{prop:doubling}~(2). We use the notation of Theorem~\ref{thm:jacobi}. The Hankel determinants of the shifted sequence $(d_{k+1})$ are denoted $\Delta_k^+$.

\begin{theorem}\label{thm:doubling-jacobi}
Let $X\simeq X^*$ have a law $\mu$ with infinite support and vanishing odd moments, with orthogonal polynomials $p_n$ and Jacobi parameters $\beta_n=0$ and $\lambda_n$. Then the law of $X\otimes X$ has orthogonal polynomials $q_n(t^2)=p_{2n}(t)$ and Jacobi parameters
\[
    \beta'_0=\lambda_1,\qquad\beta'_n=\lambda_{2n}+\lambda_{2n+1},\qquad\lambda'_n=\lambda_{2n-1}\lambda_{2n}\qquad(n\geq1),
\]
and we have $\Delta_{2k}(X)=\Delta_k(X\otimes X)\,\Delta_k^+(X\otimes X)$. Moreover, let $K\geq1$ be an integer and let $X\simeq X^*$ be any object with semisimple powers such that
\[
       X\otimes S_n\simeq S_{n+1}\oplus S_{n-1}\qquad(0\leq n<2K)
\]
for pairwise nonisomorphic simples $S_0=\one,S_1,\ldots,S_{2K}$, as in Theorem~\ref{thm:unit} with $\beta_n=0$, then $Z=X\otimes X$ satisfies
\[
   Z\otimes S_0\simeq S_2\oplus S_0,\qquad Z\otimes S_{2n}\simeq S_{2n+2}\oplus S_{2n}^{\oplus2}\oplus S_{2n-2}\qquad(1\leq n<K),
\]
and $\Delta_k(X\otimes X)=\Delta_k^+(X\otimes X)=1$ for $k\leq K$.
\end{theorem}

\begin{proof}
The polynomials $p_{2n}$ are even, so $p_{2n}(t)=q_n(t^2)$ with $q_n$ monic of degree $n$, and the $q_n$ are orthogonal for $\mathrm{sq}_*\mu$. Applying twice the recurrence $tp_n=p_{n+1}+\lambda_np_{n-1}$ gives
\[
 t^2p_{2n}=p_{2n+2}+(\lambda_{2n}+\lambda_{2n+1})p_{2n}+\lambda_{2n-1}\lambda_{2n}p_{2n-2},
\]
with $\lambda_0=0$, which is the recurrence of the $q_n$. The determinant $\Delta_{2k}(X)$ is the product of the squared norms of $p_0,\ldots,p_{2k-1}$. Those of even index are the norms of the $q_n$, with product $\Delta_k(X\otimes X)$, and those of odd index give the Gram determinant of $t,t^3,\ldots,t^{2k-1}$, which is $\Delta_k^+(X\otimes X)$. For the last assertion, tensoring twice with $X$ gives the decompositions. So the classes of the objects $Z^{\otimes i}$, $i<k$, are unitriangular combinations of $[S_0],[S_2],\ldots,[S_{2k-2}]$, and those of the objects $Z^{\otimes i}\otimes X$, $i<k$, of $[S_1],[S_3],\ldots,[S_{2k-1}]$. By the proof of Theorem~\ref{thm:stieltjes} the 2 determinants are the Gram determinants of these families for $h$, and they are equal to $1$.
\end{proof}

The last part is the fattening operation at the level of fusion rules. The parameters $\lambda_n=1$ are those of the semicircle law and of $O_N^+$, and $\beta'=(1,2,2,\ldots)$, $\lambda'_n=1$ are those of the free Poisson law and of $S_N^+$, as in Example~\ref{ex:unit}. For a finite path the same computation relates the Verlinde ring of $\mathfrak{sl}_2$ at level $2l+1$, generated by the class $x$ of the simple of highest weight $1$, to its even part: the Jacobi matrix of $x^2$ on the even vertices has diagonal $(1,2,\ldots,2)$, and $x^2-1=X_1$ has the Jacobi matrix $J_l$ of Section~\ref{sec:recognition}. So the sequences recognized by Theorem~\ref{thm:small-recognition} are shifted doubled sequences, $\tau(X_1^n)=\sum_k(-1)^{n-k}\binom nkd_k(x)$.

\section{Mixed moments and fusion rings}\label{sec:complex}
\subsection{A local characterization for normal objects}

For an object $X$ define its \emph{ordered mixed moments} by
\[
 c_{a,b}(X)=\dim\Hom\bigl(\one,X^{\otimes a}\otimes (X^*)^{\otimes b}\bigr).
\]
A \emph{complex spectral measure} is a positive measure on $\CC$ satisfying $\int z^a\bar z^b\,d\mu=c_{a,b}(X)$ for all $a,b\geq0$ \cite[Definition~1.1]{FFK}. For representations of compact groups this is the character law. In a general tensor category, other orders of the same letters can have different invariant dimensions. Put $x=[X]$ and $y=[X^*]$ in the split Grothendieck ring.

\begin{definition}\label{def:normal}
An object $X$ of a tensor category is \emph{normal} if $X^{**}\simeq X$ and $xy=yx$. It has \emph{moderate growth} if there are $C,R>0$ with $\dim\Hom(\one,W)\leq CR^n$ for every word $W$ of length $n$ in $X,X^*$.
\end{definition}

Self-dual objects are normal. In a pivotal category $X^{**}\simeq X$ is automatic, and in a braided category every object is normal, by the Drinfeld isomorphism $X\simeq X^{**}$ and the braiding $X\otimes X^*\simeq X^*\otimes X$. In a fusion category duality on simple classes is the based-ring involution,
so $X^{**}\simeq X$ and normality amounts to normality of $[X]$ in the fusion ring. If $X$ is normal, every word $W$ with $a$ letters $X$ and $b$ letters $X^*$ has class $x^ay^b$, so $\dim\Hom(\one,W)=c_{a,b}(X)$ depends only on $(a,b)$. The support estimates of Section~\ref{sec:criterion} apply to a word of length $n$ as well: its invariant dimension is at most $(\dim F(X))^n$ for a fiber functor $F$, or $\FPdim(X)^n$ in a finite tensor category, since $\FPdim(X^*)=\FPdim(X)$. Both bounds give moderate growth.

For $f=\sum_{a,b}\alpha_{ab}z^a\bar z^b\in\ZZ_{\geq0}[z,\bar z]$ let $f(X)=\bigoplus_{a,b}(X^{\otimes a}\otimes X^{*\otimes b})^{\oplus\alpha_{ab}}$, and let $\mathcal P_{X,X^*}$ be the class of these objects. For normal $X$ it contains, up to isomorphism, all finite sums of words in $X,X^*$. Let $L_X$ be the linear functional on $\CC[z,\bar z]$ with $L_X(z^a\bar z^b)=c_{a,b}(X)$, and write $\bar f(z)=\overline{f(z)}$, which conjugates the coefficients and exchanges $z,\bar z$.

\begin{lemma}\label{lem:complexbridge}
If $X$ is normal, then $\dim\Hom(f(X),g(X))=L_X(g\bar f)$ for all $f,g\in\ZZ_{\geq0}[z,\bar z]$.
\end{lemma}

\begin{proof}
By additivity, take $f=z^a\bar z^b$ and $g=z^c\bar z^d$, with $A,B$ the corresponding words. By \eqref{eq:rigidity}, $\Hom(A,B)\simeq\Hom(\one,B\otimes A^*)$. Since duality reverses tensor products and $X^{**}\simeq X$, the object $A^*$ is a word with $b$ letters $X$ and $a$ letters $X^*$, so $B\otimes A^*$ has class $x^{c+b}y^{d+a}$, and $\dim\Hom(A,B)=c_{c+b,d+a}=L_X(z^c\bar z^d\cdot\bar z^az^b)$.
\end{proof}

Forey--Fresán--Kowalski use an inequality on all objects of a symmetric tensor category to obtain mixed-moment positivity, then apply radial Carleman growth to obtain a measure \cite[Proposition~2.2]{FFK}. Under the same growth condition, the following theorem gives a necessary and sufficient inequality on sums of the tensor words of one normal object, with no braiding assumption.

\begin{theorem}\label{thm:complex}
Let $X$ be normal and suppose its radial moments satisfy
\[
   \sum_{n=1}^{\infty}c_{n,n}(X)^{-1/(2n)}=\infty,
\]
where a zero moment contributes $+\infty$. Then $X$ has a complex spectral measure if and only if
\[
 \bigl(\dim\Hom(A,B)\bigr)^2\leq\dim\End(A)\,\dim\End(B)\qquad(A,B\in\mathcal P_{X,X^*}).
\]
The measure is unique. If $c_{n,n}\le CR^{2n}$, it is supported on the disk $|z|\le R$.
\end{theorem}

\begin{proof}
If $\mu$ exists, Lemma~\ref{lem:complexbridge} gives $\dim\Hom(f(X),g(X))=\langle g,f\rangle_{L^2(\mu)}$, and the inequality is Cauchy--Schwarz in $L^2(\mu)$. Conversely, assume the inequality.

\emph{Hermitian symmetry and positivity.}
The polynomial objects are closed under direct sums, so the
argument of Theorem~\ref{thm:ext} applied to $A$ and $A^{\oplus n}\oplus B$
in both orders gives $h(A,B)=h(B,A)$ on $\mathcal P_{X,X^*}$. Taking
$A=\one$ and $B$ of class $x^ay^b$ gives $c_{a,b}=c_{b,a}$. Hence
$L_X(\bar f)=\overline{L_X(f)}$.
For an integral polynomial $u=a-b$ with $a,b$ polynomials with coefficients $\geq0$,
Lemma~\ref{lem:complexbridge} and the assumed inequality give
\[
  L_X(u\bar u)\ge
  \left(\sqrt{h(a(X),a(X))}-\sqrt{h(b(X),b(X))}\right)^2.
\]
The extension to real coefficients is by homogeneity and continuity in
each fixed degree, as in Theorem~\ref{thm:main}. For real-coefficient
$u,v$ the numbers $L_X(u\bar v)$ and $L_X(v\bar u)$ agree by symmetry,
so $L_X(|u+iv|^2)=L_X(|u|^2)+L_X(|v|^2)\ge0$.

\emph{Existence and uniqueness.} Put $s=(z+\bar z)/2$ and $t=(z-\bar z)/(2i)$. The positive functional just constructed restricts to a real functional on $\RR[s,t]$. Since $(s^2+t^2)^n-s^{2n}$ and $(s^2+t^2)^n-t^{2n}$ are sums of squares, both coordinate moments of order $2n$ are bounded by $L_X((s^2+t^2)^n)=c_{n,n}$. The radial hypothesis thus implies both coordinate Carleman conditions. Nussbaum's multidimensional Carleman theorem gives existence and uniqueness on $\RR^2$ \cite[Theorem~14.20]{Sch}. Equivalently, the complex radial theorem applies directly to $L_X$ \cite[Theorem~15.11]{Sch}. Finally $\int|z|^{2n}\,d\mu=c_{n,n}\le CR^{2n}$ forces $|z|\le R$ almost everywhere under the indicated extra hypothesis.
\end{proof}

If $X\simeq X^*$, then $c_{a,b}=c_{a+b}$, so $L_X(|z-\bar z|^2)=2c_2-c_2-c_2=0$, the measure lives on $\RR$, and we recover Theorem~\ref{thm:main} under the radial Carleman hypothesis. In the Deligne category for $GL_t$ at a transcendental complex parameter, the standard object has $c_{a,b}=0$ for $a\ne b$ and $c_{n,n}=n!$, giving the standard complex Gaussian \cite[\S3.3]{FFK}. The Carleman sum diverges, although no exponential bound holds. The global inequality yields the following consequence.

\begin{corollary}\label{cor:complexglobal}
In a tensor category satisfying global Cauchy--Schwarz, for instance a semisimple one, every normal object satisfying the radial Carleman condition has a unique complex spectral measure.
\end{corollary}

\begin{example}\label{ex:complex}
(1) In $\Vect_{\ZZ/n}$, the generator $\delta_g$ has $c_{a,b}=1$ if $a\equiv b\pmod n$ and $0$ otherwise, so $\mu$ is the uniform measure on the $n$th roots of unity.

(2) In $\Rep(\ZZ/n\times\Ga)$, which satisfies global Cauchy--Schwarz by \eqref{eq:min}, since its representations are sums of Jordan blocks on the reductive isotypic components, let $X=\delta\boxtimes J_2$, with $\delta$ a generating character of $\ZZ/n$. Then $c_{a,b}(X)=c_{a+b}(J_2)$ if $a\equiv b\pmod n$, and $0$ otherwise, so $\mu_X$ is the law of $T\omega$, with $T$ following the tilted arcsine law of $J_2$ and $\omega$ an independent uniform $n$th root of unity.

(3) For Sweedler's algebra, $X=P_0$ is normal, with $X^*=P_1$, since $P_i\otimes P_j\simeq R$ for all $i,j$. But $c_{1,0}=\dim\Hom(\one,P_0)=0$ while $c_{0,1}=\dim\Hom(\one,P_1)=1$, so Hermitian symmetry fails, and $P_0$ has no complex spectral measure. Explicitly, $A=\one$ and $B=\one\oplus P_1$ give $h(A,B)^2=4>3=h(A,A)h(B,B)$. By Proposition~\ref{prop:sweedler}, however, $P_0\oplus P_0^*=R$ has the real
law $\frac34\delta_0+\frac14\delta_4$. For comparison, in the symmetric
setting of \cite[Proposition~2.8(2)]{FFK} a complex law of $M$ induces
the real law $(2\,\mathrm{Re}\,z)_*\mu_M$ of $M\oplus D(M)$. The Sweedler
example rules out inferring existence for $X$ from existence for
$X\oplus X^*$ in general tensor categories.
\end{example}

\subsection{Ordered and full moments in a based ring}
Let $\tau$ be the unit coefficient and $N_x$ the regular matrix of left multiplication. In a finite-rank fusion ring the ordered mixed moments of $x$ are $\tau(x^a(x^*)^b)$. A probability law for these words already forces normality.

\begin{theorem}\label{thm:ordered-normal}
Let $R$ be a finite-rank fusion ring, $x\in R_{\CC}$, and $\tau$ the coefficient of $1$.
There is a positive measure on $\CC$ with mixed moments
$\tau(x^a x^{*b})$ for all $a,b\geq0$ if and only if $xx^*=x^*x$.
In that case the measure is unique. More precisely, if $q(x)=0$ with $\deg q=m\ge1$, a positive measure agreeing in all bidegrees $a+b\le2m$ already forces normality. The necessity, including this degree bound, holds for any algebraic element of a complex $*$-algebra with a faithful positive tracial state.
\end{theorem}
\begin{proof}
Let $q$ be an annihilator of degree $m\ge1$. A representing measure $\mu$ satisfies
$\int |q(z)|^2\,d\mu=\tau(q(x)q(x)^*)=0$, so its support lies in the
finite zero set of $q$. Interpolate $p(z)=\bar z$ there using a polynomial of degree at most $m-1$. Expanding both sides and using $\tau(x^*x)=\tau(xx^*)$ gives
\[
 \tau\big((x^*-p(x))(x^*-p(x))^*\big)
       =\int |\bar z-p(z)|^2\,d\mu(z)=0.
\]
The first square uses moments of total degree at most $2m$, and the second uses at most $2\max(1,m-1)\le2m$. Faithfulness gives $x^*=p(x)$, hence normality. For sufficiency,
apply the spectral theorem to the normal left multiplication matrix $N_x$ at $1$.
The masses on the zero set of $q$ are determined by polynomial interpolation.
\end{proof}

The algebraicity hypothesis is essential even among semisimple tensor categories.
For the fundamental representation $u$ of $U_N^+$, $N\geq2$, the tensor powers
$u^{\otimes a}$ are pairwise inequivalent simples, whereas
$u\otimes u^*\simeq\one\oplus r_{\alpha\beta}$ and
$u^*\otimes u\simeq\one\oplus r_{\beta\alpha}$ with
$r_{\alpha\beta}\not\simeq r_{\beta\alpha}$
by the free unitary fusion rule \cite[Theorem~1]{BanU}.
Thus $[u]$ is nonnormal, yet $c_{a,b}=h(u^{\otimes b},u^{\otimes a})=\delta_{ab}$.
Its ordered mixed moments are those of Haar probability on the unit circle.

\paragraph{Full noncommutative moments.}
For a unital based ring $R$ with $1$ in its distinguished basis, possibly of infinite rank, and $x\in R_{\CC}$, evaluation gives a functional
\[
 \varphi_x(p)=\tau\bigl(p(x,x^*)\bigr)
 \quad\text{on }\CC\langle z,z^*\rangle.
\]
It is unital, tracial and positive, since $\tau(aa^*)$ is the sum of squared absolute values of the basis coefficients of $a$. The same formula applies to the Grothendieck ring of a semisimple tensor category with any number of simples. Thus every object has a full noncommutative $*$-distribution recording the invariant dimensions of all words in $X,X^*$. In the $U_N^+$ example above,
\[
 \varphi_u(zz^*zz^*)=\dim\End(u\otimes u^*)=2,
 \qquad \int_{|z|=1}z\bar z z\bar z\,d\mu(z)=1.
\]
The decomposition $u\otimes u^*=\one\oplus r_{\alpha\beta}$ gives the first equality. Hence the ordered moments determine the circle law while the alternating word detects additional tensor-product information.

\begin{corollary}\label{cor:allnormal}
All elements of the positive cone $\ZZ_{\geq0}B$ of a finite-rank fusion ring
have a complex spectral measure if and only if $b+b^*$ is central for every $b\in B$.
\end{corollary}
\begin{proof}
The coefficients of $[x,x^*]=xx^*-x^*x$ are quadratic polynomials in the
real basis coordinates of $x$. If they vanish for all integral coordinates $\geq0$,
they vanish throughout $R_{\RR}$. For $s=s^*$ and $k=-k^*$ in $R_{\RR}$,
normality of $s+k$ now gives $[s,k]=0$. Given another self-adjoint $t$, take
$k=[s,t]$. Since $k$ commutes with $t$, traciality gives
$\tau(k^2)=\tau((st-ts)k)=0$. As $k^*=-k$, faithfulness implies $k=0$.
Every real self-adjoint element is therefore central. These elements are spanned by
the $b+b^*$. Conversely, centrality of the self-adjoint part of $x$
implies that it commutes with the skew-adjoint part, so $x$ is normal.
Apply Theorem~\ref{thm:ordered-normal}.
\end{proof}

For group rings the condition that every real element be normal occurs in the classification of Bovdi and Siciliano \cite{BS}. For example, in $\ZZ[Q_8]$ each $g+g^{-1}$ is central, so every object of
$\Vect_{Q_8}$ admits a complex measure. In $\ZZ[S_3]$ each basis element
is unitary, but $x=(12)+(123)$ has $xx^*=2\cdot1+2(13)$ and
$x^*x=2\cdot1+2(23)$. Thus the direct sum of 2 invertible simples in
$\Vect_{S_3}$ can already fail the mixed moment problem.

\paragraph{Formal codegrees.}\label{thm:plancherel}
For a normal $x$, the weights of the preceding law follow from the standard block formula $\tau|_\pi=f_\pi^{-1}\operatorname{Tr}_\pi$, where $f_\pi$ is the formal codegree of a simple block $\pi$ of $R_\CC$ \cite[Lemma~2.3 and Example~2.4]{Ost}. Applying it to the spectral projections gives
\[
 \nu_x=\sum_{\pi}\frac{1}{f_\pi}
   \sum_{\lambda\in\operatorname{Spec}(\pi(x))}
     \dim\ker(\pi(x)-\lambda I)\,\delta_\lambda ,
\]
with distinct eigenvalues in the inner sum. In the commutative case this is $\sum_\chi f_\chi^{-1}\delta_{\chi(x)}$, where $f_\chi=\sum_{b\in B}|\chi(b)|^2$. For a modular category with unitarily normalized $S$-matrix, $\chi_j(b)=S_{bj}/S_{0j}$ and $f_j^{-1}=|S_{0j}|^2$ \cite[Chapter~8]{EGNO}.

\subsection{Spectral measures and return counts in the census}\label{sec:census}
Write $B$ for the distinguished basis of a fusion ring $R$. We compare two invariants on the 353 multiplicity-free tables of ranks one through nine in \cite{Census,VSrings}. In a multiplicity-free table every $N_{ij}^k$ is zero or one. For categorifiability through rank six and the pivotal census through rank seven see \cite{LPR,VScategories}.

For the first invariant, collect the laws $\nu_b$ for normal $b\in B$, inserting a symbol $\varnothing$ for each nonnormal $b$. We call the resulting unordered multiset the \emph{basic-law multiset}. By Theorem~\ref{thm:ordered-normal} the symbol records the failure of the positive mixed-moment problem. For a normal $b$, Cramer's rule and the spectral theorem give
\[
 G_b(z)=\big\langle(zI-N_b)^{-1}1,1\big\rangle
 =\frac{\det(zI-N_b^{(0)})}{\det(zI-N_b)}
 =\sum_{\lambda\in\supp\nu_b}\frac{\nu_b(\{\lambda\})}{z-\lambda},
\]
where $N_b^{(0)}$ deletes the unit row and column. The poles of the reduced fraction give the atoms, with their residues as masses. Thus the multiset is computed by testing $bb^*=b^*b$ and sorting the reduced fractions and markers.

For the second invariant, retain also the unit coefficients of powers of nonnormal basis elements. For $d\geq2$, let
\[
   \mathcal F_d(R)=
      \bigl\{\!\bigl\{\,(\tau(b^2),\ldots,\tau(b^d)):b\in B\,\bigr\}\!\bigr\}
\]
be the unordered multiset, with multiplicities. Its entries count paths in the multiplication graph of $b$ beginning and ending at the unit. In a fusion categorification they are invariant dimensions of powers of the corresponding simple. The return counts are defined for every $b$, including those for which a positive complex law fails. For $b=b^*$, the transition probabilities $N_{bj}^k\FPdim(k)/(\FPdim(b)\FPdim(j))$, with $j,k\in B$, give return probability $\tau(b^n)/\FPdim(b)^n$ at time $n$.

\begin{proposition}[Recognition in the multiplicity-free census]\label{thm:census-return}
On the 353 tables through rank nine, $\mathcal F_5$ determines the based ring. Through rank eight, $\mathcal F_4$ suffices. At rank nine, $\mathcal F_4$ takes 137 distinct values on 142 rings, so degree five is sharp for this invariant.
\end{proposition}
\begin{proof}
For each $b$ start at $v=e_0$ and iterate $v\leftarrow N_bv$ five times. Record the unit coordinate at steps $2,3,4,5$ and sort the tuples over all $b$. The numbers of distinct multisets are listed in Table~\ref{tab:spectra-census}. This is exact integer arithmetic, reproduced by \texttt{verify\_census.py} on the bundled tables. As an explicit witness to sharpness, rank-nine source lines 30 and 114 give the same degree-four multiset. The tuple $(1,0,4,1)$ occurs once for line 30 and is absent for line 114. Even the auxiliary cubic $\tau((\sum_{b\in B}b)^3)$ equals 132 for both rings.
\end{proof}

\begin{table}[htbp]\centering\small
\caption{Basic-law multisets and return-count multisets in the multiplicity-free census.}\label{tab:spectra-census}
\begin{tabular}{@{}rrrrrrrrr@{}}
\toprule
& & & & \multicolumn{2}{c}{Nonunit basis elements} & \multicolumn{3}{c}{Distinct multisets}\\
\cmidrule(lr){5-6}\cmidrule(l){7-9}
Rank & Rings & Comm. & All & Self-dual & Nonnormal & Basic laws & $\mathcal F_4$ & $\mathcal F_5$\\
\midrule
1&1&1&1&0&0&1&1&1\\
2&2&2&2&2&0&2&2&2\\
3&4&4&4&6&0&4&4&4\\
4&10&10&10&22&0&10&10&10\\
5&16&16&16&50&0&16&16&16\\
6&39&37&37&141&0&39&39&39\\
7&43&40&40&190&0&43&43&43\\
8&96&77&80&480&10&96&96&96\\
9&142&116&119&802&40&139&137&142\\
\bottomrule
\end{tabular}
\end{table}

In the table, \emph{Comm.} counts commutative rings and \emph{All} counts rings for which every element of the positive cone has a complex law (Corollary~\ref{cor:allnormal}). The basis counts run over all the listed rings of that rank, omitting the unit and retaining different basis elements in the same automorphism orbit.

The basic-law multiset distinguishes the 211 rings through rank eight. At rank nine the sole coincidences are one triple and one pair, all separated by the cubic of $u=\sum_{b\in B} b$. Indeed $\tau(ub)=1$ for each $b$, so $\tau(u^3)=\sum_{i,j,k}N_{ij}^k$, with values
\[
\begin{array}{c|rrr|rr}
 \text{rank-nine source line}&1923&2051&2075&1936&2080\\ \hline
 \tau(u^3)&100&116&164&163&227 .
\end{array}
\]
In a finite realization of fusion type, Proposition~\ref{prop:chevalley-moments} identifies this cubic as the upper bound for $c_3(U)$, attained exactly when the simple products split. As an example of the information lost by inserting $\varnothing$, the four nonnormal basis elements of line 1923 all give $(\tau(b^3),\tau(b^4))=(0,1)$, whereas those of line 2051 give $(1,0)$. Nonnormal basis elements first occur at rank eight. On line 2887 there, $b_1^*=b_2$, $b_1b_2=1+b_6$ and $b_2b_1=1+b_5$.

The limitation to the given census is essential. In a group ring $\tau(g^n)=\mathbf1_{\operatorname{ord}(g)\mid n}$, and $g$ has the uniform law on the roots of unity of its order. Both $C_4\times C_4$ and $Q_8\times C_2$ have $1,3,12$ elements of orders $1,2,4$, respectively. Their rank-$16$ group rings have identical basic-law multisets and identical $\mathcal F_d$ for every $d$, although only the former is commutative. Their total cubics also coincide, with value $16^2$.

\paragraph{Computational data.}
The archive includes the 353 tables with their source-line identifiers, upstream revision and checksums. The program \texttt{census.py} computes exact resolvents for the basic laws and the subring and graph tests of Table~\ref{tab:braided-census}. The independent program \texttt{verify\_census.py} computes $\mathcal F_4,\mathcal F_5$ and checks the reported coincidences of laws using only unit coefficients. For a normal element at rank $r$ the law has at most $r$ atoms. Equality of holomorphic moments through degree $2r-1$ determines equality of two such measures by Vandermonde interpolation on the union of their supports.

{\footnotesize\begin{thebibliography}{99}
\bibitem{AS} N. Andruskiewitsch and H.-J. Schneider, Pointed Hopf algebras, in New Directions in Hopf Algebras, Math. Sci. Res. Inst. Publ., vol.~43, Cambridge Univ. Press, 2002, 1--68.
\bibitem{ARS} M. Auslander, I. Reiten and S. O. Smal\o,
Representation Theory of Artin Algebras,
Cambridge Studies in Advanced Mathematics, vol.~36, Cambridge University Press, 1995.
\bibitem{BanU} T. Banica, The free unitary compact quantum group,
\emph{Comm. Math. Phys.} \textbf{190} (1997), 143--172.
\bibitem{Banica}
T. Banica, Representations of compact quantum groups and subfactors,
\emph{J. Reine Angew. Math.} \textbf{509} (1999), 167--198.
\bibitem{BCS} T. Banica and B. Collins, Integration over quantum permutation groups, \emph{J. Funct. Anal.} \textbf{242} (2007), 641--657.
\bibitem{BV} H. Bercovici and D. Voiculescu, Regularity questions for free convolution, in Nonselfadjoint Operator Algebras, Operator Theory, and Related Topics, Operator Theory: Advances and Applications, vol.~104, Birkh\"auser, 1998, 37--47.
\bibitem{BW} H. Bercovici and J.-C. Wang, On freely indecomposable measures, \emph{Indiana Univ. Math. J.} \textbf{57} (2008), 2601--2610.
\bibitem{BS} V. Bovdi and S. Siciliano, Normality in group rings, \emph{St. Petersburg Math. J.} \textbf{19} (2008), 159--165.
\bibitem{Brauer} R. Brauer, On algebras which are connected with the semisimple continuous groups, \emph{Ann. of Math.} \textbf{38} (1937), 857--872.
\bibitem{BLV} A. Brugui\'eres, S. Lack and A. Virelizier, Hopf monads on monoidal categories, \emph{Adv. Math.} \textbf{227} (2011), 745--800. \href{https://doi.org/10.1016/j.aim.2011.02.008}{doi:10.1016/j.aim.2011.02.008}.
\bibitem{CVZ} H. Chen, F. Van Oystaeyen and Y. Zhang, The Green rings of Taft algebras,
\emph{Proc. Amer. Math. Soc.} \textbf{142} (2014), 765--775.
\bibitem{CK} X.-W. Chen and H. Krause, Introduction to coherent sheaves on weighted projective lines,
\href{https://arxiv.org/pdf/0911.4473}{\textcolor{black}{\mbox{arXiv:0911.4473}}} (2010).
\bibitem{Chi} T. S. Chihara, An Introduction to Orthogonal Polynomials, Gordon and Breach, 1978.
\bibitem{Conrad} B. Conrad, Reductive group schemes, in Autour des sch\'emas en groupes, Vol. I, Panoramas et Synth\`eses, vols.~42--43, Soci\'et\'e Math\'ematique de France, 2014, 93--444.
\bibitem{CF} R. E. Curto and L. A. Fialkow, Recursiveness, positivity, and truncated moment problems, \emph{Houston J. Math.} \textbf{17} (1991), 603--635.
\bibitem{CJNP} A. Czenky, D. Jaklitsch, D. Nikshych, J. Plavnik, D. Reutter, S. Sanford and H. Yadav, Pivotal Brauer--Picard groupoids and graded extensions, \href{https://arxiv.org/pdf/2511.04482}{\textcolor{black}{\mbox{arXiv:2511.04482}}} (2025).
\bibitem{DG} V. Drensky and C. K. Gupta, Constants of Weitzenb\"ock derivations and invariants of unipotent transformations acting on relatively free algebras,
\emph{J. Algebra} \textbf{292} (2005), 393--428.
\bibitem{EMfield} C. Edie-Michell and S. Morrison, A field guide to categories with $A_n$ fusion rules, \href{https://arxiv.org/pdf/1710.07362}{\textcolor{black}{\mbox{arXiv:1710.07362}}} (2017).
\bibitem{Eriksen}
E. Eriksen, Iterated extensions and uniserial length categories,
\emph{Algebr. Represent. Theory} \textbf{24} (2021), 273--286.
\bibitem{MOEtingof} P. Etingof, Answer to S. Palcoux, Non-semisimple categorification problem of fusion rings, \href{https://mathoverflow.net/a/409926}{MathOverflow 409661} (2021).
\bibitem{EGNO} P. Etingof, S. Gelaki, D. Nikshych and V. Ostrik,
Tensor Categories, Mathematical Surveys and Monographs, vol.~205, American Mathematical Society, 2015.
\bibitem{EOsemi} P. Etingof and V. Ostrik, On semisimplification of tensor categories, in Representation Theory and Algebraic Geometry, Trends in Mathematics, Birkh\"auser, 2022, 3--35.
\bibitem{FFK} A. Forey, J. Fres\'an and E. Kowalski,
Fixed-point statistics from spectral measures on tensor envelope categories,
\emph{Selecta Math. (N.S.)} \textbf{31} (2025), article~100.
\bibitem{GR} A. M. Gainutdinov and I. Runkel, Projective objects and the modified trace in factorisable finite tensor categories, \emph{Compos. Math.} \textbf{156} (2020), 770--821.
\bibitem{GW} R. Goodman and N. R. Wallach, Symmetry, Representations, and Invariants, Graduate Texts in Mathematics, vol.~255, Springer, 2009.
\bibitem{GrabinerMagyar} D. J. Grabiner and P. Magyar, Random walks in Weyl chambers and the decomposition of tensor powers, \emph{J. Algebraic Combin.} \textbf{2} (1993), 239--260.
\bibitem{GKS} R. K. Guy, C. Krattenthaler and B. E. Sagan, Lattice paths, reflections, \& dimension-changing bijections, \emph{Ars Combin.} \textbf{34} (1992), 3--15.
\bibitem{IW} E. In\"on\"u and E. P. Wigner, On the contraction of groups and their representations, \emph{Proc. Nat. Acad. Sci. U.S.A.} \textbf{39} (1953), 510--524.
\bibitem{Jones} V. F. R. Jones, Planar algebras, I, \emph{New Zealand J. Math.} \textbf{52} (2021), 1--107.
\bibitem{Koike} K. Koike, On the decomposition of tensor products of the representations of the classical groups: by means of the universal characters, \emph{Adv. Math.} \textbf{74} (1989), 57--86.
\bibitem{LP} M. Larsen and R. Pink, Determining representations from invariant dimensions,
\emph{Invent. Math.} \textbf{102} (1990), 377--398.
\bibitem{LiNichols} B. Li, Reflections of Nichols algebras in braided monoidal categories, \href{https://arxiv.org/abs/2611.08900}{arXiv:2611.08900} (2026).
\bibitem{LPR}
Z. Liu, S. Palcoux and Y. Ren, Classification of Grothendieck rings of complex fusion categories of multiplicity one up to rank 6, \emph{Lett. Math. Phys.} \textbf{112} (2022), article~54.
\bibitem{LPW} Z. Liu, S. Palcoux and J. Wu, Fusion bialgebras and Fourier analysis: analytic obstructions for unitary categorification, \emph{Adv. Math.} \textbf{390} (2021), article~107905.
\bibitem{Milne} J. S. Milne, Algebraic Groups, Cambridge Studies in Advanced Mathematics, vol.~170, Cambridge University Press, 2017.
\bibitem{NP} C. Negron and J. Plavnik, Cohomology of finite tensor categories: duality and Drinfeld centers, \emph{Trans. Amer. Math. Soc.} \textbf{375} (2022), 2069--2112.
\bibitem{NS} A. Nica and R. Speicher, Lectures on the Combinatorics of Free Probability, London Mathematical Society Lecture Note Series, vol.~335, Cambridge University Press, 2006.
\bibitem{NZ} W. D. Nichols and M. B. Zoeller, A Hopf algebra freeness theorem, \emph{Amer. J. Math.} \textbf{111} (1989), 381--385.
\bibitem{O2} V. Ostrik, Fusion categories of rank 2, \emph{Math. Res. Lett.} \textbf{10} (2003), 177--183.
\bibitem{Ost} V. Ostrik, On formal codegrees of fusion categories, \emph{Math. Res. Lett.} \textbf{16} (2009), 895--901.
\bibitem{O3} V. Ostrik, Pivotal fusion categories of rank 3, \emph{Mosc. Math. J.} \textbf{15} (2015), 373--396.
\bibitem{Census} S. Palcoux, Fusion Ring Census,
\href{https://github.com/sebastienpalcoux/fusion-ring-census}{\textcolor{black}{\texttt{github.com/sebastienpalcoux/fusion-ring-census}}} (2026).
\bibitem{Sch} K. Schm\"udgen, The Moment Problem, Graduate Texts in Mathematics, vol.~277, Springer, 2017.
\bibitem{Shimizu}
K. Shimizu, Non-degeneracy conditions for braided finite tensor categories,
\emph{Adv. Math.} \textbf{355} (2019), \mbox{article}~106778.
\bibitem{Simon} B. Simon, The Christoffel--Darboux kernel, in Perspectives in Partial Differential Equations, Harmonic Analysis and Applications, Proc. Sympos. Pure Math., vol.~79, American Mathematical Society, 2008, 295--335.
\bibitem{Timashev} D. A. Timashev, Equivariant embeddings of homogeneous spaces, \href{https://www-fourier.ujf-grenoble.fr/sites/default/files/REF_611.pdf}{Pr\'epublication de l'Institut Fourier no.~611} (2003).
\bibitem{VSrings} G. Vercleyen and J. K. Slingerland, On low rank fusion rings,
\emph{J. Math. Phys.} \textbf{64} (2023), article~091703.
\bibitem{VScategories} G. Vercleyen and J. Slingerland, On low-rank multiplicity-free complex fusion categories, \emph{Comm. Math. Phys.} \textbf{407} (2026), article~231. \href{https://doi.org/10.1007/s00220-026-05705-2}{doi:10.1007/s00220-026-05705-2}.
\bibitem{WHLZ} Z. Wang, L. Li and Y. Zhang, Green rings of pointed rank one Hopf algebras of nilpotent type,
\emph{Algebr. Represent. Theory} \textbf{17} (2014), 1901--1924.
\bibitem{WLZ}
Z. Wang, L. Li and Y. Zhang, A criterion for the Jacobson semisimplicity of the Green ring of a finite tensor category,
\emph{Glasg. Math. J.} \textbf{60} (2018), 253--272.
\bibitem{Zhou} M. Zhou, Return probabilities of random walks,
\href{https://arxiv.org/pdf/1512.04142}{\textcolor{black}{\mbox{arXiv:1512.04142}}} (2015).
\end{thebibliography}}
\end{document}
